\documentclass{mathnote_cmyk}
\usepackage{mathnote_cmyk}

\begin{document}

% タイトルページの呼び出し
\tituloum
  {Linear Algebra}
  {Namachan}
  {Mathematics Notes}
  {A concise introduction to matrices, linear maps, determinants, and vector spaces}
  {\today}

%\include{Preambulo/Capa}

\cleardoublepage
\pagenumbering{arabic}
\pagecolor{white}

% 目次の表示
\localtableofcontents

\phantomsection
\part{行列と連立一次方程式}

\begin{abstract}
  この章では，線型代数の出発点となる\textbf{行列}を定義し，その基本的な扱いを整理する．
  まず行列とその相等・転置などの基本的な用語を確認し，和・スカラー倍・積という演算を定義して，
  その計算規則を調べる．とくに，行列の積が交換律を満たさないことに注意する．
  次に，正方行列に対して逆行列を定義し，その一意性と$2$次正方行列における具体形を確かめる．
  最後に，連立一次方程式を係数行列を用いて表し，加減法が行基本変形として言い換えられること，
  および逆行列を用いた解法を見る．
\end{abstract}

\section{行列}

\phantomsection
\subsection{行列の定義}

以下，本書では$K$は$\mathbb{R}$または$\mathbb{C}$を表すと約束する．
また，本書では自然数に$0$を含めない．すなわち，自然数とは正の整数$1, 2, 3, \dots$のことであり，自然数全体の集合を$\mathbb{N} = \{1, 2, 3, \dots\}$と表す．

まず，本書を通じて主役となる\textbf{行列}\index{ぎょうれつ@行列}を定義しよう．

\begin{definition}{行列}{行列}
  $mn$個の数$a_{ij} \in K$~$(1 \le i \le m,~1 \le j \le n)$を，縦に$m$個，横に$n$個並べた表
  \[
    A =
    \begin{pmatrix}
      a_{11} & a_{12} & \cdots & a_{1n} \\
      a_{21} & a_{22} & \cdots & a_{2n} \\
      \vdots & \vdots & \ddots & \vdots \\
      a_{m1} & a_{m2} & \cdots & a_{mn}
    \end{pmatrix}
  \]
  を\textbf{$m \times n$行列}（$m \times n$型の行列）という．
  横の並びを\textbf{行}\index{ぎょう@行}，縦の並びを\textbf{列}\index{れつ@列}とよび，第$i$行と第$j$列の交わる位置にある数$a_{ij}$を$A$の\textbf{$(i, j)$成分}\index{せいぶん@成分}という．
  また，この$A$を$A = (a_{ij})$と略記することがある．
\end{definition}

高校までの学習で，列ベクトル
\[
  \begin{pmatrix} x_1 \\ x_2 \\ \vdots \\ x_n \end{pmatrix}
\]
を学ぶが，これは$n \times 1$行列である．同様に行ベクトル
\[
  \begin{pmatrix} x_1 & x_2 & \cdots & x_n \end{pmatrix}
\]
は$1 \times n$行列である．なお，以下では行ベクトルを
\[
  (x_1, x_2, \dots, x_n)
\]
のようにカンマ区切りで表す表記も用いることにする．

\begin{example}{行列}{行列}
  \[
    A =
    \begin{pmatrix}
      -1 & 5 & 3 & 4 \\
      9 & -2 & 4 & 1 \\
      6 & -10 & -3 & 5
    \end{pmatrix}
  \]
  としたとき，$A$は$3 \times 4$行列である．
  \[
    \begin{pmatrix} -1 & 5 & 3 & 4 \end{pmatrix}
  \]
  を第$1$行といい，同様に第$2$行，第$3$行が定まる．
  また，
  \[
    \begin{pmatrix} -1 \\ 9 \\ 6 \end{pmatrix}
  \]
  を第$1$列といい，同様に第$2$列から第$4$列までが定まる．
  たとえば，$A$の$(2, 3)$成分は$4$である．
\end{example}

\begin{mycolumn}
  行列を単なる「数の表」として定義したことに，戸惑うかもしれない．
  また，このあと「2つの行列が等しい」と定めたり，行列どうしの和や積を定義したりすることにも，
  はじめは人為的な印象を受けるだろう．
  しかし，これらの定義の意味は，のちに\textbf{線型写像}を学ぶことで明らかになる．
  行列とは線型写像を数の表として書き下したものであり，たとえば行列の積は写像の合成に対応するように定められている．
  当面は天下り的に見える定義も，この見通しを頭の片隅に置いて読み進めてほしい．
\end{mycolumn}

新しい対象を定義したら，次に「$2$つの対象がいつ等しいとみなされるのか」を約束しておく必要がある．
行列については，次のように定める．
 
\begin{definition}{行列の相等}{行列の相等}
  $2$つの行列$A = (a_{ij})$，$B = (b_{ij})$が\textbf{等しい}\index{ぎょうれつのそうとう@行列の相等}とは，$A$と$B$の型が等しく，かつ対応するすべての成分が等しいこと，すなわち$A$，$B$がともに$m \times n$行列であって
  \[
    a_{ij} = b_{ij} \quad (1 \le i \le m,~1 \le j \le n)
  \]
  が成り立つことをいう．このとき$A = B$とかく．
\end{definition}
 
型が等しいことも条件に含まれている点に注意してほしい．
たとえば，$2 \times 2$の零行列と$2 \times 3$の零行列は，どちらも$O$という同じ記号で表されるが，等しい行列ではない．
 
また，この定義によって，$1$つの行列の等式$A = B$は$mn$個の数の等式$a_{ij} = b_{ij}$を束ねたものとなる．
のちに逆行列を求める場面などで，行列の等式から成分を比較して連立方程式を取り出す議論が現れるが，それを支えているのがこの定義である．
 
\begin{remark}[行列の相等を考える意義]
  「等しい」という関係をわざわざ定義することを，不思議に思う方もいるかもしれない．
  しかし，これは数学の議論を始めるうえで欠かせない約束である．
 
  実数の世界を思い出してみよう．円周率$\pi$は，小数表示$3.1415\ldots$のほかに，無限級数
  \[
    4 \left( 1 - \frac{1}{3} + \frac{1}{5} - \frac{1}{7} + \cdots \right)
  \]
  としても，積分
  \[
    \int_0^1 \frac{4}{1 + x^2} \, dx
  \]
  としてもかける．見た目のまったく異なるこれらの式を「$=$」で結べるのは，実数における「$=$」が「表記のかたちが同じ」ではなく「同じ数を表している」という意味に定められているからである．
 
  行列でも事情は同じで，$1$つの行列がさまざまな計算の結果として現れうる．
  そこで，どのような経緯で得られた行列であっても，「型が等しく，すべての成分が等しければ同じ行列とみなす」と約束しておくのである．
  逆にいえば，この約束があるからこそ，これから定義する和や積の計算結果を$1$つの行列として自信をもって「$=$」で結ぶことができる．
\end{remark}

\begin{definition}{正方行列}{正方行列}
  行の数と列の数が等しい行列を\textbf{正方行列}\index{せいほうぎょうれつ@正方行列}という．
  すなわち，$n$を自然数として，$n \times n$行列を\textbf{$n$次正方行列}という．
\end{definition}

\begin{example}{正方行列}{正方行列}
  \[
    \begin{pmatrix} 2 & 5 \\ -8 & 7 \end{pmatrix}, \quad
    \begin{pmatrix} -1 & 4 & 0 \\ 2 & 1 & 9 \\ -8 & 5 & 1 \end{pmatrix}
  \]
  はそれぞれ$2$次，$3$次の正方行列である．
\end{example}

行列全体の集合にも記号を用意しておこう．

\begin{definition}{行列全体の集合}{行列全体の集合}
  $K$の元\footnote{集合に属する個々の対象を，その集合の\textbf{元}（\textbf{要素}）とよぶ．本書では「元」の語に統一する．}を成分とする$m \times n$行列全体の集合を
  \[
    M_{m,n}(K)
  \]
  と表す．とくに，$n$次正方行列全体の集合を
  \[
    M_n(K) = M_{n,n}(K)
  \]
  と表す．
\end{definition}

この記法は，のちに「行列を変数とする写像」を考える場面（たとえば行列式）や，「$m \times n$行列全体が一つの線型空間をなす」ことを述べる場面（\partref{線型空間}）で用いられる．

のちの議論で頻繁に現れる「特別な行列」にも名前をつけておこう．

\begin{definition}{零行列と単位行列}{零行列と単位行列}
  すべての成分が$0$である行列を\textbf{零行列}\index{れいぎょうれつ@零行列}とよび，$O$と表す．

  また，$(1,1)$成分から$(n,n)$成分までの対角線上の成分がすべて$1$で，他の成分がすべて$0$である$n$次正方行列
  \[
    E_n =
    \begin{pmatrix}
      1 & 0 & \cdots & 0 \\
      0 & 1 & \cdots & 0 \\
      \vdots & \vdots & \ddots & \vdots \\
      0 & 0 & \cdots & 1
    \end{pmatrix}
  \]
  を\textbf{$n$次単位行列}\index{たんいぎょうれつ@単位行列}とよぶ\footnote{単位行列を表す文字は$E$のほかに$I$も使われる．英語表記の``Identity Matrix''の頭文字の$I$をとるか，もしくはドイツ語表記の``Einheitsmatrix''の頭文字をとるかの選択である．}．
  サイズが文脈から明らかなときは，単に$E$とかく．
\end{definition}

\begin{definition}{転置行列}{転置行列}
  $m \times n$行列$A = (a_{ij})$に対して，$A$の行と列を入れ替えて得られる$n \times m$行列，すなわち$(i, j)$成分が$a_{ji}$である行列を$A$の\textbf{転置行列}\index{てんちぎょうれつ@転置行列}とよび，$A^\mathsf{T}$と表す\footnote{行列$A$の転置行列の表記として，$A^\mathsf{T}$のほかに${}^t A$もよく使われる．}．
\end{definition}

\begin{example}{転置行列}{転置行列}
  \[
    A = \begin{pmatrix} 1 & 2 & 3 \\ 4 & 5 & 6 \end{pmatrix}
    \quad \text{のとき，} \quad
    A^\mathsf{T} = \begin{pmatrix} 1 & 4 \\ 2 & 5 \\ 3 & 6 \end{pmatrix}
  \]
  である．$A$の第$1$行が$A^\mathsf{T}$の第$1$列になっていることに注意してほしい．
\end{example}

いきなり行列という概念が出てきて，戸惑っている方も多いと思う．
これから行列の計算規則や意味などを確認していこう．

\phantomsection
\subsection{行列の和とスカラー倍}

\begin{definition}{行列の和}{行列の和}
  $m \times n$行列$A = (a_{ij})$，$B = (b_{ij})$に対して，その\textbf{和}\index{ぎょうれつのわ@行列の和}$A + B$を
  \[
    A + B =
    \begin{pmatrix}
      a_{11} + b_{11} & a_{12} + b_{12} & \cdots & a_{1n} + b_{1n} \\
      a_{21} + b_{21} & a_{22} + b_{22} & \cdots & a_{2n} + b_{2n} \\
      \vdots & \vdots & \ddots & \vdots \\
      a_{m1} + b_{m1} & a_{m2} + b_{m2} & \cdots & a_{mn} + b_{mn}
    \end{pmatrix}
  \]
  と定義する．すなわち，成分ごとに和をとる．
\end{definition}

型の異なる行列に対しては，和は定義されない．

\begin{definition}{行列のスカラー倍}{行列のスカラー倍}
  $\alpha \in K$と$m \times n$行列$A = (a_{ij})$に対して，$A$の\textbf{スカラー倍}\index{ぎょうれつのすからーばい@行列のスカラー倍}$\alpha A$を
  \[
    \alpha A =
    \begin{pmatrix}
      \alpha a_{11} & \alpha a_{12} & \cdots & \alpha a_{1n} \\
      \alpha a_{21} & \alpha a_{22} & \cdots & \alpha a_{2n} \\
      \vdots & \vdots & \ddots & \vdots \\
      \alpha a_{m1} & \alpha a_{m2} & \cdots & \alpha a_{mn}
    \end{pmatrix}
  \]
  と定義する．すなわち，すべての成分を$\alpha$倍する．
\end{definition}

\begin{example}{行列のスカラー倍}{行列のスカラー倍}
  $A = \begin{pmatrix} 2 & 3 \\ 9 & 1 \\ -8 & 3 \end{pmatrix}$に$5$をかけると，
  \[
    5A =
    \begin{pmatrix}
      5 \cdot 2 & 5 \cdot 3 \\
      5 \cdot 9 & 5 \cdot 1 \\
      5 \cdot (-8) & 5 \cdot 3
    \end{pmatrix}
    =
    \begin{pmatrix}
      10 & 15 \\
      45 & 5 \\
      -40 & 15
    \end{pmatrix}
  \]
  となる．
\end{example}

\begin{remark}[スカラー倍の定義の設計]
  行列のスカラー倍を，たとえば「第$1$行だけを$\alpha$倍する」操作として定めることも，定義としては可能である．
  この操作を仮に$\alpha \ast A$とかくと，$1 \ast A = A$や$\alpha \ast (A + B) = \alpha \ast A + \alpha \ast B$は成り立つが，
  \[
    (\alpha + \beta) \ast A = \alpha \ast A + \beta \ast A
  \]
  は一般に成り立たない．実際，左辺の第$2$行以降は$A$のままであるのに対し，右辺の第$2$行以降は$A$の対応する行の$2$倍になる．
  \deref{行列のスカラー倍}が「すべての成分を$\alpha$倍する」という形をしているのは，こうした計算規則が成り立つように選ばれた結果である．
  定義は天下りに与えられるものではなく，望ましい性質から逆算して設計されるものである．
\end{remark}

\begin{definition}{行列の減法}{行列の減法}
  $A$，$B$を型の等しい行列とする．
  \deref{行列のスカラー倍}に従って行列$(-1)B$を考え，
  \[
    A + (-1)B
  \]
  を$A - B$とかく．
  以下，$(-1)B$を$-B$と略記する．
\end{definition}

\deref{行列の減法}は，減法を和とスカラー倍から組み立てる定義である．
実際に成分を計算すれば，$A - B$の$(i, j)$成分は$a_{ij} - b_{ij}$であり，成分ごとに差をとることと一致する．

\begin{example}{行列の和と差の計算}{行列の和と差の計算}
  \[
    A = \begin{pmatrix} 2 & 0 & -1 \\ 1 & 1 & 4 \end{pmatrix}, \quad
    B = \begin{pmatrix} 1 & 2 & 3 \\ 0 & 1 & -2 \end{pmatrix}, \quad
    C = \begin{pmatrix} 0 & 1 \\ 3 & 5 \\ 2 & 0 \end{pmatrix}
  \]
  としたとき，
  \[
    A + B =
    \begin{pmatrix}
      2 + 1 & 0 + 2 & -1 + 3 \\
      1 + 0 & 1 + 1 & 4 + (-2)
    \end{pmatrix}
    =
    \begin{pmatrix}
      3 & 2 & 2 \\
      1 & 2 & 2
    \end{pmatrix}, \quad
    A - B =
    \begin{pmatrix}
      1 & -2 & -4 \\
      1 & 0 & 6
    \end{pmatrix}
  \]
  である．一方，$A$と$C$は型が異なるので，$A + C$や$A - C$は定義されない．
\end{example}

行列の和について成り立つ計算規則をまとめておこう．

\begin{prop}{行列の和の計算規則}{行列の和の計算規則}
  $A$，$B$，$C$を型の等しい行列とし，$O$を同じ型の零行列とする．このとき，次が成り立つ：
  \begin{align*}
    (A + B) + C &= A + (B + C) && \text{（和の結合律）} \\
    A + B &= B + A && \text{（和の交換律）} \\
    A + O &= A && \text{（加法単位元の存在）} \\
    A + (-A) &= O && \text{（加法逆元の存在）}
  \end{align*}
\end{prop}

\begin{leftbar}
\begin{proof}
  行列の和とスカラー倍は成分ごとの計算として定義されているから，いずれの等式も，両辺の$(i, j)$成分を比較すれば，数$a_{ij}, b_{ij}, c_{ij} \in K$についての対応する等式に帰着される．
  たとえば第$1$式については，両辺の$(i, j)$成分はともに$a_{ij} + b_{ij} + c_{ij}$であり，これは数の和の結合律にほかならない．
  他の式も同様である．
\end{proof}
\end{leftbar}

つまり，行列の加法については，両辺が定義されるかぎり，結合律および交換律が成り立つ．
次節では行列の積を定義し，積についてもこれらの性質が成り立つかを考えることにする．

\phantomsection
\subsection{行列の積}

\begin{definition}{行列の積}{行列の積}
  $l \times m$行列$A = (a_{ij})$と$m \times n$行列$B = (b_{jk})$に対して，$A$と$B$の\textbf{積}\index{ぎょうれつのせき@行列の積}$AB$を，$(i, k)$成分が
  \[
    c_{ik} = \sum_{j=1}^{m} a_{ij} b_{jk}
    \quad (i = 1, 2, \dots, l,~ k = 1, 2, \dots, n)
  \]
  で与えられる$l \times n$行列$C = (c_{ik})$として定義する．
\end{definition}

\begin{wrapfigure}[7]{r}{0.5\textwidth}
  \centering
  \vspace{-\baselineskip}
  \begin{tikzpicture}[
      mat/.style={matrix of math nodes,
                  left delimiter=(, right delimiter=),
                  nodes={inner sep=1.5pt, font=\scriptsize},
                  column sep=5pt, row sep=4pt},
      waku/.style={draw, magenta, thick, rounded corners=2pt, inner sep=2.5pt}]
    \matrix (A) [mat] {
      a_{11} & \cdots & a_{1m} \\
      \vdots &        & \vdots \\
      |[magenta]| a_{i1} & |[magenta]| \cdots & |[magenta]| a_{im} \\
      a_{l1} & \cdots & a_{lm} \\
    };
    \matrix (B) [mat, right=14pt of A] {
      b_{11} & \cdots & |[magenta]| b_{1k} & \cdots & b_{1n} \\
      \vdots &        & |[magenta]| \vdots &        & \vdots \\
      b_{m1} & \cdots & |[magenta]| b_{mk} & \cdots & b_{mn} \\
    };
    \node[waku, fit=(A-3-1)(A-3-3)] (row) {};
    \node[waku, fit=(B-1-3)(B-3-3)] (col) {};
    \node[magenta, font=\tiny, left=8pt of row] {第$i$行};
    \node[magenta, font=\tiny, above=4pt of col] {第$k$列};
  \end{tikzpicture}
    \caption{行列の積の計算}
  \label{fig:matrix-product}
\end{wrapfigure}

積$AB$が定義されるのは，左の行列$A$の列の数と，右の行列$B$の行の数が一致するときに限ることに注意したい．
そして，$AB$の$(i, k)$成分は，$A$の第$i$行と$B$の第$k$列の成分どうしを順にかけて足し合わせたもの，すなわち
\[
  a_{i1}b_{1k} + a_{i2}b_{2k} + \cdots + a_{im}b_{mk}
\]
である．

\deref{行列の積}は複雑である．
しかし，のちに線型写像について考察すると，この定義が自然であることが明らかとなる\footnote{行列の積が写像の合成に対応することについては，\prref{行列の積と写像の合成}で確認する．}．
ここでは具体例で確認してみよう．

\begin{example}{行列の積}{行列の積}
  \[
    A = \begin{pmatrix} 1 & 2 & 0 \\ 3 & 2 & 0 \end{pmatrix}, \quad
    B = \begin{pmatrix} 1 & 0 \\ 0 & 2 \\ 1 & 1 \end{pmatrix}
  \]
  のとき，$A$は$2 \times 3$行列，$B$は$3 \times 2$行列であるため，積$AB$は定義され，
  \[
    AB =
    \begin{pmatrix}
      1 \cdot 1 + 2 \cdot 0 + 0 \cdot 1 & 1 \cdot 0 + 2 \cdot 2 + 0 \cdot 1 \\
      3 \cdot 1 + 2 \cdot 0 + 0 \cdot 1 & 3 \cdot 0 + 2 \cdot 2 + 0 \cdot 1
    \end{pmatrix}
    =
    \begin{pmatrix}
      1 & 4 \\
      3 & 4
    \end{pmatrix}
  \]
  である．
\end{example}

一般に，積$AB$が定義されても$BA$が定義されるとは限らない．
ただし，$A$が$m\times n$行列，$B$が$n\times m$行列であれば，
$AB$と$BA$はともに定義される．
とくに，同じ次数の正方行列どうしでは，常に$AB$と$BA$の両方が定義される．

任意の$a, b \in \mathbb{R}$に対して，以下の$2$つの事実が成り立つことは既知であろう：
\begin{gather*}
  ab = ba, \\
  ab = 0 \limp (a = 0) \lor (b = 0).
\end{gather*}
つまり，$\mathbb{R}$では積の交換律が成り立ち，$a \ne 0$かつ$b \ne 0$なのに$ab = 0$となるような$a, b$は存在しない．
では，行列の場合はどうであろうか．

\begin{example}{行列の積の非可換性と零因子}{行列の積の非可換性と零因子}
  行列$A$，$B$を
  \[
    A = \begin{pmatrix} 1 & 0 \\ 1 & 0 \end{pmatrix}, \quad
    B = \begin{pmatrix} 0 & 0 \\ 0 & 1 \end{pmatrix}
  \]
  で定めると，
  \[
    AB = \begin{pmatrix} 0 & 0 \\ 0 & 0 \end{pmatrix}, \quad
    BA = \begin{pmatrix} 0 & 0 \\ 1 & 0 \end{pmatrix}
  \]
  となる．
\end{example}

この例から，行列においては，一般に積の交換律$AB = BA$は成立しないことがわかる．
さらに，$A \ne O$かつ$B \ne O$であるのに$AB = O$となることも起こりうる\footnote{このように，零行列でないのに，零行列でない行列をかけると積が零行列になるような行列を\textbf{零因子}\index{れいいんし@零因子}という．}．
このように，行列の積は数の積とは異なる振る舞いをすることがあるので，計算の際には注意が必要である．

一方で，行列の積は数の積と似た性質をもつ場面もある．

\begin{prop}{行列の乗法結合律}{行列の乗法結合律}
  $A$が$l \times m$行列，$B$が$m \times n$行列，$C$が$n \times p$行列であれば，
  \[
    (AB)C = A(BC)
  \]
  が成り立つ．
\end{prop}

\begin{leftbar}
\begin{proof}
  $A = (a_{ij})$，$B = (b_{jk})$，$C = (c_{kq})$とかく．
  $AB$の$(i, k)$成分は
  \[
    \sum_{j=1}^{m} a_{ij} b_{jk}
    \quad (i = 1, 2, \dots, l,~ k = 1, 2, \dots, n)
  \]
  であるから，$(AB)C$の$(i, q)$成分は
  \[
    \sum_{k=1}^{n} \left( \sum_{j=1}^{m} a_{ij} b_{jk} \right) c_{kq}
    = \sum_{k=1}^{n} \sum_{j=1}^{m} a_{ij} b_{jk} c_{kq}
  \]
  である．同様に，$A(BC)$の$(i, q)$成分は
  \[
    \sum_{j=1}^{m} \sum_{k=1}^{n} a_{ij} b_{jk} c_{kq}
  \]
  である．この$2$つの式は，$\sum$記号の順序が違うだけで一致している．
  これが証明すべきことであった．
\end{proof}
\end{leftbar}

交換律は成り立たないが結合律は成り立つ，という場面はほかにもある．
たとえば，写像の合成について，
\[
  f \circ g = g \circ f
\]
は一般には成り立たないが，
\[
  h \circ (g \circ f) = (h \circ g) \circ f
\]
はつねに成り立つ．

\begin{example}{合成関数と交換律}{合成関数と交換律}
  実数全体で定義された関数$f(x) = 1 + x$，$g(x) = x^3$を考える．
  $f$と$g$，$g$と$f$はいずれも合成可能で，
  \[
    (g \circ f)(x) = (1 + x)^3, \quad (f \circ g)(x) = 1 + x^3
  \]
  となり，$g \circ f \ne f \circ g$である．
\end{example}

実は，行列の積と写像の合成のこの類似は偶然ではない．
このことは，のちに線型写像と行列の関係を調べる際に明らかになる（\prref{行列の積と写像の合成}）．

行列の積について成り立つ計算規則をまとめておこう．

\begin{prop}{行列の積の計算規則}{行列の積の計算規則}
  $A$，$B$，$C$を行列とし，$E$を単位行列とする．両辺が定義されるかぎり，次が成り立つ：
  \begin{align*}
    (AB)C &= A(BC) && \text{（積の結合律）} \\
    AE &= A, \quad EA = A && \text{（乗法単位元の存在）} \\
    A(B + C) &= AB + AC, \quad (A + B)C = AC + BC && \text{（分配律）}
  \end{align*}
\end{prop}

\begin{leftbar}
\begin{proof}
  第$1$式は\prref{行列の乗法結合律}で示した通りである．
  第$2$式と第$3$式は，\prref{行列の和の計算規則}の証明と同様に，両辺の成分を比較することで確かめられる．
\end{proof}
\end{leftbar}

\phantomsection
\subsection{逆行列}

数$a \ne 0$に対しては，$ax = xa = 1$を満たす数$x = a^{-1}$が存在した．
この「逆数」にあたる概念を，正方行列に対して考えよう．

\begin{definition}{逆行列}{逆行列}
  $n$次正方行列$A$と$n$次単位行列$E$に対して，
  \[
    AX = XA = E
  \]
  となる$n$次正方行列$X$が存在するとき，$A$は\textbf{正則}\index{せいそく@正則}であるという．
  このような$X$を$A$の\textbf{逆行列}\index{ぎゃくぎょうれつ@逆行列}とよび，$A^{-1}$とかく．
\end{definition}

\begin{prop}{逆行列の一意性}{逆行列の一意性}
  正方行列$A$に逆行列が存在するとき，それはただ一つに限る．
\end{prop}

\begin{leftbar}
\begin{proof}
  $X$，$Y$がともに$A$の逆行列であるとする．このとき，
  \[
    AY = YA = E
  \]
  が成り立つ．このもとで，
  \[
    X = XE = X(AY) = (XA)Y = EY = Y
  \]
  となるため，逆行列の一意性が示された．
\end{proof}
\end{leftbar}

\begin{remark}[定義は最小限にとどめる]
  逆行列は，定義から一意であることがわかった．
  ここで，「常に逆行列が一意であるならば，逆行列の定義に一意性を入れたほうが定義の情報量が多くなってよいはず」と考える方もおられるだろう．

  しかし，数学においては「定義は最小限の情報におさえて，その他のことは命題として示す」という慣習がある．
  定義を最小限にしたほうが，「逆行列が一意であることは，定義から証明できる事実である」ということが明確になるのである．
\end{remark}

\begin{example}{2次正方行列の逆行列}{2次正方行列の逆行列}
  $A = \begin{pmatrix} a & b \\ c & d \end{pmatrix}$（成分は$K$の元）に対して，$AX = E$となる行列$X$が存在する条件を求めよう．
  $X = \begin{pmatrix} x_{11} & x_{12} \\ x_{21} & x_{22} \end{pmatrix}$とすると，
  \[
    AX =
    \begin{pmatrix}
      ax_{11} + bx_{21} & ax_{12} + bx_{22} \\
      cx_{11} + dx_{21} & cx_{12} + dx_{22}
    \end{pmatrix}
  \]
  であり，これが$E$と等しくなればよいので，
  \[
    \begin{cases}
      ax_{11} + bx_{21} = 1 \\
      cx_{11} + dx_{21} = 0
    \end{cases}, \quad
    \begin{cases}
      ax_{12} + bx_{22} = 0 \\
      cx_{12} + dx_{22} = 1
    \end{cases}
  \]
  を得る．$ad - bc \ne 0$のとき，この$2$組の連立方程式は加減法によりただ一組の解をもち，
  \[
    x_{11} = \frac{d}{ad - bc}, \quad
    x_{12} = -\frac{b}{ad - bc}, \quad
    x_{21} = -\frac{c}{ad - bc}, \quad
    x_{22} = \frac{a}{ad - bc}
  \]
  となる．この$X$に対して$XA$を計算すると，たしかに$XA = E$となる．
  よって，$ad - bc \ne 0$のとき$A$は正則であり，その逆行列は
  \[
    A^{-1} = \frac{1}{ad - bc} \begin{pmatrix} d & -b \\ -c & a \end{pmatrix}
  \]
  である\footnote{逆に$ad - bc = 0$のときは$A$は正則でないが，このことは上の計算（$ad - bc \ne 0$の場合の解の構成）だけからは直ちには従わない．\partref{行列式}で行列式を導入したあとに，一般の$n$次正方行列が正則であるための必要十分条件として証明される．}．
\end{example}

\phantomsection
\section{連立一次方程式}

\phantomsection
\subsection{加減法による解法}

次のような連立一次方程式を考える：
\begin{equation}
  \begin{cases}
    4x + 3y = 8 \\
    2x + 3y = 4
  \end{cases} \label{eq:連立一次方程式の例}
\end{equation}

初等的な解法を確認しておこう．
第$1$式から第$2$式を辺々引くと$2x = 4$を得るので，$x = 2$である．
これを第$2$式に代入すると$4 + 3y = 4$，すなわち$3y = 0$となるので，$y = 0$である．
よって，\eqref{eq:連立一次方程式の例}の解は$x = 2$，$y = 0$である．

この解法は「加減法」としてよく知られているものである．
本節では，この解法を行列のことばで表現し直すことを考える．

\phantomsection
\subsection{係数行列と行基本変形}

一般に，$m$個の式と$n$個の未知数からなる連立一次方程式
\[
  \begin{cases}
    a_{11}x_1 + a_{12}x_2 + \cdots + a_{1n}x_n = b_1 \\
    a_{21}x_1 + a_{22}x_2 + \cdots + a_{2n}x_n = b_2 \\
    \hspace{2em} \vdots \\
    a_{m1}x_1 + a_{m2}x_2 + \cdots + a_{mn}x_n = b_m
  \end{cases}
\]
は，行列の積を用いて
\[
  A\bm{x} = \bm{b}, \quad
  A = (a_{ij}), \quad
  \bm{x} = \begin{pmatrix} x_1 \\ x_2 \\ \vdots \\ x_n \end{pmatrix}, \quad
  \bm{b} = \begin{pmatrix} b_1 \\ b_2 \\ \vdots \\ b_m \end{pmatrix}
\]
とかける．

\begin{definition}{係数行列・拡大係数行列}{係数行列}
  上の連立一次方程式に対して，係数を並べた$m \times n$行列$A = (a_{ij})$を\textbf{係数行列}\index{けいすうぎょうれつ@係数行列}という．
  さらに，$A$の右側に$\bm{b}$を付け加えた$m \times (n + 1)$行列
  \[
    \begin{pmatrix}
      a_{11} & a_{12} & \cdots & a_{1n} & b_1 \\
      a_{21} & a_{22} & \cdots & a_{2n} & b_2 \\
      \vdots & \vdots & \ddots & \vdots & \vdots \\
      a_{m1} & a_{m2} & \cdots & a_{mn} & b_m
    \end{pmatrix}
  \]
  を\textbf{拡大係数行列}\index{かくだいけいすうぎょうれつ@拡大係数行列}という．
\end{definition}

たとえば，\eqref{eq:連立一次方程式の例}は
\[
  \begin{pmatrix} 4 & 3 \\ 2 & 3 \end{pmatrix}
  \begin{pmatrix} x \\ y \end{pmatrix}
  =
  \begin{pmatrix} 8 \\ 4 \end{pmatrix}
\]
とかけ，係数行列は$\begin{pmatrix} 4 & 3 \\ 2 & 3 \end{pmatrix}$，拡大係数行列は$\begin{pmatrix} 4 & 3 & 8 \\ 2 & 3 & 4 \end{pmatrix}$である．

\begin{definition}{行基本変形}{行基本変形}
  行列に対する以下の$3$種類の操作を\textbf{行基本変形}\index{ぎょうきほんへんけい@行基本変形}という．
  \begin{enumerate}[label=(\arabic*),leftmargin=3.2em]
    \item 第$i$行を$c$倍する（$c \ne 0$）\footnote{$c = 0$倍を許すと，第$i$行の情報が失われてしまい，もとの連立方程式と同値でなくなる．そのため$c \ne 0$という条件を課している．}．
    \item 第$i$行と第$j$行を入れ替える．
    \item 第$i$行に第$j$行の定数倍を加える（$i \ne j$）．
  \end{enumerate}
\end{definition}

拡大係数行列に行基本変形を施すことは，連立方程式に対して「両辺を定数倍する」「式の順序を入れ替える」「ある式に別の式の定数倍を加える」という操作を行うことに対応する．
いずれの操作でも連立方程式の解は変わらないので，行基本変形を繰り返して拡大係数行列を簡単な形に変形すれば，解が読み取れる．

たとえば，\eqref{eq:連立一次方程式の例}の拡大係数行列は，以下のように変形できる：
\begin{align*}
  \begin{pmatrix}
    4 & 3 & 8 \\
    2 & 3 & 4
  \end{pmatrix}
  &\xrightarrow{\text{第1行に第2行の$(-1)$倍を加える}}
  \begin{pmatrix}
    2 & 0 & 4 \\
    2 & 3 & 4
  \end{pmatrix} \\
  &\xrightarrow{\text{第2行に第1行の$(-1)$倍を加える}}
  \begin{pmatrix}
    2 & 0 & 4 \\
    0 & 3 & 0
  \end{pmatrix} \\
  &\xrightarrow{\text{第1行を$1/2$倍，第2行を$1/3$倍する}}
  \begin{pmatrix}
    1 & 0 & 2 \\
    0 & 1 & 0
  \end{pmatrix}
\end{align*}
最後の行列は連立方程式
\[
  \begin{cases}
    x = 2 \\
    y = 0
  \end{cases}
\]
を表しており，これがそのまま解を与えている．
この操作は，本質的には加減法による解法と同じである．

\phantomsection
\subsection{逆行列を用いた解法}

係数行列が正則な場合には，逆行列を用いて解を表すこともできる．
\eqref{eq:連立一次方程式の例}を
\begin{equation}
  A\bm{x} = \bm{b}, \quad
  A = \begin{pmatrix} 4 & 3 \\ 2 & 3 \end{pmatrix}, \quad
  \bm{x} = \begin{pmatrix} x \\ y \end{pmatrix}, \quad
  \bm{b} = \begin{pmatrix} 8 \\ 4 \end{pmatrix} \label{eq:連立一次方程式の行列表示}
\end{equation}
とかく．$4 \cdot 3 - 3 \cdot 2 = 6 \ne 0$なので，\exref{2次正方行列の逆行列}より$A$は正則であり，
\[
  A^{-1} = \frac{1}{6} \begin{pmatrix} 3 & -3 \\ -2 & 4 \end{pmatrix}
\]
である．これを\eqref{eq:連立一次方程式の行列表示}の両辺に左からかけると，
\begin{gather*}
  A^{-1}A\bm{x} = A^{-1}\bm{b} \\
  \therefore~ \bm{x} = A^{-1}\bm{b}
\end{gather*}
であり，
\[
  A^{-1}\bm{b} = \frac{1}{6}
  \begin{pmatrix} 3 & -3 \\ -2 & 4 \end{pmatrix}
  \begin{pmatrix} 8 \\ 4 \end{pmatrix}
  =
  \begin{pmatrix} 2 \\ 0 \end{pmatrix}
\]
であるから，与えられた連立一次方程式の解は$x = 2$，$y = 0$とわかる．
これは加減法および行基本変形で得た解と一致している．
ここまでで，行列の演算と，連立一次方程式の行列による扱いを整理した．
次章では，行列と密接に関わる概念である線型写像を導入し，行列がなにを表しているのかを見直そう．

\phantomsection
\part{線型写像}

\begin{abstract}
  この章では，行列と密接に関わる概念である\textbf{線型写像}を定義する．
  まず一般の写像を定義したうえで，「和とスカラー倍を保つ」写像として線型写像を定義し，
  対称移動・正射影・回転といった幾何的な変換から，多項式の微分・積分や数列の階差・部分和にいたるまで，
  さまざまな例を確認する．
  次に，$K^n$から$K^m$への線型写像が$m \times n$行列とちょうど一対一に対応することを示し，
  行列の積が写像の合成に対応していることを見る．
  最後に，線型写像に付随する基本的な集合として像と核を導入し，
  それらが連立一次方程式のことばで言い換えられることを確かめる．
\end{abstract}

\phantomsection
\section{線型写像の定義と例}

\phantomsection
\subsection{写像と線型写像}

本節では，行列と密接に関わる概念である\textbf{線型写像}\index{せんけいしゃぞう@線型写像}を定義する．
その準備として，まず一般の\textbf{写像}\index{しゃぞう@写像}を定義しよう．

\begin{definition}{写像}{写像}
  $V$，$W$を集合とする．
  対応$f \colon V \to W$が\textbf{写像}であるとは，すべての$x \in V$に対して，$f(x) = y$となる$y \in W$がただひとつ存在することである．
  このとき，$V$を$f$の\textbf{定義域}\index{ていぎいき@定義域}，$W$を$f$の\textbf{終域}\index{しゅういき@終域}という．
\end{definition}

写像のうち，「和とスカラー倍を保つ」という性質をもつものを線型写像とよぶ．

\begin{definition}{線型写像・線型変換}{線型写像}
  写像$f \colon K^n \to K^m$が，任意の$\bm{x}, \bm{y} \in K^n$と任意の$\alpha \in K$に対して
  \[
    \begin{cases}
      f(\bm{x} + \bm{y}) = f(\bm{x}) + f(\bm{y}) \\
      f(\alpha\bm{x}) = \alpha f(\bm{x})
    \end{cases}
  \]
  を満たすとき，$f$を$K^n$から$K^m$への\textbf{線型写像}という．
  とくに$m = n$のとき，$f$を$K^n$の\textbf{線型変換}\index{せんけいへんかん@線型変換}という．
\end{definition}

\deref{線型写像}に現れる$2$つの条件は，次のひとつの条件にまとめることもできる：
任意の$\bm{x}, \bm{y} \in K^n$と任意の$\alpha, \beta \in K$に対して，
\[
  f(\alpha\bm{x} + \beta\bm{y}) = \alpha f(\bm{x}) + \beta f(\bm{y})
\]
が成り立つ\footnote{実際，この条件で$\alpha = \beta = 1$とすれば第$1$の条件が，$\beta = 0$とすれば第$2$の条件が得られる．逆に$2$つの条件を順に用いれば$f(\alpha\bm{x} + \beta\bm{y}) = f(\alpha\bm{x}) + f(\beta\bm{y}) = \alpha f(\bm{x}) + \beta f(\bm{y})$となる．}．

定義だけでは掴みにくいので，最も簡単な具体例を確かめておこう．

\begin{example}{線型変換}{線型変換}
  $a \in K$を定数とし，写像$f \colon K \to K$を
  \[
    f(x) = ax
  \]
  で定義する．このとき，$f$は$K$の線型変換である．
  実際，任意の$x, y \in K$と$\alpha \in K$に対して
  \begin{align*}
    f(x + y) &= a(x + y) = ax + ay = f(x) + f(y), \\
    f(\alpha x) &= a(\alpha x) = \alpha(ax) = \alpha f(x)
  \end{align*}
  が成り立つからである．
\end{example}

この例は$n = m = 1$の場合であった．
以下，$K^n$から$K^m$への線型写像の例を，いくつか具体的に見ていこう．
とくに断らない限り，本節の幾何的な例では$K = \mathbb{R}$とし，$\mathbb{R}^2$の元
\[
  \bm{x} = \begin{pmatrix} x \\ y \end{pmatrix}
\]
を座標平面上の点$(x, y)$と同一視する．

\begin{example}{恒等変換と零写像}{恒等変換と零写像}
  写像$\mathrm{id} \colon K^n \to K^n$を$\mathrm{id}(\bm{x}) = \bm{x}$で定めると，
  \[
    \mathrm{id}(\bm{x} + \bm{y}) = \bm{x} + \bm{y} = \mathrm{id}(\bm{x}) + \mathrm{id}(\bm{y}), \quad
    \mathrm{id}(\alpha\bm{x}) = \alpha\bm{x} = \alpha\,\mathrm{id}(\bm{x})
  \]
  であるから，$\mathrm{id}$は線型変換である．これを\textbf{恒等変換}\index{こうとうへんかん@恒等変換}という．

  また，すべての$\bm{x} \in K^n$に対して$f(\bm{x}) = \bm{0}$と定めた写像$f \colon K^n \to K^m$も，
  \[
    f(\bm{x} + \bm{y}) = \bm{0} = \bm{0} + \bm{0} = f(\bm{x}) + f(\bm{y}), \quad
    f(\alpha\bm{x}) = \bm{0} = \alpha\bm{0} = \alpha f(\bm{x})
  \]
  より線型写像である．これを\textbf{零写像}\index{れいしゃぞう@零写像}という．
\end{example}

\begin{example}{相似の中心が原点である拡大・縮小}{相似変換}
  $c \in K$を定数とし，写像$f \colon K^n \to K^n$を
  \[
    f(\bm{x}) = c\bm{x}
  \]
  で定める．このとき，
  \[
    f(\bm{x} + \bm{y}) = c(\bm{x} + \bm{y}) = c\bm{x} + c\bm{y} = f(\bm{x}) + f(\bm{y}), \quad
    f(\alpha\bm{x}) = c(\alpha\bm{x}) = \alpha(c\bm{x}) = \alpha f(\bm{x})
  \]
  が成り立つので，$f$は線型変換である．
  $K = \mathbb{R}$，$n = 2$のときこれは，原点を中心とする相似比$c$の拡大・縮小にあたる．
  とくに$c = 1$のときが恒等変換，$c = 0$のときが零写像，$c = -1$のときが原点に関する対称移動である．
\end{example}

次に，座標平面上の図形的な変換の例を見よう．いずれも高校数学で図形の移動として扱われるものである．

\begin{example}{対称移動と正射影}{対称移動と正射影}
  $K = \mathbb{R}$とし，写像$\mathbb{R}^2 \to \mathbb{R}^2$を次のように定める：
  \[
    f_1 \begin{pmatrix} x \\ y \end{pmatrix} = \begin{pmatrix} x \\ -y \end{pmatrix}, \quad
    f_2 \begin{pmatrix} x \\ y \end{pmatrix} = \begin{pmatrix} y \\ x \end{pmatrix}, \quad
    f_3 \begin{pmatrix} x \\ y \end{pmatrix} = \begin{pmatrix} x \\ 0 \end{pmatrix}.
  \]
  $f_1$は$x$軸に関する対称移動，$f_2$は直線$y = x$に関する対称移動，$f_3$は$x$軸への正射影（$x$軸に下ろした垂線の足への対応）である．
  これらはいずれも線型変換である．

  たとえば$f_1$について確かめよう．
  $\bm{x} = \begin{pmatrix} x_1 \\ x_2 \end{pmatrix}$，$\bm{y} = \begin{pmatrix} y_1 \\ y_2 \end{pmatrix}$，$\alpha \in \mathbb{R}$に対して，
  \[
    f_1(\bm{x} + \bm{y})
    = \begin{pmatrix} x_1 + y_1 \\ -(x_2 + y_2) \end{pmatrix}
    = \begin{pmatrix} x_1 \\ -x_2 \end{pmatrix} + \begin{pmatrix} y_1 \\ -y_2 \end{pmatrix}
    = f_1(\bm{x}) + f_1(\bm{y}),
  \]
  \[
    f_1(\alpha\bm{x})
    = \begin{pmatrix} \alpha x_1 \\ -\alpha x_2 \end{pmatrix}
    = \alpha \begin{pmatrix} x_1 \\ -x_2 \end{pmatrix}
    = \alpha f_1(\bm{x})
  \]
  が成り立つ．$f_2$，$f_3$についても同様である．
\end{example}

\begin{example}{原点のまわりの回転}{原点のまわりの回転}
  $K = \mathbb{R}$とし，$\theta$を定数とする．
  座標平面上の点を原点のまわりに角$\theta$だけ回転させる対応は，
  \[
    f\begin{pmatrix} x \\ y \end{pmatrix}
    = \begin{pmatrix} x\cos\theta - y\sin\theta \\ x\sin\theta + y\cos\theta \end{pmatrix}
  \]
  と表される\footnote{点$(x, y)$を極形式で$x = r\cos\varphi$，$y = r\sin\varphi$とかけば，回転後の点は$(r\cos(\varphi + \theta),\ r\sin(\varphi + \theta))$である．三角関数の加法定理を用いて展開すれば，上の式が得られる．}．
  この$f$は$\mathbb{R}^2$の線型変換である．実際，
  \[
    f\left( \begin{pmatrix} x_1 \\ x_2 \end{pmatrix} + \begin{pmatrix} y_1 \\ y_2 \end{pmatrix} \right)
    = \begin{pmatrix} (x_1 + y_1)\cos\theta - (x_2 + y_2)\sin\theta \\ (x_1 + y_1)\sin\theta + (x_2 + y_2)\cos\theta \end{pmatrix}
    = f\begin{pmatrix} x_1 \\ x_2 \end{pmatrix} + f\begin{pmatrix} y_1 \\ y_2 \end{pmatrix}
  \]
  であり，また$\alpha \in \mathbb{R}$に対して
  \[
    f\left( \alpha \begin{pmatrix} x_1 \\ x_2 \end{pmatrix} \right)
    = \begin{pmatrix} \alpha x_1 \cos\theta - \alpha x_2 \sin\theta \\ \alpha x_1 \sin\theta + \alpha x_2 \cos\theta \end{pmatrix}
    = \alpha f\begin{pmatrix} x_1 \\ x_2 \end{pmatrix}
  \]
  が成り立つからである．

  なお，この事実は幾何的にも納得できる．
  和$\bm{x} + \bm{y}$は平行四辺形の対角線として得られるが，平行四辺形全体を回転させても平行四辺形のままであり，
  対角線は対角線にうつる．これが$f(\bm{x} + \bm{y}) = f(\bm{x}) + f(\bm{y})$の意味するところである．
\end{example}

高校で学ぶベクトルの内積も，線型写像の例を与える．

\begin{example}{内積が定める線型写像}{内積が定める線型写像}
  $K = \mathbb{R}$とし，$\bm{a} = \begin{pmatrix} a_1 \\ a_2 \end{pmatrix} \in \mathbb{R}^2$を定ベクトルとする．
  写像$f \colon \mathbb{R}^2 \to \mathbb{R}$を，内積を用いて
  \[
    f(\bm{x}) = \bm{a} \cdot \bm{x} = a_1 x_1 + a_2 x_2
  \]
  で定めると，$f$は線型写像である．
  これは，高校で学ぶ内積の性質
  \[
    \bm{a} \cdot (\bm{x} + \bm{y}) = \bm{a} \cdot \bm{x} + \bm{a} \cdot \bm{y}, \qquad
    \bm{a} \cdot (\alpha\bm{x}) = \alpha (\bm{a} \cdot \bm{x})
  \]
  そのものにほかならない．

  さらに，$\bm{u}$を大きさ$1$のベクトルとするとき，$\bm{x}$から$\bm{u}$の定める直線への正射影は
  \[
    g(\bm{x}) = (\bm{u} \cdot \bm{x})\,\bm{u}
  \]
  で与えられる．内積の上の性質より，
  \[
    g(\bm{x} + \bm{y}) = \bigl(\bm{u} \cdot (\bm{x} + \bm{y})\bigr)\bm{u}
    = (\bm{u} \cdot \bm{x})\bm{u} + (\bm{u} \cdot \bm{y})\bm{u} = g(\bm{x}) + g(\bm{y})
  \]
  となり，スカラー倍についても同様であるから，$g$は$\mathbb{R}^2$の線型変換である．
  $\bm{u} = \bm{e}_1$とすれば，\exref{対称移動と正射影}の$f_3$が得られる．
\end{example}

線型性の定義から直ちに従う，簡単だが有用な性質を注意しておこう．

\begin{prop}{線型写像は原点を原点にうつす}{線型写像と零ベクトル}
  $f \colon K^n \to K^m$が線型写像ならば，
  \[
    f(\bm{0}) = \bm{0}
  \]
  が成り立つ．
\end{prop}

\begin{leftbar}
\begin{proof}
  \deref{線型写像}の第$2$の条件で$\alpha = 0$，$\bm{x} = \bm{0}$とすれば，
  \[
    f(\bm{0}) = f(0 \cdot \bm{0}) = 0 \cdot f(\bm{0}) = \bm{0}
  \]
  が成り立つ．これが証明すべきことであった．
\end{proof}
\end{leftbar}

\prref{線型写像と零ベクトル}の対偶をとれば，「$f(\bm{0}) \ne \bm{0}$ならば$f$は線型写像でない」となる．
これは，与えられた写像が線型でないことを示す際の，もっとも手軽な判定法である．

ただし，$f(\bm{0}) = \bm{0}$であることだけから，$f$が線型写像であるとは限らない．
線型写像でないことを示すには，線型写像の定義に現れる条件のうち，いずれかが成り立たないことを示せばよい．
以下では，いくつかの具体例を通して，線型写像でないことを確かめてみよう．

\begin{example}{線型写像でない写像}{線型写像でない写像}
  以下，$K = \mathbb{R}$とする．
  \begin{enumerate}[label=(\arabic*),leftmargin=3.2em]
    \item （平行移動）$\bm{b} \ne \bm{0}$を定ベクトルとし，$f \colon \mathbb{R}^2 \to \mathbb{R}^2$を$f(\bm{x}) = \bm{x} + \bm{b}$で定めると，
      $f(\bm{0}) = \bm{b} \ne \bm{0}$であるから，\prref{線型写像と零ベクトル}より$f$は線型写像でない．
    \item （$2$乗）$f \colon \mathbb{R} \to \mathbb{R}$を$f(x) = x^2$で定めると，
      \[
        f(1 + 1) = 4, \quad f(1) + f(1) = 2
      \]
      であり，$f(x + y) = f(x) + f(y)$が成り立たない．よって$f$は線型写像でない．
    \item （絶対値）$f \colon \mathbb{R} \to \mathbb{R}$を$f(x) = |x|$で定めると，$f(0) = 0$は成り立つが，
      \[
        f\bigl((-1) \cdot 1\bigr) = 1, \quad (-1) \cdot f(1) = -1
      \]
      であり，$f(\alpha x) = \alpha f(x)$が成り立たない．よって$f$は線型写像でない．
    \item （成分の積）$f \colon \mathbb{R}^2 \to \mathbb{R}$を$f\begin{pmatrix} x \\ y \end{pmatrix} = xy$で定めると，$f(\bm{0}) = 0$であるが，
      $\bm{x} = \begin{pmatrix} 1 \\ 1 \end{pmatrix}$に対して
      \[
        f(2\bm{x}) = 2 \cdot 2 = 4, \quad 2 f(\bm{x}) = 2 \cdot 1 = 2
      \]
      となり，やはり線型写像でない．
  \end{enumerate}
  (3)，(4)からわかるように，$f(\bm{0}) = \bm{0}$は線型写像であるための必要条件にすぎず，十分条件ではない．
\end{example}

\begin{mycolumn}
  高校数学では$y = ax + b$を\textbf{一次関数}とよぶ．
  しかし，$b \ne 0$のとき，この対応$f(x) = ax + b$は\prref{線型写像と零ベクトル}より線型写像ではない．
  「一次」という言葉と「線型」という言葉は，ここでずれるのである．

  グラフでいえば，線型写像$\mathbb{R} \to \mathbb{R}$に対応するのは「原点を通る直線」だけであり，
  切片$b$の分だけ平行移動された直線は線型写像を与えない．
  一般に，線型写像$f$と定ベクトル$\bm{b}$を用いて
  \[
    g(\bm{x}) = f(\bm{x}) + \bm{b}
  \]
  と表される写像を\textbf{アフィン写像}\index{あふぃんしゃぞう@アフィン写像}という．
  一次関数は，線型写像そのものではなく，アフィン写像の例なのである．

  なお，$b \ne 0$のときの$f(x) = ax + b$も，$1$次元分の「下駄」を履かせて
  \[
    \begin{pmatrix} x \\ 1 \end{pmatrix} \longmapsto \begin{pmatrix} ax + b \\ 1 \end{pmatrix}
  \]
  という$\mathbb{R}^2$の線型変換とみなすことができる．
  コンピュータグラフィックスで平行移動を行列で扱う際には，この考え方（同次座標）が用いられる．
\end{mycolumn}

\phantomsection
\subsection{さまざまな線型写像}

幾何以外にも，高校数学のなかには線型写像が数多く隠れている．

\begin{example}{合計と平均}{合計と平均}
  $n$個のデータ$x_1, x_2, \dots, x_n$を$\bm{x} \in \mathbb{R}^n$とみなし，写像$S, M \colon \mathbb{R}^n \to \mathbb{R}$を
  \[
    S(\bm{x}) = x_1 + x_2 + \cdots + x_n, \qquad
    M(\bm{x}) = \frac{x_1 + x_2 + \cdots + x_n}{n}
  \]
  で定める．すなわち$S$は合計，$M$は平均である．
  このとき$S$，$M$はいずれも線型写像である．
  「$2$組のデータを足し合わせたものの平均は，それぞれの平均の和に等しい」「すべてのデータを$\alpha$倍すれば平均も$\alpha$倍になる」という，
  データの分析でおなじみの性質が，そのまま線型性である．

  一方，分散
  \[
    V(\bm{x}) = \frac{1}{n}\sum_{i=1}^{n} \bigl(x_i - M(\bm{x})\bigr)^2
  \]
  は線型写像ではない．実際，$V(\alpha\bm{x}) = \alpha^2 V(\bm{x})$であって$\alpha V(\bm{x})$ではない．
\end{example}

\begin{example}{多項式の微分と定積分}{多項式の微分と定積分}
  $2$次以下の多項式$a_0 + a_1 x + a_2 x^2$を，その係数を並べたベクトル
  \[
    \begin{pmatrix} a_0 \\ a_1 \\ a_2 \end{pmatrix} \in \mathbb{R}^3
  \]
  と同一視する．同様に，$1$次以下の多項式$b_0 + b_1 x$を$\begin{pmatrix} b_0 \\ b_1 \end{pmatrix} \in \mathbb{R}^2$と同一視する．

  多項式を微分する操作は，
  \[
    (a_0 + a_1 x + a_2 x^2)' = a_1 + 2a_2 x
  \]
  であるから，この同一視のもとで写像
  \[
    D \colon \mathbb{R}^3 \to \mathbb{R}^2, \quad
    D\begin{pmatrix} a_0 \\ a_1 \\ a_2 \end{pmatrix} = \begin{pmatrix} a_1 \\ 2a_2 \end{pmatrix}
  \]
  を定める．$D$は線型写像であり，これは高校で学ぶ微分の性質
  \[
    (f + g)' = f' + g', \qquad (\alpha f)' = \alpha f'
  \]
  にほかならない．

  同様に，定積分
  \[
    \int_0^1 (a_0 + a_1 x + a_2 x^2)\,dx = a_0 + \frac{a_1}{2} + \frac{a_2}{3}
  \]
  を対応させる写像$I \colon \mathbb{R}^3 \to \mathbb{R}$も線型写像である．
  これは，積分の線型性
  \[
    \int_0^1 \bigl(f(x) + g(x)\bigr) dx = \int_0^1 f(x)\,dx + \int_0^1 g(x)\,dx, \qquad
    \int_0^1 \alpha f(x)\,dx = \alpha \int_0^1 f(x)\,dx
  \]
  の言い換えである．
\end{example}

数列もまた，線型写像の格好の題材である．
以下では，有限個の項からなる数列$a_1, a_2, \dots, a_n$を，そのままベクトル
\[
  \bm{a} = \begin{pmatrix} a_1 \\ a_2 \\ \vdots \\ a_n \end{pmatrix} \in K^n
\]
とみなす．

\begin{example}{階差数列をとる写像}{階差数列をとる写像}
  写像$\Delta \colon K^n \to K^{n-1}$を，階差数列を対応させる写像
  \[
    \Delta \bm{a} =
    \begin{pmatrix}
      a_2 - a_1 \\ a_3 - a_2 \\ \vdots \\ a_n - a_{n-1}
    \end{pmatrix}
  \]
  で定める．このとき$\Delta$は線型写像である．
  実際，第$i$成分だけを見れば
  \[
    (a_{i+1} + b_{i+1}) - (a_i + b_i) = (a_{i+1} - a_i) + (b_{i+1} - b_i), \qquad
    \alpha a_{i+1} - \alpha a_i = \alpha (a_{i+1} - a_i)
  \]
  であり，これがすべての$i$について成り立つからである．

  「$2$つの数列を足してから階差をとっても，それぞれ階差をとってから足しても同じ」という，
  高校で意識せずに使っている事実が，この写像の線型性にほかならない．
\end{example}

\begin{example}{部分和をとる写像}{部分和をとる写像}
  写像$\Sigma \colon K^n \to K^n$を，初項から第$k$項までの和
  \[
    S_k = a_1 + a_2 + \cdots + a_k \quad (1 \le k \le n)
  \]
  を並べたベクトルを対応させる写像
  \[
    \Sigma \bm{a} = \begin{pmatrix} S_1 \\ S_2 \\ \vdots \\ S_n \end{pmatrix}
  \]
  で定める．各$S_k$は$a_1, \dots, a_n$の線型結合であるから，$\Sigma$は線型変換である．
  これは，和の記号についての性質
  \[
    \sum_{i=1}^{k} (a_i + b_i) = \sum_{i=1}^{k} a_i + \sum_{i=1}^{k} b_i, \qquad
    \sum_{i=1}^{k} \alpha a_i = \alpha \sum_{i=1}^{k} a_i
  \]
  の言い換えである．

  なお，$\Delta$が微分の，$\Sigma$が積分の離散版にあたることに注意しよう．
  高校で学ぶ関係式$a_n = S_n - S_{n-1}$は，「部分和をとってから階差をとると元に戻る」ことを述べており，
  微分積分学の基本定理の数列版とみることができる．
\end{example}

\begin{example}{等差数列を作る写像}{等差数列を作る写像}
  初項$a$と公差$d$の組$\begin{pmatrix} a \\ d \end{pmatrix} \in K^2$に対して，
  等差数列の第$n$項までを並べたベクトルを対応させる写像$\varphi \colon K^2 \to K^n$を
  \[
    \varphi \begin{pmatrix} a \\ d \end{pmatrix}
    = \begin{pmatrix} a \\ a + d \\ a + 2d \\ \vdots \\ a + (n-1)d \end{pmatrix}
  \]
  で定める．各成分は$a$と$d$の線型結合であるから，$\varphi$は線型写像である．

  すなわち，「初項どうし・公差どうしを足した等差数列は，もとの$2$つの等差数列を項ごとに足したものに等しい」ということである．

  一方，初項$a$と公比$r$の組から等比数列を作る写像
  \[
    \psi \begin{pmatrix} a \\ r \end{pmatrix} = \begin{pmatrix} a \\ ar \\ ar^2 \end{pmatrix}
  \]
  は線型写像ではない．実際，$\bm{x} = \begin{pmatrix} 1 \\ 2 \end{pmatrix}$に対して
  \[
    \psi(2\bm{x}) = \begin{pmatrix} 2 \\ 8 \\ 32 \end{pmatrix}, \qquad
    2\psi(\bm{x}) = \begin{pmatrix} 2 \\ 4 \\ 8 \end{pmatrix}
  \]
  であり，一致しない．
  公比が数列の「作られ方」そのものに関わるため，公比を変数とみるかぎり線型性は失われるのである．
  ただし，公比$r$を固定して初項だけを動かす写像$a \mapsto (a, ar, ar^2)$は線型写像である．
\end{example}

\begin{example}{数列に関する線型でない操作}{数列と線型でない操作}
次の$2$つは，数列に対するごく自然な操作であるが，線型写像ではない．
\begin{enumerate}[label=(\arabic*),leftmargin=3.2em]
\item 各項を$2$乗する写像$\bm{a} \mapsto (a_1^2, a_2^2, \dots, a_n^2)$．
$\alpha$倍すると各成分は$\alpha^2$倍になるので，$f(\alpha\bm{a}) = \alpha f(\bm{a})$が成り立たない．
\item 隣り合う項の比$\bm{a} \mapsto (a_2/a_1, \dots, a_n/a_{n-1})$．
そもそも$a_i = 0$のとき定義されず，$K^n$全体で定義された写像にすらならない．
\end{enumerate}
線型写像は加法とスカラー倍を保つため，このような乗法や除法を含む操作とは一般に相性がよくない．
等差数列（差が一定）が線型な世界となじみ，等比数列（比が一定）がなじまないのは，このあたりの事情による．
\end{example}

\begin{mycolumn}
  \exref{多項式の微分と定積分}では多項式を係数の並びとみなし，本項の数列の例では数列を項の並びとみなすことで，それぞれ数ベクトルに翻訳した．
  同様の翻訳は，関数などさまざまな対象に対して可能である．
  \partref{線型空間}では，このような「翻訳」を経由せずに，多項式全体や関数全体をそのまま扱える枠組みとして\textbf{線型空間}を導入する．
  そこでは，微分や積分は線型空間のあいだの線型写像として，改めて位置づけられることになる．
\end{mycolumn}

\phantomsection
\subsection{線型写像の和・スカラー倍・合成}

次に，線型写像から新しい線型写像をつくる操作として，和・スカラー倍・合成を考える．

\begin{definition}{線型写像の和とスカラー倍}{線型写像の和とスカラー倍}
  $f, g \colon K^n \to K^m$を線型写像とする．
  このとき，
  \[
    h(\bm{x}) = f(\bm{x}) + g(\bm{x})
  \]
  で定まる写像$h \colon K^n \to K^m$を$f$と$g$の\textbf{和}\index{せんけいしゃぞうのわ@線型写像の和}とよび，$f + g$と表す．

  また，$\alpha \in K$と線型写像$f \colon K^n \to K^m$に対して，
  \[
    h(\bm{x}) = \alpha \cdot f(\bm{x})
  \]
  で定まる写像$h \colon K^n \to K^m$を$f$の\textbf{$\alpha$倍}\index{せんけいしゃぞうのすからーばい@線型写像のスカラー倍}とよび，$\alpha f$と表す．
\end{definition}

\begin{prop}{線型写像の和とスカラー倍の線型性}{線型写像の和とスカラー倍の線型性}
  $f, g \colon K^n \to K^m$が線型写像ならば，$f + g \colon K^n \to K^m$も線型写像である．
  また，$\alpha \in K$に対して，$\alpha f \colon K^n \to K^m$も線型写像である．
\end{prop}

\begin{leftbar}
\begin{proof}
  前半の主張を示す．任意の$\bm{x}, \bm{y} \in K^n$に対して，
  \begin{align}
    (f + g)(\bm{x} + \bm{y}) &= f(\bm{x} + \bm{y}) + g(\bm{x} + \bm{y}) \tag{$f+g$の定義} \\
    &= f(\bm{x}) + f(\bm{y}) + g(\bm{x}) + g(\bm{y}) \tag{$f,g$の線型性} \\
    &= \bigl(f(\bm{x}) + g(\bm{x})\bigr) + \bigl(f(\bm{y}) + g(\bm{y})\bigr) \nonumber \\
    &= (f + g)(\bm{x}) + (f + g)(\bm{y}) \tag{$f+g$の定義}
  \end{align}
  が成り立ち，さらに任意の$\alpha \in K$に対して，
  \begin{align}
    (f + g)(\alpha\bm{x}) &= f(\alpha\bm{x}) + g(\alpha\bm{x}) \tag{$f+g$の定義} \\
    &= \alpha f(\bm{x}) + \alpha g(\bm{x}) \tag{$f,g$の線型性} \\
    &= \alpha\bigl(f(\bm{x}) + g(\bm{x})\bigr) \nonumber \\
    &= \alpha(f + g)(\bm{x}) \tag{$f+g$の定義}
  \end{align}
  が成り立つ．よって$f + g$は\deref{線型写像}の条件を満たし，線型写像である．

  後半の主張も同様にして証明される．
\end{proof}
\end{leftbar}

\begin{prop}{線型写像の合成}{線型写像の合成}
  $f \colon K^n \to K^m$，$g \colon K^m \to K^l$を線型写像とする．
  このとき，合成写像$g \circ f \colon K^n \to K^l$も線型写像である．
\end{prop}

\begin{leftbar}
\begin{proof}
  任意の$\bm{x}, \bm{y} \in K^n$に対して，
  \begin{align}
    (g \circ f)(\bm{x} + \bm{y}) &= g\bigl(f(\bm{x} + \bm{y})\bigr) \nonumber \\
    &= g\bigl(f(\bm{x}) + f(\bm{y})\bigr) \tag{$f$の線型性} \\
    &= g\bigl(f(\bm{x})\bigr) + g\bigl(f(\bm{y})\bigr) \tag{$g$の線型性} \\
    &= (g \circ f)(\bm{x}) + (g \circ f)(\bm{y}) \nonumber
  \end{align}
  が成り立ち，さらに任意の$\alpha \in K$に対して，
  \begin{align}
    (g \circ f)(\alpha\bm{x}) &= g\bigl(f(\alpha\bm{x})\bigr) \nonumber \\
    &= g\bigl(\alpha f(\bm{x})\bigr) \tag{$f$の線型性} \\
    &= \alpha g\bigl(f(\bm{x})\bigr) \tag{$g$の線型性} \\
    &= \alpha(g \circ f)(\bm{x}) \nonumber
  \end{align}
  が成り立つ．これが証明すべきことであった．
\end{proof}
\end{leftbar}

\phantomsection
\section{線型写像と行列}

\phantomsection
\subsection{線型写像の行列表示}

本節では，線型写像と行列の関係を調べる．
これまでにも，行列と写像にはいくつかの類似性があることを示唆してきた．ここでは，両者の関係をより詳しく調べていこう．
結論からいえば，$K^n$から$K^m$への線型写像は$m \times n$行列とちょうど一対一に対応する．

以下では，第$j$成分が$1$で他の成分がすべて$0$であるベクトル
\[
  \bm{e}_j =
  \begin{pmatrix}
    0 \\ \vdots \\ 1 \\ \vdots \\ 0
  \end{pmatrix}
  \in K^n
  \quad (j = 1, 2, \dots, n)
\]
を\textbf{標準基底ベクトル}\index{ひょうじゅんきていべくとる@標準基底ベクトル}とよぶ\footnote{「基底」という概念そのものは，\partref{線型空間}で正式に定義される（\deref{基底}）．そこで見るように，$(\bm{e}_1, \bm{e}_2, \dots, \bm{e}_n)$は$K^n$の基底のもっとも代表的な例であり，「標準基底ベクトル」という名称もそれに由来する．現段階では，「第$j$成分だけが$1$で他の成分がすべて$0$のベクトル」を表す名前だと思っておけばよい．}．

\begin{prop}{行列の定める線型写像}{行列の定める線型写像}
  $A$を$m \times n$行列とする．
  $\bm{x} \in K^n$に対して$A\bm{x} \in K^m$を対応させる写像を$T_A$とかく．
  このとき，$T_A$は$K^n$から$K^m$への線型写像である．
\end{prop}

\begin{leftbar}
\begin{proof}
  任意の$\bm{x}, \bm{y} \in K^n$と$\alpha \in K$に対して，
  \begin{align}
    T_A(\bm{x} + \bm{y}) &= A(\bm{x} + \bm{y}) \tag{$T_A$の定義} \\
    &= A\bm{x} + A\bm{y} \tag{行列の積の分配律} \\
    &= T_A(\bm{x}) + T_A(\bm{y}) \nonumber
  \end{align}
  が成り立ち，さらに，
  \begin{align}
    T_A(\alpha\bm{x}) &= A(\alpha\bm{x}) \tag{$T_A$の定義} \\
    &= \alpha(A\bm{x}) \tag{行列の積の性質} \\
    &= \alpha T_A(\bm{x}) \nonumber
  \end{align}
  が成り立つ．
  よって$T_A$は線型写像である．
\end{proof}
\end{leftbar}

\prref{行列の定める線型写像}は「行列が与えられれば線型写像が定まる」ことを主張している．
実は，この逆も成り立つ．すなわち，任意の線型写像はある行列から定まる線型写像として表され，しかもそのような行列はただひとつしか存在しない．

\begin{theorem}{線型写像の行列表示}{線型写像の行列表示}
  $T \colon K^n \to K^m$を線型写像とする．
  このとき，ある$m \times n$行列$A$が存在して，任意の$\bm{x} \in K^n$に対して
  \[
    T(\bm{x}) = A\bm{x}
  \]
  が成り立つ．さらに，このような$A$はただひとつしか存在しない．
\end{theorem}

\begin{leftbar}
\begin{proof}
  まず，存在を示す．
  \[
    \bm{a}_j = T(\bm{e}_j) \quad (1 \le j \le n)
  \]
  とおき，これらを列として並べた$m \times n$行列
  \[
    A = (\bm{a}_1, \bm{a}_2, \dots, \bm{a}_n)
  \]
  を考える．$A$の定める線型写像を$T_A \colon K^n \to K^m$とすると，$1 \le j \le n$の範囲で
  \begin{equation}
    T_A(\bm{e}_j) = A\bm{e}_j = \bm{a}_j = T(\bm{e}_j) \label{eq:行列表示の基底での一致}
  \end{equation}
  が成り立つ．ここで，任意の$\bm{x} = \begin{pmatrix} x_1 \\ x_2 \\ \vdots \\ x_n \end{pmatrix} \in K^n$は
  \[
    \bm{x} = \sum_{j=1}^{n} x_j \bm{e}_j
  \]
  とかけるので，$T_A$と$T$の線型性および\eqref{eq:行列表示の基底での一致}より，
  \begin{align*}
    T_A(\bm{x}) &= T_A\left(\sum_{j=1}^{n} x_j \bm{e}_j\right)
    = \sum_{j=1}^{n} x_j T_A(\bm{e}_j)
    = \sum_{j=1}^{n} x_j T(\bm{e}_j)
    = T\left(\sum_{j=1}^{n} x_j \bm{e}_j\right)
    = T(\bm{x})
  \end{align*}
  が成り立つ．よって$T = T_A$であり，$A$の存在が示された．

  次に，一意性を示す．$A$，$B$がともに$m \times n$行列であって$T_A = T_B$を満たすとし，
  \[
    A = (\bm{a}_1, \bm{a}_2, \dots, \bm{a}_n), \quad
    B = (\bm{b}_1, \bm{b}_2, \dots, \bm{b}_n)
  \]
  とかく．このとき，$1 \le j \le n$の範囲で
  \begin{align}
    \bm{a}_j &= A\bm{e}_j \nonumber \\
    &= T_A(\bm{e}_j) \tag{$T_A$の定義} \\
    &= T_B(\bm{e}_j) \tag{$T_A = T_B$} \\
    &= B\bm{e}_j = \bm{b}_j \nonumber
  \end{align}
  が成り立つ．すべての列が一致するので$A = B$であり，一意性が示された．
  これが証明すべきことであった．
\end{proof}
\end{leftbar}

まとめると，行列$A$が与えられたとき，それに対応する写像$T_A$は線型写像であり，逆に，$n$次元のベクトルを$m$次元のベクトルにうつす任意の線型写像は，必ずただひとつの行列によって表される．
この意味で，「$K^n$から$K^m$への線型写像」と「$m \times n$行列」は同一視できる．

\begin{remark}[線型写像は基底ベクトルの行き先で決まる]
  \thref{線型写像の行列表示}の証明の核心は，「線型写像は標準基底ベクトルの行き先だけで完全に決まる」という点にある．
  実際，任意のベクトルは標準基底ベクトルの線型結合（\partref{線型独立性と階数}で定義する）として表せるので，線型性を用いれば，$T(\bm{e}_1), T(\bm{e}_2), \dots, T(\bm{e}_n)$の値から$T(\bm{x})$の値がすべて復元できる．
  この「基底での値がすべてを決める」という考え方は，線型代数の随所で現れる重要な視点である．
\end{remark}

この同一視のもとで，行列の積は写像の合成に対応する．

\begin{prop}{行列の積と写像の合成}{行列の積と写像の合成}
  $A$を$l \times m$行列，$B$を$m \times n$行列とする．このとき，
  \[
    T_A \circ T_B = T_{AB}
  \]
  が成り立つ．
\end{prop}

\begin{leftbar}
\begin{proof}
  任意の$\bm{x} \in K^n$に対して，
  \begin{align}
    (T_A \circ T_B)(\bm{x}) &= T_A\bigl(T_B(\bm{x})\bigr) \nonumber \\
    &= A(B\bm{x}) \tag{$T_A$，$T_B$の定義} \\
    &= (AB)\bm{x} \tag{\prref{行列の乗法結合律}} \\
    &= T_{AB}(\bm{x}) \nonumber
  \end{align}
  が成り立つ\footnote{$\bm{x} \in K^n$を$n \times 1$行列とみなせば，$A(B\bm{x}) = (AB)\bm{x}$は行列の乗法結合律そのものである．}．
  これが証明すべきことであった．
\end{proof}
\end{leftbar}

\prref{行列の積と写像の合成}によって，一見複雑に見えた行列の積の定義（\deref{行列の積}）が，「写像の合成に対応するように定められたもの」として自然に理解できる．
行列の積が交換律を満たさないのは，写像の合成が交換律を満たさないことの反映なのである．

\phantomsection
\subsection{線型変換を表す行列}

\thref{線型写像の行列表示}の証明は，行列$A$の第$j$列が$T(\bm{e}_j)$であることを教えてくれる．
これを用いれば，前節で見た線型変換を表す行列は，標準基底ベクトルの行き先を調べるだけで求められる．

\begin{example}{図形の移動を表す行列}{図形の移動を表す行列}
  $K = \mathbb{R}$とする．
  \exref{対称移動と正射影}の$f_1$（$x$軸に関する対称移動）については，
  \[
    f_1(\bm{e}_1) = \begin{pmatrix} 1 \\ 0 \end{pmatrix}, \quad
    f_1(\bm{e}_2) = \begin{pmatrix} 0 \\ -1 \end{pmatrix}
  \]
  であるから，これらを列として並べて
  \[
    f_1(\bm{x}) = \begin{pmatrix} 1 & 0 \\ 0 & -1 \end{pmatrix} \bm{x}
  \]
  となる．同様にして，
  \[
    f_2(\bm{x}) = \begin{pmatrix} 0 & 1 \\ 1 & 0 \end{pmatrix} \bm{x}
    \quad \text{（直線$y = x$に関する対称移動）}, \qquad
    f_3(\bm{x}) = \begin{pmatrix} 1 & 0 \\ 0 & 0 \end{pmatrix} \bm{x}
    \quad \text{（$x$軸への正射影）}
  \]
  が得られる．
  また，\exref{原点のまわりの回転}の回転については，
  \[
    f(\bm{e}_1) = \begin{pmatrix} \cos\theta \\ \sin\theta \end{pmatrix}, \quad
    f(\bm{e}_2) = \begin{pmatrix} -\sin\theta \\ \cos\theta \end{pmatrix}
  \]
  であるから，
  \[
    R(\theta) = \begin{pmatrix} \cos\theta & -\sin\theta \\ \sin\theta & \cos\theta \end{pmatrix}
  \]
  とおけば$f(\bm{x}) = R(\theta)\bm{x}$である．この$R(\theta)$を\textbf{回転行列}\index{かいてんぎょうれつ@回転行列}という．
\end{example}

\begin{example}{回転行列と加法定理}{回転行列と加法定理}
  角$\varphi$の回転をしてから角$\theta$の回転をすれば，全体としては角$\theta + \varphi$の回転になる．
  これを\prref{行列の積と写像の合成}によって行列の言葉に翻訳すると，
  \[
    R(\theta) R(\varphi) = R(\theta + \varphi)
  \]
  となる．左辺を\deref{行列の積}に従って計算すれば，
  \[
    R(\theta)R(\varphi) =
    \begin{pmatrix}
      \cos\theta\cos\varphi - \sin\theta\sin\varphi & -(\cos\theta\sin\varphi + \sin\theta\cos\varphi) \\
      \sin\theta\cos\varphi + \cos\theta\sin\varphi & \cos\theta\cos\varphi - \sin\theta\sin\varphi
    \end{pmatrix}
  \]
  であり，右辺は
  \[
    R(\theta + \varphi) =
    \begin{pmatrix}
      \cos(\theta + \varphi) & -\sin(\theta + \varphi) \\
      \sin(\theta + \varphi) & \cos(\theta + \varphi)
    \end{pmatrix}
  \]
  である．両辺の成分を比較すれば，三角関数の加法定理
  \[
    \cos(\theta + \varphi) = \cos\theta\cos\varphi - \sin\theta\sin\varphi, \qquad
    \sin(\theta + \varphi) = \sin\theta\cos\varphi + \cos\theta\sin\varphi
  \]
  が得られる．
  行列の積という「一見複雑な定義」が，写像の合成という自然な操作を正確に写し取っていることの，ひとつの現れである．
\end{example}

\begin{example}{回転行列の逆行列}{回転行列の逆行列}
  角$\theta$の回転を打ち消すのは角$-\theta$の回転であるから，
  \[
    R(\theta) R(-\theta) = R(0) = E_2
  \]
  が成り立つ．よって$R(\theta)$は正則であり，
  \[
    R(\theta)^{-1} = R(-\theta) = \begin{pmatrix} \cos\theta & \sin\theta \\ -\sin\theta & \cos\theta \end{pmatrix}
    = R(\theta)^\mathsf{T}
  \]
  である．
  一方，\exref{対称移動と正射影}の正射影$f_3$を表す行列$\begin{pmatrix} 1 & 0 \\ 0 & 0 \end{pmatrix}$は正則でない．
  実際，$\bm{e}_2 \ne \bm{0}$が$\bm{0}$にうつるため，$f_3$は逆をもちえない．
  「つぶれてしまう変換は元に戻せない」という直観が，逆行列の存在・非存在に対応しているのである．
\end{example}

以上の例から，図形の移動を線型変換として捉えることで，その性質を行列の計算によって調べられることがわかる．
このような線型変換と行列の結びつきは，図形の移動だけに現れるものではない．
少し視点を変えると，高校数学で学んだ複素数の乗法も，同じ枠組みで捉えることができる．

\begin{mycolumn}
  $\alpha = p + qi$を定数とし，複素数$z = x + yi$に$\alpha z$を対応させる写像を考えよう．
  \[
    \alpha z = (p + qi)(x + yi) = (px - qy) + (qx + py)i
  \]
  であるから，$z = x + yi$を$\begin{pmatrix} x \\ y \end{pmatrix} \in \mathbb{R}^2$と同一視すれば，この対応は
  \[
    \begin{pmatrix} x \\ y \end{pmatrix} \longmapsto
    \begin{pmatrix} p & -q \\ q & p \end{pmatrix}
    \begin{pmatrix} x \\ y \end{pmatrix}
  \]
  という$\mathbb{R}^2$の線型変換にほかならない．
  とくに$\alpha = \cos\theta + i\sin\theta$のとき，この行列は回転行列$R(\theta)$と一致する．
  「複素数を極形式でかけ算すると偏角が足される」という高校で学ぶ事実と，
  \exref{回転行列と加法定理}の$R(\theta)R(\varphi) = R(\theta + \varphi)$とが，同じ内容を述べていることがわかる．
\end{mycolumn}

\phantomsection
\subsection{数列と線型変換}

数列に関する例も，行列の言葉で見直してみよう．

\begin{example}{階差行列と部分和行列}{階差行列と部分和行列}
  \exref{部分和をとる写像}の$\Sigma \colon K^n \to K^n$を表す行列を求めよう．
  $\Sigma \bm{e}_j$は「第$j$項だけが$1$で他が$0$の数列」の部分和であるから，
  第$j$成分以降がすべて$1$，それより前が$0$のベクトルである．
  これらを列として並べると，$n = 4$のとき
  \[
    L =
    \begin{pmatrix}
      1 & 0 & 0 & 0 \\
      1 & 1 & 0 & 0 \\
      1 & 1 & 1 & 0 \\
      1 & 1 & 1 & 1
    \end{pmatrix}
  \]
  となる（対角成分より下がすべて$1$の行列である）．

  一方，\exref{階差数列をとる写像}の$\Delta$を少し変形し，「第$1$項はそのまま残し，第$2$項以降は階差をとる」写像
  \[
    \widetilde{\Delta} \bm{a} =
    \begin{pmatrix} a_1 \\ a_2 - a_1 \\ a_3 - a_2 \\ a_4 - a_3 \end{pmatrix}
  \]
  を考えると，これは$K^4$の線型変換であり，その行列は
  \[
    D =
    \begin{pmatrix}
      1 & 0 & 0 & 0 \\
      -1 & 1 & 0 & 0 \\
      0 & -1 & 1 & 0 \\
      0 & 0 & -1 & 1
    \end{pmatrix}
  \]
  である．実際に積を計算すれば
  \[
    DL = LD = E_4
  \]
  となり，$D = L^{-1}$であることがわかる．
  これは，「部分和をとってから階差をとると元の数列に戻る」という\exref{部分和をとる写像}の観察を，
  逆行列の言葉で述べ直したものである．
\end{example}

\begin{example}{漸化式が定める線型変換}{漸化式が定める線型変換}
  $p, q \in K$を定数とし，数列$\{a_n\}$が三項間漸化式
  \[
    a_{n+2} = p\,a_{n+1} + q\,a_n
  \]
  を満たすとする．
  連続する$2$項を並べたベクトル$\begin{pmatrix} a_n \\ a_{n+1} \end{pmatrix} \in K^2$に着目すると，
  漸化式は
  \[
    \begin{pmatrix} a_{n+1} \\ a_{n+2} \end{pmatrix}
    = \begin{pmatrix} 0 & 1 \\ q & p \end{pmatrix}
      \begin{pmatrix} a_{n} \\ a_{n+1} \end{pmatrix}
  \]
  と表される．
  すなわち，「$1$項だけ先へ進める」という操作が$K^2$の線型変換にほかならない．

  この行列を$A$とおくと，\prref{行列の積と写像の合成}をくり返し用いて
  \[
    \begin{pmatrix} a_{n+1} \\ a_{n+2} \end{pmatrix}
    = A^{n} \begin{pmatrix} a_{1} \\ a_{2} \end{pmatrix}
  \]
  が得られる．
  たとえばフィボナッチ数列（$F_1 = F_2 = 1$，$F_{n+2} = F_{n+1} + F_n$）であれば，
  \[
    A = \begin{pmatrix} 0 & 1 \\ 1 & 1 \end{pmatrix}, \qquad
    \begin{pmatrix} F_{n+1} \\ F_{n+2} \end{pmatrix}
    = A^{n} \begin{pmatrix} 1 \\ 1 \end{pmatrix}
  \]
  である．漸化式を$n$回追いかける代わりに，行列の$n$乗を計算すればよいことになる．
\end{example}

\begin{mycolumn}
\exref{漸化式が定める線型変換}によって，数列の一般項を求める問題は「行列$A$の$n$乗を求める問題」に翻訳された．
では，$A^n$はどう計算すればよいのだろうか．

鍵となるのは，$A\bm{x} = \lambda \bm{x}$を満たすベクトル$\bm{x} \ne \bm{0}$である．
このとき$A^n \bm{x} = \lambda^n \bm{x}$となり，べき乗が一気に見通しよくなる．
この$\lambda$を$A$の固有値とよび，のちの章で詳しく扱う．
高校で三項間漸化式を解くときに現れる$2$次方程式$t^2 = pt + q$（特性方程式）は，
実は$A$の固有値を求める方程式にほかならない．

また，漸化式$a_{n+2} = p a_{n+1} + q a_n$を満たす数列全体は，項ごとの和とスカラー倍について閉じている．
高校で学ぶ「$2$つの解の線型結合もまた解である」という事実は，この集合が線型空間をなすことを意味する．
さらに，初項と第$2$項で解が決まるためその次元は$2$であり，一般解が$2$つの解の線型結合で表されることも説明できる．
これらの概念は\partref{線型空間}で整備される．
\end{mycolumn}

\phantomsection
\section{線型写像の像と核}

\phantomsection
\subsection{像と核}

線型写像$f \colon K^n \to K^m$が与えられると，「$f$の行き先として実際に現れるベクトル全体」と「$f$によって$\bm{0}$にうつされるベクトル全体」という二つの集合が自然に定まる．
本節では，この二つを\textbf{像}・\textbf{核}として定式化する．
これらは，\partref{線型空間}の次元公式（\thref{線型写像の次元公式}）をはじめ，以後の議論の随所で用いられる基本概念である．

その準備として，まず一般の写像に関する用語を整えておこう．

\begin{definition}{単射・全射・全単射}{単射と全射}
  $f \colon V \to W$を写像とする．
  \begin{enumerate}[label=(\roman*),leftmargin=3.2em]
    \item 任意の$x, y \in V$に対して「$f(x) = f(y)$ならば$x = y$」が成り立つとき，$f$は\textbf{単射}\index{たんしゃ@単射}であるという．
    \item 任意の$y \in W$に対して，$f(x) = y$となる$x \in V$が存在するとき，$f$は\textbf{全射}\index{ぜんしゃ@全射}であるという．
    \item $f$が単射かつ全射であるとき，$f$は\textbf{全単射}\index{ぜんたんしゃ@全単射}であるという．
  \end{enumerate}
\end{definition}

単射とは「異なる元は必ず異なる元にうつる」こと，全射とは「終域のどの元も行き先として現れる」ことである\footnote{全単射は，\partref{行列式}で置換を定義する際にも用いられる．}．

\begin{definition}{像と核}{像と核}
  $f \colon K^n \to K^m$を線型写像とする．
  \[
    \Im f = \{ f(\bm{x}) \mid \bm{x} \in K^n \} \subset K^m
  \]
  を$f$の\textbf{像}\index{ぞう@像}\index{Im@$\Im$}（image）という．
  また，
  \[
    \Ker f = \{ \bm{x} \in K^n \mid f(\bm{x}) = \bm{0} \} \subset K^n
  \]
  を$f$の\textbf{核}\index{かく@核}\index{Ker@$\Ker$}（kernel）という．
\end{definition}

像が終域$K^m$の部分集合であるのに対し，核は定義域$K^n$の部分集合である．
両者の「住んでいる場所」の違いに注意してほしい\footnote{像は，線型とは限らない一般の写像$f \colon V \to W$に対しても同じ式で定義される．一方，核の定義には終域に$\bm{0}$という特別な元があることが必要であり，こちらは線型写像に固有の概念である．}．

\prref{線型写像と零ベクトル}より$f(\bm{0}) = \bm{0}$であるから，つねに$\bm{0} \in \Ker f$かつ$\bm{0} \in \Im f$であり，とくに像も核も空集合ではない．
また，定義から直ちに，$f$が全射であることと$\Im f = K^m$であることは同値である．

\begin{example}{像と核}{像と核の例}
  \begin{enumerate}[label=(\arabic*),leftmargin=3.2em]
    \item 恒等変換$\mathrm{id} \colon K^n \to K^n$（\exref{恒等変換と零写像}）については，
      \[
        \Im \mathrm{id} = K^n, \qquad \Ker \mathrm{id} = \{\bm{0}\}
      \]
      である．一方，零写像$f \colon K^n \to K^m$については，
      \[
        \Im f = \{\bm{0}\}, \qquad \Ker f = K^n
      \]
      である．
    \item $x$軸への正射影$f_3 \colon \mathbb{R}^2 \to \mathbb{R}^2$（\exref{対称移動と正射影}）については，
      \[
        \Im f_3 = \left\{ \begin{pmatrix} x \\ 0 \end{pmatrix} ~\middle|~ x \in \mathbb{R} \right\}, \qquad
        \Ker f_3 = \left\{ \begin{pmatrix} 0 \\ y \end{pmatrix} ~\middle|~ y \in \mathbb{R} \right\}
      \]
      である．すなわち，像は$x$軸，核は$y$軸である．
      平面全体が$x$軸へとつぶされるとき，「届く範囲」が像であり，「原点へつぶれる方向」が核なのである．
    \item 多項式の微分$D \colon \mathbb{R}^3 \to \mathbb{R}^2$（\exref{多項式の微分と定積分}）については，
      \[
        \Ker D = \left\{ \begin{pmatrix} a_0 \\ 0 \\ 0 \end{pmatrix} ~\middle|~ a_0 \in \mathbb{R} \right\}
      \]
      である．これは「（$2$次以下の）多項式のうち，微分して$0$になるのは定数にかぎる」という事実の言い換えである．
      また，$1$次以下の任意の多項式$b_0 + b_1 x$は，たとえば$b_0 x + \frac{b_1}{2} x^2$の微分として得られる（原始関数の存在）から，$\Im D = \mathbb{R}^2$である．
  \end{enumerate}
\end{example}

\phantomsection
\subsection{像・核と連立一次方程式}

連立一次方程式のことばを使えば，行列の定める線型写像$T_A$の像と核は，すでによく知っている対象であることがわかる．
以下，右辺の定数項がすべて$0$である連立一次方程式$A\bm{x} = \bm{0}$を\textbf{斉次連立一次方程式}\index{せいじれんりついちじほうていしき@斉次連立一次方程式}（同次連立一次方程式）という．

\begin{prop}{連立一次方程式と像・核}{連立一次方程式と像と核}
  $A$を$m \times n$行列とし，$T_A \colon K^n \to K^m$を$A$の定める線型写像（\prref{行列の定める線型写像}）とする．このとき，次が成り立つ．
  \begin{enumerate}[label=(\roman*),leftmargin=3.2em]
    \item $\Ker T_A$は，斉次連立一次方程式$A\bm{x} = \bm{0}$の解全体の集合である．
    \item $\bm{b} \in K^m$に対して，$\bm{b} \in \Im T_A$であることと，連立一次方程式$A\bm{x} = \bm{b}$が解をもつことは同値である．
  \end{enumerate}
\end{prop}

\begin{leftbar}
\begin{proof}
  (i)について，$\bm{x} \in \Ker T_A$であることは，$T_A$の定義より$A\bm{x} = \bm{0}$と同値であり，これは$\bm{x}$が$A\bm{x} = \bm{0}$の解であることにほかならない．

  (ii)について，$\bm{b} \in \Im T_A$であることは，ある$\bm{x} \in K^n$が存在して$A\bm{x} = \bm{b}$となること，すなわち$A\bm{x} = \bm{b}$が解をもつことと同値である．
  いずれも定義の言い換えである．これが証明すべきことであった．
\end{proof}
\end{leftbar}

このように，核は「斉次連立一次方程式の解全体」を，像は「右辺として実現可能なベクトル全体」を，写像のことばで言い直したものである．

核のもっとも重要な役割のひとつは，線型写像の単射性の判定である．

\begin{prop}{単射性と核}{単射性と核}
  線型写像$f \colon K^n \to K^m$に対して，次の二つは同値である．
  \begin{enumerate}[label=(\roman*),leftmargin=3.2em]
    \item $f$は単射である．
    \item $\Ker f = \{\bm{0}\}$である．
  \end{enumerate}
\end{prop}

\begin{leftbar}
\begin{proof}
  (i)$\limp$(ii)：$\bm{0} \in \Ker f$はつねに成り立つので，逆の包含を示せばよい．
  $\bm{x} \in \Ker f$とすると，\prref{線型写像と零ベクトル}より
  \[
    f(\bm{x}) = \bm{0} = f(\bm{0})
  \]
  である．$f$は単射であるから$\bm{x} = \bm{0}$となり，$\Ker f = \{\bm{0}\}$である．

  (ii)$\limp$(i)：$f(\bm{x}) = f(\bm{y})$とする．$f$の線型性より
  \[
    f(\bm{x} - \bm{y}) = f\bigl(\bm{x} + (-1)\bm{y}\bigr) = f(\bm{x}) + (-1)f(\bm{y}) = f(\bm{x}) - f(\bm{y}) = \bm{0}
  \]
  であるから，$\bm{x} - \bm{y} \in \Ker f = \{\bm{0}\}$である．
  よって$\bm{x} - \bm{y} = \bm{0}$，すなわち$\bm{x} = \bm{y}$となり，$f$は単射である．
  これが証明すべきことであった．
\end{proof}
\end{leftbar}

一般の写像の単射性を確かめるには，すべての値の組$f(\bm{x})$，$f(\bm{y})$を比較しなければならない．
これに対して線型写像では，\prref{単射性と核}により「$\bm{0}$にうつるベクトルが$\bm{0}$だけかどうか」を調べれば十分である．
これが核という概念の威力である．

\exref{回転行列の逆行列}では，正射影$f_3$について「$\bm{e}_2 \ne \bm{0}$が$\bm{0}$にうつるため，$f_3$は逆をもちえない」と述べた．
これはまさに$\bm{e}_2 \in \Ker f_3$，すなわち$\Ker f_3 \ne \{\bm{0}\}$という主張であり，\prref{単射性と核}より$f_3$は単射でない．
「つぶれてしまう変換は元に戻せない」という直観は，核のことばで正確に表現されるのである．

\begin{mycolumn}
  像と核は，この後の章で中心的な役割を果たす．
  \partref{線型空間}では，$\Im f$と$\Ker f$が単なる部分集合ではなく，それぞれ$K^m$，$K^n$の\textbf{線型部分空間}であることを確かめる（\prref{像と核の線型部分空間性}）．
  さらに「次元」を導入すると，両者のあいだには
  \[
    \dim \Ker f + \dim \Im f = n
  \]
  という関係が成り立つ（線型写像の次元公式，\thref{線型写像の次元公式}）．
  定義域のもつ$n$次元分の情報のうち，核の分だけが$\bm{0}$につぶれ，残りが像として生き残る，という直観を正確に述べたのがこの公式である．
\end{mycolumn}

ここまでで，線型写像と行列が一対一に対応すること，および線型写像に付随する像と核について見てきた．
次章では，ベクトルの組の「無駄のなさ」を測る概念として線型独立性を導入し，行列の階数を定義しよう．

\phantomsection
\part{線型独立性と階数}

\begin{abstract}
この章では，ベクトルの組の「無駄のなさ」を測る概念を整理する．
まず\textbf{線型結合}と線型関係を定義し，線型写像が線型結合を保つことを確かめる．
次に，\textbf{線型独立}と\textbf{線型従属}を定義し，その言い換えと，以後くり返し用いる「線型独立な組の拡大」「取替補題」を扱う．
最後に，行列の\textbf{階数}を線型独立な列の最大個数として定義し，行階数と列階数の一致，および行基本変形による階数の計算法を示す．
\end{abstract}

\phantomsection
\section{線型結合と線型独立性}

\phantomsection
\subsection{線型結合と線型関係}

本節では，複数のベクトルからつくられる基本的な演算である\textbf{線型結合}\index{せんけいけつごう@線型結合}を定義する\footnote{ここでは$K^m$のベクトルに対して定義するが，\partref{線型空間}で導入する一般の線型空間の元に対しても，まったく同様に定義される．}．

\begin{definition}{線型結合・線型関係}{線型結合}
  $n$個のベクトル$\bm{x}_1, \bm{x}_2, \dots, \bm{x}_n \in K^m$を考える．
  \begin{equation}
    \lambda_1 \bm{x}_1 + \lambda_2 \bm{x}_2 + \cdots + \lambda_n \bm{x}_n
    \quad (\lambda_i \in K) \label{eq:線型結合の形}
  \end{equation}
  の形のベクトルを，$\bm{x}_1, \bm{x}_2, \dots, \bm{x}_n$の\textbf{線型結合}という．

  また，\eqref{eq:線型結合の形}に関連して，等式
  \[
    \lambda_1 \bm{x}_1 + \lambda_2 \bm{x}_2 + \cdots + \lambda_n \bm{x}_n = \bm{0}
  \]
  を$\bm{x}_1, \bm{x}_2, \dots, \bm{x}_n$の\textbf{線型関係}\index{せんけいかんけい@線型関係}という．
\end{definition}

任意のベクトル$\bm{x}_1, \bm{x}_2, \dots, \bm{x}_n$に対して，係数をすべて$0$とした
\[
  0\bm{x}_1 + 0\bm{x}_2 + \cdots + 0\bm{x}_n = \bm{0}
\]
はつねに成り立つ線型関係である．これを\textbf{自明な線型関係}\index{じめいなせんけいかんけい@自明な線型関係}とよぶ．

\partref{線型写像}で導入した線型写像は，線型結合と相性がよい．
すなわち，線型写像は線型結合を保つ：

\begin{prop}{線型写像と線型結合}{線型写像と線型結合}
  $f \colon K^n \to K^m$を線型写像とする．
  このとき，任意の$\bm{x}_1, \bm{x}_2, \dots, \bm{x}_k \in K^n$と任意の$\lambda_1, \lambda_2, \dots, \lambda_k \in K$に対して，
  \[
    f(\lambda_1 \bm{x}_1 + \lambda_2 \bm{x}_2 + \cdots + \lambda_k \bm{x}_k)
    = \lambda_1 f(\bm{x}_1) + \lambda_2 f(\bm{x}_2) + \cdots + \lambda_k f(\bm{x}_k)
  \]
  が成り立つ．
\end{prop}

\begin{leftbar}
\begin{proof}
  $k$に関する帰納法で示す．
  $k = 1$のときは，$f$の線型性の第$2$の条件そのものである．

  $k - 1$のとき主張が成り立つと仮定する．このとき，
  \begin{align}
    f\left(\sum_{i=1}^{k} \lambda_i \bm{x}_i\right)
    &= f\left(\sum_{i=1}^{k-1} \lambda_i \bm{x}_i + \lambda_k \bm{x}_k\right) \nonumber \\
    &= f\left(\sum_{i=1}^{k-1} \lambda_i \bm{x}_i\right) + f(\lambda_k \bm{x}_k) \tag{$f$の線型性} \\
    &= \sum_{i=1}^{k-1} \lambda_i f(\bm{x}_i) + \lambda_k f(\bm{x}_k) \tag{帰納法の仮定と$f$の線型性} \\
    &= \sum_{i=1}^{k} \lambda_i f(\bm{x}_i) \nonumber
  \end{align}
  となり，$k$のときも主張が成り立つ．これが証明すべきことであった．
\end{proof}
\end{leftbar}

\enlargethispage{4.5mm}

\begin{remark}[$A\bm{x}$は列ベクトルの線型結合]
  行列とベクトルの積は，線型結合のことばで見直すことができる．
  $m \times n$行列$A$の列ベクトルを$\bm{a}_1, \bm{a}_2, \dots, \bm{a}_n \in K^m$とすると，任意の$\bm{x} = \begin{pmatrix} x_1 \\ x_2 \\ \vdots \\ x_n \end{pmatrix} \in K^n$に対して，
  \[
    A\bm{x} = x_1 \bm{a}_1 + x_2 \bm{a}_2 + \cdots + x_n \bm{a}_n
  \]
  が成り立つ．実際，$\bm{x} = \sum_{j=1}^{n} x_j \bm{e}_j$と\prref{線型写像と線型結合}より，
  \[
    A\bm{x} = T_A\left(\sum_{j=1}^{n} x_j \bm{e}_j\right)
    = \sum_{j=1}^{n} x_j T_A(\bm{e}_j)
    = \sum_{j=1}^{n} x_j \bm{a}_j
  \]
  となる．
  すなわち，$A\bm{x}$は$A$の列ベクトルの線型結合であり，その係数は$\bm{x}$の成分にほかならない．
  この見方によれば，連立一次方程式$A\bm{x} = \bm{b}$が解を持つことは，$\bm{b}$が$A$の列ベクトルの線型結合として表せることと同値である．
  一方，\prref{連立一次方程式と像と核}(ii)より，$A\bm{x} = \bm{b}$が解を持つことは$\bm{b} \in \Im T_A$とも同値であった．
  あわせると，像$\Im T_A$とは「$A$の列ベクトルの線型結合全体の集合」にほかならないことがわかる．
\end{remark}


\phantomsection
\subsection{線型独立と線型従属}

線型結合の定義をもとに，ベクトルの組の「無駄のなさ」を測る概念として\textbf{線型独立性}を定義する．
線型
\newpage

\noindent 
関係のうち，自明なものしか存在しないとき，ベクトルの組は\textbf{線型独立}\index{せんけいどくりつ@線型独立}であるという．

\begin{definition}{線型独立}{線型独立}
  $\bm{x}_1, \bm{x}_2, \dots, \bm{x}_n \in K^m$が\textbf{線型独立}であるとは，
  \[
    \lambda_1 \bm{x}_1 + \lambda_2 \bm{x}_2 + \cdots + \lambda_n \bm{x}_n = \bm{0}
  \]
  ならば，
  \[
    \lambda_1=\lambda_2=\cdots=\lambda_n = 0
  \]
  であることである．
  線型独立は\textbf{一次独立}\index{いちじどくりつ@一次独立}ともいう．
\end{definition}

\begin{prop}{線型独立の言い換え}{線型独立の言い換え}
  $\bm{x}_1,\bm{x}_2,\dots,\bm{x}_n\in K^m$が線型独立であることは，
  少なくとも一つが$0$でない任意の
  $\lambda_1,\lambda_2,\dots,\lambda_n\in K$に対して，
  \[
    \lambda_1\bm{x}_1+\lambda_2\bm{x}_2+\cdots
    +\lambda_n\bm{x}_n\ne\bm{0}
  \]
  となることと同値である．
\end{prop}

\begin{leftbar}
\begin{proof}
  \deref{線型独立}の条件の対偶を取れば明らかである．
\end{proof}
\end{leftbar}

線型独立でないベクトルの組は，\textbf{線型従属}\index{せんけいじゅうぞく@線型従属}であるという．

\begin{definition}{線型従属}{線型従属}
  $\bm{x}_1, \bm{x}_2, \dots, \bm{x}_n \in K^m$が線型独立でないとき，すなわち，少なくとも一つが$0$でない$\lambda_1, \lambda_2, \dots, \lambda_n \in K$が存在して
  \[
    \lambda_1 \bm{x}_1 + \lambda_2 \bm{x}_2 + \cdots + \lambda_n \bm{x}_n = \bm{0}
  \]
  が成り立つとき，$\bm{x}_1, \bm{x}_2, \dots, \bm{x}_n$は\textbf{線型従属}であるという．
  線型従属は\textbf{一次従属}\index{いちじじゅうぞく@一次従属}ともいう．
\end{definition}

\begin{example}{線型独立と線型従属}{線型独立と線型従属}
  $K^2$のベクトル
  \[
    \bm{e}_1 = \begin{pmatrix} 1 \\ 0 \end{pmatrix}, \quad
    \bm{e}_2 = \begin{pmatrix} 0 \\ 1 \end{pmatrix}
  \]
  は線型独立である．実際，
  \[
    \lambda_1 \bm{e}_1 + \lambda_2 \bm{e}_2 = \begin{pmatrix} \lambda_1 \\ \lambda_2 \end{pmatrix} = \bm{0}
  \]
  ならば$\lambda_1 = \lambda_2 = 0$となるからである．

  一方，$K^2$のベクトル
  \[
    \bm{x}_1 = \begin{pmatrix} 1 \\ 2 \end{pmatrix}, \quad
    \bm{x}_2 = \begin{pmatrix} 2 \\ 4 \end{pmatrix}
  \]
  は線型従属である．実際，$(\lambda_1, \lambda_2) = (2, -1) \ne (0, 0)$に対して
  \[
    2\bm{x}_1 + (-1)\bm{x}_2 = \bm{0}
  \]
  という非自明な線型関係が成り立つからである．
\end{example}

最後に，線型従属性の重要な言い換えを示そう．
これは「線型従属なベクトルの組では，どれかひとつのベクトルが残りのベクトルの線型結合として表せる」という主張であり，のちの章でも繰り返し用いられる．

\begin{prop}{線型従属の言い換え}{線型従属の言い換え}
  $n \ge 2$とする．
  $\bm{x}_1, \bm{x}_2, \dots, \bm{x}_n \in K^m$が線型従属であることは，ある$i$が存在して，$\bm{x}_i$が残りの$n - 1$個のベクトルの線型結合として表せることと同値である．
\end{prop}

\begin{leftbar}
\begin{proof}
  まず，$\bm{x}_1, \bm{x}_2, \dots, \bm{x}_n$が線型従属であるとする．
  このとき，少なくとも一つが$0$でないスカラー$\lambda_1, \lambda_2, \dots, \lambda_n \in K$が存在して
  \[
    \lambda_1 \bm{x}_1 + \lambda_2 \bm{x}_2 + \cdots + \lambda_n \bm{x}_n = \bm{0}
  \]
  が成り立つ．$\lambda_i \ne 0$となる$i$をとると，
  \[
    \bm{x}_i = \sum_{\substack{1 \le j \le n \\ j \ne i}} \left(-\frac{\lambda_j}{\lambda_i}\right) \bm{x}_j
  \]
  となり，$\bm{x}_i$は残りのベクトルの線型結合として表せる．

  逆に，ある$i$に対して
  \[
    \bm{x}_i = \sum_{\substack{1 \le j \le n \\ j \ne i}} \lambda_j \bm{x}_j
    \quad (\lambda_j \in K)
  \]
  と表せるとする．このとき，
  \[
    \lambda_1 \bm{x}_1 + \cdots + \lambda_{i-1} \bm{x}_{i-1} + (-1)\bm{x}_i + \lambda_{i+1} \bm{x}_{i+1} + \cdots + \lambda_n \bm{x}_n = \bm{0}
  \]
  が成り立つ．$\bm{x}_i$の係数は$-1 \ne 0$であるから，これは非自明な線型関係である．
  よって$\bm{x}_1, \bm{x}_2, \dots, \bm{x}_n$は線型従属である．

  以上の考察により，命題が示された．
\end{proof}
\end{leftbar}

すなわち，線型従属なベクトルの組には，残りのベクトルの線型結合として「すでに表せてしまう」ベクトルが存在するのである．

なお，与えられたベクトルの線型結合として表されるベクトル全体がなす集合は，\partref{線型空間}において\textbf{線型包}として詳しく扱う．
そこでは，本章で定義した線型結合の概念が，数ベクトルに限らず一般の線型空間においてもそのまま通用することを見る．

\phantomsection
\subsection{線型独立性に関する基本補題}

線型独立性について，以後くり返し用いる二つの補題を用意しておこう．
いずれも線型結合と線型独立の定義のみから導かれる．

なお，これらの補題は，\partref{線型空間}において一般の線型空間に対してあらためて定式化され，証明される（\lemref{一般の取替補題}）．
ここでは，\partref{行列式}までの議論で必要となる数ベクトル空間$K^m$の場合を先に示しておく．

\begin{lemma}{線型独立な組の拡大}{線型独立な組の拡大}
  $\bm{v}_1, \bm{v}_2, \dots, \bm{v}_k \in K^m$を線型独立とし，$\bm{w} \in K^m$は
  $\bm{v}_1, \bm{v}_2, \dots, \bm{v}_k$の線型結合として表せないとする．
  このとき，$\bm{v}_1, \bm{v}_2, \dots, \bm{v}_k, \bm{w}$は線型独立である．
\end{lemma}

\begin{leftbar}
\begin{proof}
  \[
    \lambda_1 \bm{v}_1 + \cdots + \lambda_k \bm{v}_k + \mu \bm{w} = \bm{0}
  \]
  とする．もし$\mu \ne 0$ならば，
  \[
    \bm{w} = \sum_{i=1}^{k} \left(-\frac{\lambda_i}{\mu}\right) \bm{v}_i
  \]
  となり，$\bm{w}$が線型結合として表せることになって仮定に反する．
  よって$\mu = 0$であり，このとき$\sum_{i=1}^{k} \lambda_i \bm{v}_i = \bm{0}$となるから，
  線型独立性より$\lambda_1 = \cdots = \lambda_k = 0$である．
  これが証明すべきことであった．
\end{proof}
\end{leftbar}

次の補題は，「線型独立なベクトルは，それを生成するベクトルの個数を超えて取れない」ことを述べている．

\begin{lemma}{取替補題}{取替補題}
  $\bm{v}_1, \bm{v}_2, \dots, \bm{v}_n \in K^m$とする．
  $\bm{w}_1, \bm{w}_2, \dots, \bm{w}_l \in K^m$がいずれも$\bm{v}_1, \bm{v}_2, \dots, \bm{v}_n$の線型結合として表され，
  かつ線型独立であるならば，
  \[
    l \le n
  \]
  が成り立つ．
\end{lemma}

\begin{leftbar}
\begin{proof}
  $k = 0, 1, \dots, l$に対して，次の主張$(\ast_k)$を$k$についての帰納法で示す．

  \begin{quote}
    $(\ast_k)$：$k \le n$であり，かつ$\bm{v}_1, \dots, \bm{v}_n$の番号を適当に付け替えることで，
    $\bm{v}_1, \dots, \bm{v}_n$の線型結合として表せるベクトルはすべて
    \[
      \bm{w}_1, \dots, \bm{w}_k, \bm{v}_{k+1}, \dots, \bm{v}_n
    \]
    の線型結合としても表せる，とできる．
  \end{quote}

  $k = 0$のときは主張の内容が空であり，成り立つ．

  $k < l$として$(\ast_k)$を仮定する．
  $\bm{w}_{k+1}$は$\bm{v}_1, \dots, \bm{v}_n$の線型結合であるから，帰納法の仮定より
  \begin{equation}
    \bm{w}_{k+1}
    = \sum_{i=1}^{k} a_i \bm{w}_i + \sum_{j=k+1}^{n} b_j \bm{v}_j
    \label{eq:取替補題の展開}
  \end{equation}
  と表せる．

  ここで，$b_{k+1} = \cdots = b_n = 0$であると仮定してみる
  （$k = n$のときは第$2$の和が空なので，自動的にこの場合にあたる）．
  このとき\eqref{eq:取替補題の展開}は
  \[
    a_1 \bm{w}_1 + \cdots + a_k \bm{w}_k + (-1)\bm{w}_{k+1} = \bm{0}
  \]
  と書き直せる．$\bm{w}_{k+1}$の係数は$-1 \ne 0$であるから，これは非自明な線型関係であり，
  $\bm{w}_1, \dots, \bm{w}_l$の線型独立性に矛盾する．

  よって$k < n$，すなわち$k + 1 \le n$であり，さらに$b_j \ne 0$となる$j$が存在する．
  残っている$\bm{v}_{k+1}, \dots, \bm{v}_n$の番号を付け替えて$b_{k+1} \ne 0$としてよい．
  このとき\eqref{eq:取替補題の展開}を$\bm{v}_{k+1}$について解けば，
  \[
    \bm{v}_{k+1}
    = \frac{1}{b_{k+1}}\left( \bm{w}_{k+1} - \sum_{i=1}^{k} a_i \bm{w}_i - \sum_{j=k+2}^{n} b_j \bm{v}_j \right)
  \]
  となり，$\bm{v}_{k+1}$は$\bm{w}_1, \dots, \bm{w}_{k+1}, \bm{v}_{k+2}, \dots, \bm{v}_n$の線型結合として表せる．
  したがって，$\bm{w}_1, \dots, \bm{w}_k, \bm{v}_{k+1}, \dots, \bm{v}_n$の線型結合として表せるベクトルは，
  すべて$\bm{w}_1, \dots, \bm{w}_{k+1}, \bm{v}_{k+2}, \dots, \bm{v}_n$の線型結合としても表せる．
  よって$(\ast_{k+1})$が成り立つ．

  以上より$(\ast_l)$が成り立ち，とくに$l \le n$である．
  これが証明すべきことであった．
\end{proof}
\end{leftbar}

\begin{cor}{$K^m$における線型独立なベクトルの個数}{Km における線型独立なベクトルの個数}
  $K^m$の線型独立なベクトルは，高々$m$個である．
  すなわち，$m + 1$個以上のベクトルは必ず線型従属である．
\end{cor}

\begin{leftbar}
\begin{proof}
  任意の$\bm{x} = (x_i) \in K^m$は
  \[
    \bm{x} = \sum_{i=1}^{m} x_i \bm{e}_i
  \]
  と標準基底ベクトルの線型結合で表せる．
  よって\lemref{取替補題}を$\bm{v}_i = \bm{e}_i$として適用すればよい．
  これが証明すべきことであった．
\end{proof}
\end{leftbar}

\begin{cor}{$K^n$の$n$個の線型独立なベクトル}{Kn の n 個の線型独立なベクトル}
  $\bm{a}_1, \bm{a}_2, \dots, \bm{a}_n \in K^n$が線型独立ならば，
  任意の$\bm{b} \in K^n$は$\bm{a}_1, \bm{a}_2, \dots, \bm{a}_n$の線型結合として表せる．
\end{cor}

\begin{leftbar}
\begin{proof}
  ある$\bm{b} \in K^n$が線型結合として表せないと仮定すると，\lemref{線型独立な組の拡大}より
  $\bm{a}_1, \dots, \bm{a}_n, \bm{b}$は線型独立である．
  これは$K^n$の$n + 1$個の線型独立なベクトルであり，\coref{Km における線型独立なベクトルの個数}に矛盾する．
  これが証明すべきことであった．
\end{proof}
\end{leftbar}

\phantomsection
\section{階数}

\phantomsection
\subsection{階数}

行列に対して，その列ベクトルがどれだけ「独立な情報」をもっているかを表す量を定める．

\begin{definition}{階数}{階数}
  $A$を$m \times n$行列とし，その列ベクトルを$\bm{a}_1, \bm{a}_2, \dots, \bm{a}_n \in K^m$とする．
  $\bm{a}_1, \bm{a}_2, \dots, \bm{a}_n$のうちから選んで線型独立となるベクトルの最大個数を，
  $A$の\textbf{階数}\index{かいすう@階数}\index{rank@$\rank$}といい，$\rank A$と表す．

  ただし，$0$個のベクトルの組は線型独立であると約束する．
  したがって$A = O$のとき$\rank A = 0$である．
\end{definition}

\begin{remark}[階数の定義について]
  「線型独立な列の最大個数」という定義は，一見すると回りくどく感じられるかもしれない．
  実は，\partref{線型空間}で「次元」を定義したあとであれば，階数は「$A$の列ベクトルの線型結合として表せるベクトル全体（列空間とよばれる）の次元」と言い直すことができる．
  この列空間は，像$\Im T_A$にほかならない（\prref{線型写像と線型結合}のあとのRemarkを参照）．
  すなわち階数は，\partref{線型空間}のことばでは「線型写像$T_A$の像の次元」$\dim \Im T_A$である（\coref{階数と像の次元}）．
  しかし本章の段階では次元がまだ定義されていないため，「次元」という語を使わずに同じ数を取り出す方法として，この定義を採用しているのである．
\end{remark}

\begin{example}{階数}{階数の例}
  \[
    A = \begin{pmatrix} 1 & 2 & 0 \\ 2 & 4 & 0 \end{pmatrix}
  \]
  の列ベクトルは
  \[
    \bm{a}_1 = \begin{pmatrix} 1 \\ 2 \end{pmatrix}, \quad
    \bm{a}_2 = \begin{pmatrix} 2 \\ 4 \end{pmatrix}, \quad
    \bm{a}_3 = \begin{pmatrix} 0 \\ 0 \end{pmatrix}
  \]
  である．$\bm{a}_2 = 2\bm{a}_1$，$\bm{a}_3 = \bm{0}$であるから，
  どの$2$本を選んでも線型従属であり，一方$\bm{a}_1$の$1$本は線型独立である．
  よって$\rank A = 1$である．
\end{example}

階数の定義は列ベクトルによるものであったが，実は行ベクトルで定義しても同じ値が得られる．

\begin{theorem}{行階数と列階数の一致}{行階数と列階数の一致}
  $A$を$m \times n$行列とする．
  $A$の行ベクトルのうち線型独立なものの最大個数は，$\rank A$に等しい．
  したがって，
  \[
    \rank A^\mathsf{T} = \rank A
  \]
  が成り立つ．
\end{theorem}

\begin{leftbar}
\begin{proof}
  $\rank A = r$とし，$A$の行ベクトルのうち線型独立なものの最大個数を$r'$とする．

  まず$r' \le r$を示す．
  $A$の列ベクトル$\bm{a}_1, \dots, \bm{a}_n$のうち，線型独立な$r$本を選び，
  それらを$\bm{c}_1, \dots, \bm{c}_r$とする．
  各$j$について，$\bm{c}_1, \dots, \bm{c}_r, \bm{a}_j$は$r + 1$本であり，$r$の最大性から線型従属である．
  よって\lemref{線型独立な組の拡大}の対偶より，$\bm{a}_j$は$\bm{c}_1, \dots, \bm{c}_r$の線型結合として表せる：
  \[
    \bm{a}_j = \sum_{i=1}^{r} c_{ij} \bm{c}_i \quad (1 \le j \le n).
  \]

  そこで，$\bm{c}_1, \dots, \bm{c}_r$を列として並べた$m \times r$行列を$C$，
  $r \times n$行列$(c_{ij})$を$R$とおくと，上の式は$CR$の第$j$列が$\bm{a}_j$であることを意味するので，
  \[
    A = CR
  \]
  が成り立つ．
  行列の積の定義（\deref{行列の積}）より，$A$の第$k$行は
  \[
    (A \text{の第}k\text{行}) = \sum_{i=1}^{r} (C \text{の}(k,i)\text{成分}) \cdot (R \text{の第}i\text{行})
  \]
  とかける．すなわち，$A$の各行は$R$の$r$本の行の線型結合である．
  したがって，$A$の線型独立な行を$r'$本とれば，\lemref{取替補題}より$r' \le r$である．

  次に$r \le r'$を示す．
  前半の議論は任意の行列に対して通用する．すなわち，任意の行列$M$に対して
  \[
    (M\text{の線型独立な行の最大個数})
    \le
    (M\text{の線型独立な列の最大個数})
  \]
  が成り立つ．これを$M = A^\mathsf{T}$に適用する．
  $A^\mathsf{T}$の行は$A$の列であるから，左辺は$A$の線型独立な列の最大個数，すなわち$r$に等しい．
  また，$A^\mathsf{T}$の列は$A$の行であるから，右辺は$A$の線型独立な行の最大個数，すなわち$r'$に等しい．
  したがって$r \le r'$である．

  以上より$r = r'$であり，$A$の線型独立な行の最大個数は$\rank A$に等しい．
  さらに，$A^\mathsf{T}$の列は$A$の行であるから，
  \[
    \rank A^\mathsf{T} = r' = r = \rank A
  \]
  も成り立つ．
  これが証明すべきことであった．
\end{proof}
\end{leftbar}

\begin{cor}{階数の上限}{階数の上限}
  $m \times n$行列$A$に対して，
  \[
    \rank A \le \min\{m, n\}
  \]
  が成り立つ．
\end{cor}

\begin{leftbar}
\begin{proof}
  $A$の列は$n$本しかないので$\rank A \le n$である．
  また$A$の行は$m$本しかないので，\thref{行階数と列階数の一致}より$\rank A \le m$である．
  これが証明すべきことであった．
\end{proof}
\end{leftbar}

\begin{remark}[行階数と列階数の一致の意味]
  \thref{行階数と列階数の一致}は，一見すると当たり前ではない主張である．
  行の本数と列の本数は一般に異なるうえ，行の線型独立性と列の線型独立性は別々の条件だからである．
  それにもかかわらず両者が一致するという事実は，階数が行列の「かたち」ではなく本質的な情報を捉えていることを示している．

  なお，証明中に現れた分解$A = CR$は，「$m \times n$行列を，階数$r$の情報だけを使って$m \times r$行列と$r \times n$行列の積に分解する」ものであり，
  それ自体が応用上重要な意味をもつ．
\end{remark}

\phantomsection
\subsection{階数の計算}

階数の定義はそのままでは計算に向かない．
実際には，行基本変形（\deref{行基本変形}）によって階数を求めることができる．

その根拠となるのが，次の観察である．

\begin{prop}{列の線型関係と斉次連立一次方程式}{列の線型関係と斉次連立方程式}
  $A$を$m \times n$行列とし，その列ベクトルを$\bm{a}_1, \dots, \bm{a}_n$とする．
  このとき，$\bm{x} = (x_j) \in K^n$に対して次は同値である：
  \begin{enumerate}[label=(\roman*),leftmargin=3.2em]
    \item $A\bm{x} = \bm{0}$である．
    \item $x_1 \bm{a}_1 + x_2 \bm{a}_2 + \cdots + x_n \bm{a}_n = \bm{0}$である．
  \end{enumerate}
  すなわち，$A$の列ベクトルの線型関係と，斉次連立一次方程式$A\bm{x} = \bm{0}$の解は同じものである．
\end{prop}

\begin{leftbar}
\begin{proof}
  \prref{線型写像と線型結合}のあとのRemarkで見たとおり，
  \[
    A\bm{x} = x_1 \bm{a}_1 + x_2 \bm{a}_2 + \cdots + x_n \bm{a}_n
  \]
  が成り立つ．これが証明すべきことであった．
\end{proof}
\end{leftbar}

\begin{prop}{行基本変形と解集合}{行基本変形と解集合}
  $A$に行基本変形を施して得られる行列を$A'$とすると，
  \[
    \{\bm{x} \in K^n \mid A\bm{x} = \bm{0}\}
    = \{\bm{x} \in K^n \mid A'\bm{x} = \bm{0}\}
  \]
  が成り立つ．
\end{prop}

\begin{leftbar}
\begin{proof}
  $A\bm{x} = \bm{0}$の第$i$式を$f_i(\bm{x}) = 0$とかく．
  行基本変形の各操作で得られる新しい式は，いずれも$f_1, \dots, f_m$の線型結合である．
  よって$A\bm{x} = \bm{0}$を満たす$\bm{x}$は$A'\bm{x} = \bm{0}$も満たす．

  また，\deref{行基本変形}の$3$種類の操作はいずれも逆操作をもつ
  （第$i$行を$c$倍する操作の逆は$1/c$倍する操作，行の入れ替えの逆はもう一度入れ替える操作，
  第$i$行に第$j$行の$c$倍を加える操作の逆は$(-c)$倍を加える操作である）．
  同じ議論を逆操作に適用すれば，逆向きの包含も従う．
  これが証明すべきことであった．
\end{proof}
\end{leftbar}

\begin{theorem}{行基本変形と階数}{行基本変形と階数}
  行基本変形によって階数は変化しない．
\end{theorem}

\begin{leftbar}
\begin{proof}
  $A$に行基本変形を施して$A'$が得られたとする．
  列の番号$j_1 < j_2 < \cdots < j_k$を任意にとり，$A$の第$j_1, \dots, j_k$列だけを取り出した行列を$B$，
  $A'$の同じ列を取り出した行列を$B'$とする．
  $A$から$A'$への行基本変形は，$B$から$B'$への行基本変形をそのまま引き起こすから，
  \prref{行基本変形と解集合}より$B\bm{y} = \bm{0}$と$B'\bm{y} = \bm{0}$の解集合は一致する．

  \prref{列の線型関係と斉次連立方程式}より，
  $A$の第$j_1, \dots, j_k$列が線型独立であることは$B\bm{y} = \bm{0}$が自明な解しかもたないことと同値であり，
  $A'$についても同様である．
  よって，どの列の組についても，$A$で線型独立であることと$A'$で線型独立であることは一致する．
  線型独立な列の最大個数も一致するので，$\rank A = \rank A'$である．
  これが証明すべきことであった．
\end{proof}
\end{leftbar}

\begin{remark}[行基本変形による階数の計算]
  \thref{行基本変形と階数}により，階数は次のようにして求められる．
  行基本変形をくり返して$A$をできるだけ簡単な形に変形すると，
  最終的に，$0$でない各行の先頭の非零成分が右下に向かって階段状に並び，
  $0$のみからなる行が下側に集まった形にできる（さらに必要なら，各先頭成分を$1$にできる）．
  この形の行列では，$0$でない行はすべて線型独立であり，残りは$0$のみからなる行であるから，
  線型独立な行の最大個数は$0$でない行の本数に等しい．
  \thref{行階数と列階数の一致}より，これがそのまま階数である．

  \exref{階数の例}の行列$A$であれば，
  \[
    \begin{pmatrix} 1 & 2 & 0 \\ 2 & 4 & 0 \end{pmatrix}
    \xrightarrow{\text{第$2$行に第$1$行の$(-2)$倍を加える}}
    \begin{pmatrix} 1 & 2 & 0 \\ 0 & 0 & 0 \end{pmatrix}
  \]
  となり，$0$でない行は$1$本であるから$\rank A = 1$と読み取れる．
\end{remark}

\begin{mycolumn}
  階数は，のちに線型空間の次元を定義したあとで，
  「$A$の列ベクトルが張る線型部分空間（列空間）の次元」として捉え直される．
  \deref{階数}の「線型独立な列の最大個数」という定義は，その次元の値をあらかじめ直接に取り出したものにあたる．
\end{mycolumn}

ここまでの3つの章で，線型代数の基礎となる概念についての議論を終えた．
次章ではこれらの概念を踏まえて，行列式を定義しよう．

\input{contents/part4_determinant.tex}
\phantomsection 
\part{線型空間}

\begin{abstract}
  この章では，線型空間を導入し，ベクトルを抽象的な立場から捉える．
  これにより，数ベクトルだけでなく，行列や関数なども同じ枠組みで扱えるようになる．
  さらに，線型部分空間，基底，次元などの基本概念を導入し，線型空間の構造について考察する．
\end{abstract}

\phantomsection
\section{線型空間}

\phantomsection
\subsection{線型空間の定義}


線型代数学では，ある条件を満たす集合の元を「ベクトル」として扱う．
そこで，その枠組みとなる\textbf{線型空間}\index{せんけいくうかん@線型空間}を定義する．

\partref{行列と連立一次方程式}から\partref{行列式}までは$K=\mathbb{R}$または$K=\mathbb{C}$としていたが，
本章では$K$を任意の\textbf{体}\index{たい@体}（field）とする%
\footnote{体とは，大まかにいえば，加法・減法・乗法，および$0$でない元による除法が自由に行える代数系である．
$\mathbb{R}$や$\mathbb{C}$はその代表例である．}．

次の定義では，体$K$の加法・乗法と
$V$の加法・スカラー倍を区別するため，それぞれ異なる記号を用いる．

\begin{definition}{線型空間}{線型空間}
  $K$を体とし，その加法と乗法をそれぞれ
  \[
    \Kadd\colon K\times K\to K,
    \qquad
    \Kmul\colon K\times K\to K
  \]
  とする．また，加法の零元を$0_K$，乗法の単位元を$1_K$とする．

  集合$V$に，加法とスカラー倍
  \[
    \Vadd\colon V\times V\to V,
    \qquad
    \scalmul\colon K\times V\to V
  \]
  が定められているとする．
  このとき，次の条件を満たす$V$を体$K$上の\textbf{線型空間}という．

  \begin{enumerate}[label=(\roman*),leftmargin=3.2em]
    \item \textbf{加法に関して可換群}

    $V$は，加法$\Vadd$に関して可換群をなす．

    \item \textbf{スカラー倍に関する性質}

    任意の$\bm{v},\bm{w}\in V$および$\alpha,\beta\in K$に対して，
    \[
      1_K\scalmul\bm{v}=\bm{v},
      \qquad
      (\alpha\Kadd\beta)\scalmul\bm{v}
      =
      (\alpha\scalmul\bm{v})\Vadd(\beta\scalmul\bm{v}),
    \]
    \[
      \alpha\scalmul(\bm{v}\Vadd\bm{w})
      =
      (\alpha\scalmul\bm{v})\Vadd(\alpha\scalmul\bm{w}),
      \qquad
      (\alpha\Kmul\beta)\scalmul\bm{v}
      =
      \alpha\scalmul(\beta\scalmul\bm{v})
    \]
    が成り立つ．
  \end{enumerate}

  $V$の元を\textbf{ベクトル}という．
\end{definition}

この定義では，$4$つの演算を区別するために，
一時的に$\Kadd$，$\Kmul$，$\Vadd$，$\scalmul$という記号を用いた．
しかし，どの演算であるかは文脈から判断できるため，以下では
\[
  \alpha\Kadd\beta,\quad
  \bm{v}\Vadd\bm{w},\quad
  \alpha\Kmul\beta,\quad
  \alpha\scalmul\bm{v}
\]
を，それぞれ
\[
  \alpha+\beta,\quad
  \bm{v}+\bm{w},\quad
  \alpha\beta,\quad
  \alpha\bm{v}
\]
と表す．
たとえば
\[
  (\alpha+\beta)\bm{v}
  =
  \alpha\bm{v}+\beta\bm{v}
\]
では，左辺の$+$は$K$の加法，右辺の$+$は$V$の加法である．

\begin{remark}[加法に関する可換群の条件]
線型空間の公理のうち，
「加法に関して可換群である」という条件には，

\begin{itemize}
  \item 結合法則
  \item 交換法則
  \item 零ベクトルの存在
  \item 加法逆元の存在
\end{itemize}

が含まれている．
具体的には，任意の$\bm{v},\bm{w},\bm{u}\in V$に対して，
\[
  (\bm{v}+\bm{w})+\bm{u}
  =
  \bm{v}+(\bm{w}+\bm{u}),
\]
\[
  \bm{v}+\bm{w}
  =
  \bm{w}+\bm{v},
\]
\[
  \bm{v}+\bm{0}_V
  =
  \bm{v},
\]
\[
  \bm{v}+(-\bm{v})
  =
  \bm{0}_V
\]
が成り立つ．
ここではすでに通常の表記に戻しているため，
$V$の加法には通常の$+$を用いている．
\end{remark}

\begin{example}{線型空間の例}{線型空間の例}
次の集合はいずれも，通常の加法とスカラー倍に関して線型空間となる．

\begin{itemize}
\item $\mathbb{R}$上の線型空間$\mathbb{R}^n$
\item $\mathbb{C}$上の線型空間$\mathbb{C}^n$
\item $K$上の線型空間$M_{m,n}(K)$（$K$の元を成分とする$m\times n$行列全体，\deref{行列全体の集合}）
\item $\mathbb{R}$上の線型空間$\mathbb{R}[x]$（実数係数多項式全体）
\end{itemize}

これらは一見異なる対象であるが，
線型空間という共通の枠組みで扱うことができる\footnote{同じ集合でも，どの体の上の線型空間とみるかによって構造は異なる．たとえば$\mathbb{C}^n$は，スカラーを実数に制限すれば$\mathbb{R}$上の線型空間ともみなせる．線型空間を考えるときに係数体を明示するのはこのためである．}．
\end{example}

線型空間の公理からは，数ベクトルの場合と同様の計算規則が導かれる．

\begin{prop}{線型空間の基本性質}{線型空間の基本性質}
  $V$を体$K$上の線型空間とする．
  任意の$\bm{v} \in V$および任意の$\alpha \in K$に対して，次が成り立つ：
  \begin{enumerate}[label=(\roman*),leftmargin=3.2em]
    \item $0_K \bm{v} = \bm{0}_V$
    \item $\alpha \bm{0}_V = \bm{0}_V$
    \item $(-1_K)\bm{v} = -\bm{v}$
  \end{enumerate}
\end{prop}

\begin{leftbar}
\begin{proof}
  (i)について，スカラー倍に関する分配律より，
  \[
    0_K \bm{v} = (0_K + 0_K)\bm{v} = 0_K \bm{v} + 0_K \bm{v}
  \]
  である．両辺に$0_K \bm{v}$の加法逆元$-(0_K \bm{v})$を加えると，
  \[
    \bm{0}_V = 0_K \bm{v}
  \]
  を得る．

  (ii)について，同様に，
  \[
    \alpha \bm{0}_V = \alpha(\bm{0}_V + \bm{0}_V) = \alpha \bm{0}_V + \alpha \bm{0}_V
  \]
  であるから，両辺に$-(\alpha \bm{0}_V)$を加えて$\bm{0}_V = \alpha \bm{0}_V$を得る．

  (iii)について，(i)より，
  \[
    \bm{v} + (-1_K)\bm{v} = 1_K \bm{v} + (-1_K)\bm{v} = \bigl(1_K + (-1_K)\bigr)\bm{v} = 0_K \bm{v} = \bm{0}_V
  \]
  である．加法逆元の一意性より，$(-1_K)\bm{v} = -\bm{v}$である．
  以上の考察により，命題が示された．
\end{proof}
\end{leftbar}

\subsection{線型部分空間}

線型空間の部分集合のうち，
それ自身も線型空間となるものを線型部分空間という．

\begin{definition}{線型部分空間}{線型部分空間}
  体$K$上の線型空間$V$の部分集合$W$が，
  次の条件を満たすとき，
  $W$を$V$の\textbf{線型部分空間}という．
\begin{enumerate}[label=(\roman*),leftmargin=3.2em]
    \item \textbf{空でない}
    \[
      W\neq\varnothing
    \]
    である．

    \item \textbf{加法に関する性質}

    任意の$\bm{v},\bm{w}\in W$に対して，
    \[
      \bm{v}+\bm{w}\in W
    \]
    である．

    \item \textbf{スカラー倍に関する性質}

    任意の$\bm{v}\in W$，
    任意の$\alpha\in K$に対して，
    \[
      \alpha \bm{v}\in W
    \]
    である．
  \end{enumerate}
\end{definition}

\begin{remark}[空でないことと零ベクトル]
定義の最初の条件は
\[
\bm{0}_V\in W
\]
としてもよい．

実際，$W$が空でなければ$\bm{w} \in W$がとれ，$W$がスカラー倍について閉じていることと\prref{線型空間の基本性質}より，
\[
\bm{0}_V = 0_K \bm{w} \in W
\]
となる．
\end{remark}

線型部分空間の例をいくつか見よう．

\begin{example}{線型部分空間の例}{線型部分空間の例}
  \begin{enumerate}[label=(\arabic*),leftmargin=3.2em]
    \item 任意の線型空間$V$に対して，$\{\bm{0}_V\}$と$V$自身は，いずれも$V$の線型部分空間である．
    \item $K = \mathbb{R}$とする．原点を通る直線
      \[
        W = \left\{ \begin{pmatrix} x \\ y \end{pmatrix} \in \mathbb{R}^2 ~\middle|~ 2x - y = 0 \right\}
      \]
      は$\mathbb{R}^2$の線型部分空間である．
      実際，$\bm{0} \in W$であり，$2x_1 - y_1 = 0$かつ$2x_2 - y_2 = 0$ならば$2(x_1 + x_2) - (y_1 + y_2) = 0$，また任意の$\alpha \in \mathbb{R}$に対して$2(\alpha x_1) - \alpha y_1 = \alpha(2x_1 - y_1) = 0$となるからである．
      一方，原点を通らない直線（たとえば$2x - y = 1$で定まる直線）は$\bm{0}$を含まないので，線型部分空間ではない．
    \item 次数が$n$以下の実数係数多項式全体（零多項式を含む）は，$\mathbb{R}[x]$の線型部分空間である．
      次数$n$以下の多項式どうしの和や，そのスカラー倍によって，次数が$n$を超えることはないからである．
  \end{enumerate}
\end{example}

さらに，\partref{線型写像}で導入した線型写像の像と核（\deref{像と核}）は，線型部分空間の重要な例を与える\footnote{\partref{線型写像}では$K$を$\mathbb{R}$または$\mathbb{C}$としていたが，線型写像とその像・核の定義（\deref{線型写像}，\deref{像と核}），および\prref{線型写像と零ベクトル}や\prref{線型写像と線型結合}などの基本性質は，一般の体$K$に対してもそのまま通用する．}．

\begin{prop}{像と核の線型部分空間性}{像と核の線型部分空間性}
  $f \colon K^n \to K^m$を線型写像とする．このとき，次が成り立つ．
  \begin{enumerate}[label=(\roman*),leftmargin=3.2em]
    \item $\Ker f$は$K^n$の線型部分空間である．
    \item $\Im f$は$K^m$の線型部分空間である．
  \end{enumerate}
\end{prop}

\begin{leftbar}
\begin{proof}
  (i)について示す．
  \prref{線型写像と零ベクトル}より$f(\bm{0}) = \bm{0}$であるから，$\bm{0} \in \Ker f$であり，$\Ker f$は空でない．
  $\bm{x}, \bm{y} \in \Ker f$とすると，$f$の線型性より
  \[
    f(\bm{x} + \bm{y}) = f(\bm{x}) + f(\bm{y}) = \bm{0} + \bm{0} = \bm{0}
  \]
  であるから，$\bm{x} + \bm{y} \in \Ker f$である．
  また，任意の$\alpha \in K$に対して
  \[
    f(\alpha\bm{x}) = \alpha f(\bm{x}) = \alpha\bm{0} = \bm{0}
  \]
  であるから，$\alpha\bm{x} \in \Ker f$である．
  よって$\Ker f$は$K^n$の線型部分空間である．

  (ii)について示す．
  $\bm{0} = f(\bm{0}) \in \Im f$であるから，$\Im f$は空でない．
  $\bm{v}, \bm{w} \in \Im f$とすると，$\bm{v} = f(\bm{x})$，$\bm{w} = f(\bm{y})$となる$\bm{x}, \bm{y} \in K^n$がとれる．$f$の線型性より
  \[
    \bm{v} + \bm{w} = f(\bm{x}) + f(\bm{y}) = f(\bm{x} + \bm{y}) \in \Im f
  \]
  であり，また任意の$\alpha \in K$に対して
  \[
    \alpha\bm{v} = \alpha f(\bm{x}) = f(\alpha\bm{x}) \in \Im f
  \]
  である．よって$\Im f$は$K^m$の線型部分空間である．
  これが証明すべきことであった．
\end{proof}
\end{leftbar}

とくに，斉次連立一次方程式$A\bm{x} = \bm{0}$の解全体は$\Ker T_A$であったから（\prref{連立一次方程式と像と核}），$K^n$の線型部分空間である．
これを$A\bm{x} = \bm{0}$の\textbf{解空間}\index{かいくうかん@解空間}という．
「$2$つの解の和はふたたび解であり，解のスカラー倍もまた解である」という重ね合わせの性質こそが，部分空間性の意味するところである．

線型部分空間どうしから新たな線型部分空間を作る方法として，
共通部分と和が重要である．

\begin{definition}{共通部分・和}{共通部分・和}
  $V$を線型空間，
  $W_1,W_2$を$V$の線型部分空間とする．

  このとき，
  \[
  W_1\cap W_2
  \]
  を$W_1$と$W_2$の\textbf{共通部分}という．

  また，
  \[
  W_1+W_2
  =
  \{
  \bm{w}_1+\bm{w}_2
  \mid
  \bm{w}_1\in W_1,\
  \bm{w}_2\in W_2
  \}
  \]
  を$W_1$と$W_2$の\textbf{和}という．
\end{definition}

\begin{prop}{共通部分と和の線型部分空間性}{共通部分と和の線型部分空間性}
  \deref{共通部分・和}で定義した，
  線型部分空間の共通部分および和は，
  ともに線型部分空間である．
\end{prop}

\begin{leftbar}
\begin{proof}
  まず，共通部分$W_1\cap W_2$について示す．
  $W_1,W_2$はいずれも線型部分空間であるから，
  \[
    \bm{0}_V\in W_1,
    \quad
    \bm{0}_V\in W_2
  \]
  である．したがって，
  \[
    \bm{0}_V\in W_1\cap W_2
  \]
  であり，$W_1\cap W_2$は空でない．

  次に，$\bm{x},\bm{y}\in W_1\cap W_2$とする．
  このとき，
  \[
    (\bm{x},\bm{y}\in W_1)
    \quad \land \quad
    (\bm{x},\bm{y}\in W_2)
  \]
  である．$W_1,W_2$は加法について閉じているから，
  \[
    (\bm{x}+\bm{y}\in W_1)
    \quad \land \quad
    (\bm{x}+\bm{y}\in W_2)
  \]
  となる．ゆえに，
  \[
    \bm{x}+\bm{y}\in W_1\cap W_2
  \]
  である．

  また，任意の$\alpha\in K$に対して，
  $W_1,W_2$はスカラー倍について閉じているから，
  \[
    (\alpha\bm{x}\in W_1)
    \quad \land \quad
    (\alpha\bm{x}\in W_2)
  \]
  となる．したがって，
  \[
    \alpha\bm{x}\in W_1\cap W_2
  \]
  である．以上より，$W_1\cap W_2$は$V$の線型部分空間である．

  続いて，和$W_1+W_2$について示す．
  \[
    \bm{0}_V=\bm{0}_V+\bm{0}_V
  \]
  であり，$\bm{0}_V\in W_1$かつ$\bm{0}_V\in W_2$であるから，
  \[
    \bm{0}_V\in W_1+W_2
  \]
  である．したがって，$W_1+W_2$は空でない．

  $\bm{x},\bm{y}\in W_1+W_2$とする．
  このとき，ある$\bm{x}_1,\bm{y}_1\in W_1$および
  $\bm{x}_2,\bm{y}_2\in W_2$が存在して，
  \[
    \bm{x}=\bm{x}_1+\bm{x}_2,
    \quad
    \bm{y}=\bm{y}_1+\bm{y}_2
  \]
  と表せる．したがって，
  \begin{align*}
    \bm{x}+\bm{y}
    &=(\bm{x}_1+\bm{x}_2)+(\bm{y}_1+\bm{y}_2)\\
    &=(\bm{x}_1+\bm{y}_1)+(\bm{x}_2+\bm{y}_2)
  \end{align*}
  である．ここで，
  \[
    \bm{x}_1+\bm{y}_1\in W_1,
    \quad
    \bm{x}_2+\bm{y}_2\in W_2
  \]
  であるから，
  \[
    \bm{x}+\bm{y}\in W_1+W_2
  \]
  となる．

  さらに，任意の$\alpha\in K$に対して，
  \[
    \alpha\bm{x}
    =\alpha(\bm{x}_1+\bm{x}_2)
    =\alpha\bm{x}_1+\alpha\bm{x}_2
  \]
  である．$\alpha\bm{x}_1\in W_1$かつ
  $\alpha\bm{x}_2\in W_2$であるから，
  \[
    \alpha\bm{x}\in W_1+W_2
  \]
  である．以上より，$W_1+W_2$も$V$の線型部分空間である．
\end{proof}
\end{leftbar}

\phantomsection
\subsection{線型結合と線型包}

\partref{線型独立性と階数}では，数ベクトルに対して線型結合を定義した（\deref{線型結合}）．
その定義は，一般の線型空間の元に対しても，まったく同様に通用する．

\begin{definition}{線型結合（一般の線型空間）}{一般の線型結合}
  $V$を体$K$上の線型空間とする．
  $n$個のベクトル$\bm{v}_1, \bm{v}_2, \dots, \bm{v}_n \in V$に対して，
  \[
    \lambda_1 \bm{v}_1 + \lambda_2 \bm{v}_2 + \cdots + \lambda_n \bm{v}_n
    \quad (\lambda_i \in K)
  \]
  の形のベクトルを，$\bm{v}_1, \bm{v}_2, \dots, \bm{v}_n$の\textbf{線型結合}という．

  線型関係，線型独立，線型従属についても，\partref{線型独立性と階数}と同じ式によって定義される．
\end{definition}

ベクトルの集合が与えられたとき，その元の線型結合として表されるベクトルをすべて集めた集合を考えよう．

\begin{definition}{線型包}{線型包}
  $V$を体$K$上の線型空間とし，$S$を$V$の部分集合とする．

  $S \neq \varnothing$のとき，$S$の有限個の元の線型結合全体の集合
  \[
    \Span(S)
    = \{ \lambda_1 \bm{v}_1 + \lambda_2 \bm{v}_2 + \cdots + \lambda_n \bm{v}_n \mid n \ge 1,\ \bm{v}_1, \dots, \bm{v}_n \in S,\ \lambda_1, \dots, \lambda_n \in K \}
  \]
  を$S$の\textbf{線型包}\index{せんけいほう@線型包}\index{Span@$\Span$}（\textbf{張る空間}\index{はるくうかん@張る空間}，生成される部分空間）という．
  また，$\Span(\varnothing) = \{\bm{0}_V\}$と定める．

  有限個のベクトル$\bm{v}_1, \bm{v}_2, \dots, \bm{v}_n \in V$に対しては，
  \[
    \Span(\bm{v}_1, \bm{v}_2, \dots, \bm{v}_n) = \Span(\{\bm{v}_1, \bm{v}_2, \dots, \bm{v}_n\})
  \]
  ともかく．
\end{definition}

\begin{example}{線型包の例}{線型包の例}
  $K^3$の標準基底ベクトル$\bm{e}_1$，$\bm{e}_2$に対して，
  \[
    \Span(\bm{e}_1, \bm{e}_2)
    = \left\{ \lambda_1 \begin{pmatrix} 1 \\ 0 \\ 0 \end{pmatrix} + \lambda_2 \begin{pmatrix} 0 \\ 1 \\ 0 \end{pmatrix} ~\middle|~ \lambda_1, \lambda_2 \in K \right\}
    = \left\{ \begin{pmatrix} \lambda_1 \\ \lambda_2 \\ 0 \end{pmatrix} ~\middle|~ \lambda_1, \lambda_2 \in K \right\}
  \]
  である．$K = \mathbb{R}$のとき，これは$xyz$空間における$xy$平面にほかならない．

  また，線型空間$V$の$1$個のベクトル$\bm{v} \neq \bm{0}_V$に対して，
  \[
    \Span(\bm{v}) = \{ \lambda \bm{v} \mid \lambda \in K \}
  \]
  であり，たとえば$V = \mathbb{R}^3$のとき，これは原点を通り$\bm{v}$の方向に伸びる直線である．

  さらに，$V$を実数係数多項式全体のなす線型空間とすると，
  \[
    \Span(1, x, x^2)
    = \{ a + bx + cx^2 \mid a, b, c \in \mathbb{R} \}
  \]
  であり，これは次数が$2$以下の多項式全体（零多項式を含む）にほかならない．
  このように，線型包の概念は数ベクトルに限らず適用できる．
\end{example}

線型包は，つねに線型部分空間となる．

\begin{prop}{線型包の線型部分空間性}{線型包の線型部分空間性}
  $V$を体$K$上の線型空間，$S$を$V$の部分集合とする．
  このとき，$\Span(S)$は$V$の線型部分空間である．
\end{prop}

\begin{leftbar}
\begin{proof}
  まず，$S = \varnothing$のときを考える．
  このとき$\Span(\varnothing) = \{\bm{0}_V\}$であり，
  \[
    \bm{0}_V + \bm{0}_V = \bm{0}_V
  \]
  および，任意の$\alpha \in K$に対して\prref{線型空間の基本性質}より
  \[
    \alpha \bm{0}_V = \bm{0}_V
  \]
  が成り立つ．よって$\{\bm{0}_V\}$は空でなく，加法とスカラー倍について閉じているから，$V$の線型部分空間である．

  次に，$S \neq \varnothing$とする．
  $\bm{v} \in S$をとると，
  \[
    \bm{v} = 1_K \bm{v} \in \Span(S)
  \]
  であるから，$\Span(S)$は空でない．

  $\bm{x}, \bm{y} \in \Span(S)$とすると，ある$\bm{v}_1, \dots, \bm{v}_m \in S$，$\bm{w}_1, \dots, \bm{w}_l \in S$および$\lambda_i, \mu_j \in K$を用いて
  \[
    \bm{x} = \sum_{i=1}^{m} \lambda_i \bm{v}_i, \quad
    \bm{y} = \sum_{j=1}^{l} \mu_j \bm{w}_j
  \]
  と表せる．このとき，
  \[
    \bm{x} + \bm{y}
    = \lambda_1 \bm{v}_1 + \cdots + \lambda_m \bm{v}_m + \mu_1 \bm{w}_1 + \cdots + \mu_l \bm{w}_l
  \]
  も$S$の有限個の元の線型結合であるから，$\bm{x} + \bm{y} \in \Span(S)$である．

  さらに，任意の$\alpha \in K$に対して，スカラー倍に関する分配律より，
  \[
    \alpha\bm{x} = \sum_{i=1}^{m} (\alpha\lambda_i) \bm{v}_i
  \]
  も$S$の有限個の元の線型結合であるから，$\alpha\bm{x} \in \Span(S)$である．
  以上より，$\Span(S)$は$V$の線型部分空間である．
\end{proof}
\end{leftbar}

さらに，線型包は「$S$を含む最小の線型部分空間」として特徴づけられる．

\begin{prop}{線型包の最小性}{線型包の最小性}
  $V$を体$K$上の線型空間，$S$を$V$の部分集合とする．このとき，次が成り立つ：
  \begin{enumerate}[label=(\roman*),leftmargin=3.2em]
    \item $S \subset \Span(S)$である．
    \item $W$が$S \subset W$を満たす$V$の線型部分空間ならば，$\Span(S) \subset W$である．
  \end{enumerate}
  すなわち，$\Span(S)$は，$S$を含む最小の線型部分空間である．
\end{prop}

\begin{leftbar}
\begin{proof}
  (i)について，任意の$\bm{v} \in S$は
  \[
    \bm{v} = 1_K \bm{v}
  \]
  と表せるから，$\bm{v} \in \Span(S)$である．

  (ii)について，まず$S = \varnothing$のときを考える．
  任意の線型部分空間$W$は必ず零ベクトルを元にもつ（\deref{線型部分空間}のあとのRemarkを参照）から，
  \[
    \Span(\varnothing) = \{\bm{0}_V\} \subset W
  \]
  である．

  次に，$S \neq \varnothing$とし，$\Span(S)$の任意の元
  \[
    \bm{x} = \lambda_1 \bm{v}_1 + \lambda_2 \bm{v}_2 + \cdots + \lambda_n \bm{v}_n
    \quad (\bm{v}_i \in S,\ \lambda_i \in K)
  \]
  をとる．$S \subset W$より$\bm{v}_i \in W$であり，$W$はスカラー倍について閉じているから，
  \[
    \lambda_i \bm{v}_i \in W \quad (1 \le i \le n)
  \]
  である．さらに，$W$は加法について閉じているから，これらの和である$\bm{x}$も$W$に属する．
  よって$\Span(S) \subset W$である．
  以上の考察により，命題が示された．
\end{proof}
\end{leftbar}

\begin{remark}[線型包と共通部分]
\prref{線型包の最小性}より，$\Span(S)$は「$S$を含む線型部分空間すべての共通部分」とも一致する．

実際，\prref{線型包の線型部分空間性}と\prref{線型包の最小性}(i)より，$\Span(S)$自身が「$S$を含む線型部分空間」のひとつであるから，共通部分は$\Span(S)$に含まれる．
逆に，\prref{線型包の最小性}(ii)より，$\Span(S)$は$S$を含むどの線型部分空間にも含まれるから，共通部分にも含まれる．

なお，\prref{共通部分と和の線型部分空間性}では二つの線型部分空間の共通部分を扱ったが，同様の議論により，任意個の線型部分空間の共通部分も線型部分空間となる．
\end{remark}

和と線型包のあいだには，次の関係が成り立つ．

\begin{prop}{和と線型包}{和と線型包}
  $V$を線型空間，$W_1, W_2$を$V$の線型部分空間とする．このとき，
  \[
    W_1 + W_2 = \Span(W_1 \cup W_2)
  \]
  が成り立つ．
\end{prop}

\begin{leftbar}
\begin{proof}
  まず，$W_1 + W_2 \subset \Span(W_1 \cup W_2)$を示す．
  任意の$\bm{w}_1 \in W_1$，$\bm{w}_2 \in W_2$に対して，
  \[
    \bm{w}_1 + \bm{w}_2 = 1_K \bm{w}_1 + 1_K \bm{w}_2
  \]
  は$W_1 \cup W_2$の元の線型結合であるから，$\bm{w}_1 + \bm{w}_2 \in \Span(W_1 \cup W_2)$である．

  逆に，$\Span(W_1 \cup W_2) \subset W_1 + W_2$を示す．
  \prref{共通部分と和の線型部分空間性}より$W_1 + W_2$は$V$の線型部分空間であり，
  \[
    \bm{w}_1 = \bm{w}_1 + \bm{0}_V \in W_1 + W_2,
    \quad
    \bm{w}_2 = \bm{0}_V + \bm{w}_2 \in W_1 + W_2
  \]
  より$W_1 \cup W_2 \subset W_1 + W_2$である．
  よって，\prref{線型包の最小性}より$\Span(W_1 \cup W_2) \subset W_1 + W_2$である．
  以上の考察により，命題が示された．
\end{proof}
\end{leftbar}

\begin{remark}[合併は線型部分空間とは限らない]
一般に，線型部分空間どうしの合併$W_1 \cup W_2$は線型部分空間になるとは限らない．
たとえば，$\mathbb{R}^2$において$x$軸と$y$軸はいずれも線型部分空間であるが，
\[
  \begin{pmatrix} 1 \\ 0 \end{pmatrix}
  +
  \begin{pmatrix} 0 \\ 1 \end{pmatrix}
  =
  \begin{pmatrix} 1 \\ 1 \end{pmatrix}
\]
は$x$軸にも$y$軸にも属さないので，合併は加法について閉じていない．

\prref{和と線型包}は，合併そのものではなく，和$W_1 + W_2$こそが「$W_1$と$W_2$の両方を含む最小の線型部分空間」であることを示している．
\end{remark}

\phantomsection
\subsection{直和}

本項では「直和」とよばれる二つの構成を扱う．
まず，二つの線型空間を組にして一つの線型空間を作る方法として，外部直和を定義する．

\begin{definition}{直和（外部直和）}{直和}
  $V,W$を体$K$上の線型空間とする．

  直積集合
  \[
  V\times W
  =
  \{
  (\bm{v},\bm{w})
  \mid
  \bm{v}\in V,\
  \bm{w}\in W
  \}
  \]
  を考え，
  加法およびスカラー倍を
  \[
  (\bm{v}_1,\bm{w}_1)+(\bm{v}_2,\bm{w}_2)
  =
  (\bm{v}_1+\bm{v}_2,\bm{w}_1+\bm{w}_2),
  \]
  \[
  \alpha(\bm{v},\bm{w})
  =
  (\alpha \bm{v},\alpha \bm{w})
  \]
  によって定める．

  このとき，線型空間の各公理を成分ごとに確認することにより，
  $V\times W$はこれらの演算に関して体$K$上の線型空間となる．

  これを
  $V$と$W$の\textbf{直和}\index{ちょくわ@直和}（\textbf{外部直和}\index{がいぶちょくわ@外部直和}）といい，
  \[
  V\oplus W
  \]
  と表す．
\end{definition}

\begin{remark}[外部直和の直観]
直和
$V\oplus W$
は，
二つの線型空間を互いに独立な成分として並べた線型空間と考えることができる．
\end{remark}

一方，一つの線型空間$V$の「内部」で，二つの線型部分空間が$V$を過不足なく分担している状況を表す概念として，内部直和を定義する．

\begin{definition}{内部直和}{内部直和}
  $W_1, W_2$を体$K$上の線型空間$V$の線型部分空間とする．
  \[
    V = W_1 + W_2,
    \qquad
    W_1 \cap W_2 = \{\bm{0}_V\}
  \]
  がともに成り立つとき（和$W_1 + W_2$については\deref{共通部分・和}を参照），$V$は$W_1$と$W_2$の\textbf{内部直和}\index{ないぶちょくわ@内部直和}であるといい，
  \[
    V = W_1 \oplus W_2
  \]
  と表す．
\end{definition}

内部直和は，「ベクトルの分解の一意性」によって特徴づけられる．

\begin{prop}{内部直和と分解の一意性}{内部直和と分解の一意性}
  $W_1, W_2$を線型空間$V$の線型部分空間とする．このとき，次の二つは同値である：
  \begin{enumerate}[label=(\roman*),leftmargin=3.2em]
    \item $V = W_1 \oplus W_2$（内部直和）である．
    \item 任意の$\bm{v} \in V$は，
    \[
      \bm{v} = \bm{w}_1 + \bm{w}_2
      \quad (\bm{w}_1 \in W_1,~ \bm{w}_2 \in W_2)
    \]
    の形にただ一通りに表される．
  \end{enumerate}
\end{prop}

\begin{leftbar}
\begin{proof}
  まず，(i)を仮定して(ii)を示す．
  $V = W_1 + W_2$であるから，任意の$\bm{v} \in V$は$\bm{v} = \bm{w}_1 + \bm{w}_2$（$\bm{w}_1 \in W_1$，$\bm{w}_2 \in W_2$）の形に表せる．
  いま，
  \[
    \bm{v} = \bm{w}_1 + \bm{w}_2 = \bm{w}'_1 + \bm{w}'_2
    \quad (\bm{w}_i, \bm{w}'_i \in W_i)
  \]
  と二通りに表せたとすると，移項して
  \[
    \bm{w}_1 - \bm{w}'_1 = \bm{w}'_2 - \bm{w}_2
  \]
  を得る．左辺は$W_1$に，右辺は$W_2$に属するから，両辺が表すベクトルは$W_1 \cap W_2 = \{\bm{0}_V\}$に属する．
  よって$\bm{w}_1 = \bm{w}'_1$かつ$\bm{w}_2 = \bm{w}'_2$であり，表し方はただ一通りである．

  逆に，(ii)を仮定して(i)を示す．
  任意の$\bm{v} \in V$が$\bm{w}_1 + \bm{w}_2$の形に表せることから，$V = W_1 + W_2$である．
  次に，$\bm{v} \in W_1 \cap W_2$をとると，
  \[
    \bm{v} = \bm{v} + \bm{0}_V = \bm{0}_V + \bm{v}
  \]
  は，いずれも「$W_1$の元と$W_2$の元の和」としての$\bm{v}$の表示である．
  表し方の一意性より$\bm{v} = \bm{0}_V$となるから，$W_1 \cap W_2 = \{\bm{0}_V\}$である．
  以上の考察により，命題が示された．
\end{proof}
\end{leftbar}

\begin{remark}[内部直和と外部直和の同一視]
  外部直和と内部直和は異なる構成であるが，同じ記号$\oplus$で表される．
  これは，$V = W_1 \oplus W_2$（内部直和）のとき，外部直和$W_1 \oplus W_2 = W_1 \times W_2$の元$(\bm{w}_1, \bm{w}_2)$に$\bm{w}_1 + \bm{w}_2 \in V$を対応させる写像が全単射となり（全射性は$V = W_1 + W_2$から，単射性は\prref{内部直和と分解の一意性}の分解の一意性から従う），さらに和とスカラー倍も保たれるため，両者を自然に同一視できるからである．
\end{remark}

\phantomsection
\section{基底と次元}

\phantomsection
\subsection{基底}

前節で，ベクトルの集合$S$から生成される線型部分空間$\Span(S)$を導入した．
一方，\partref{線型独立性と階数}では，ベクトルの組に「無駄がない」ことを表す概念として線型独立を導入した（\deref{線型独立}）．

この二つを組み合わせよう．
すなわち，「$V$全体を生成し，かつ無駄がない」ベクトルの組を考える．

\begin{definition}{基底}{基底}
  $V$を体$K$上の線型空間とする．
  $V$のベクトルの組$\mathcal{B} = (\bm{v}_1, \bm{v}_2, \dots, \bm{v}_n)$が次の二つの条件を満たすとき，
  $\mathcal{B}$を$V$の\textbf{基底}\index{きてい@基底}という．
  \begin{enumerate}[label=(\roman*),leftmargin=3.2em]
    \item \textbf{線型独立性}

    $\bm{v}_1, \bm{v}_2, \dots, \bm{v}_n$は線型独立である．

    \item \textbf{生成性}
    \[
      \Span(\bm{v}_1, \bm{v}_2, \dots, \bm{v}_n) = V
    \]
    である．
  \end{enumerate}

  また，$V = \{\bm{0}_V\}$のときは，空の組$(\ )$を$V$の基底と定める．
\end{definition}

\begin{remark}[基底を「組」として定義する理由]
  基底を「集合」ではなく「組」（順序のついた有限列）として定義していることに注意してほしい．
  順序を込めて考えるのは，のちに定義する座標（\deref{座標ベクトル}）が「第何成分か」を指定して初めて意味をもつからである．
  文献によっては基底を集合として定義し，座標を考えるときにあらためて順序を固定する流儀もある．

  また，$V = \{\bm{0}_V\}$に対して空の組を基底と定めたのは，$\Span(\varnothing) = \{\bm{0}_V\}$という約束（\deref{線型包}）と整合させるためである．
  「$\bm{0}_V$のみからなる組」を基底とすることはできない．
  実際，$1_K \bm{0}_V = \bm{0}_V$は非自明な線型関係であり，$(\bm{0}_V)$は線型従属だからである．
\end{remark}

\begin{example}{基底の例}{基底の例}
  \begin{enumerate}[label=(\arabic*),leftmargin=3.2em]
    \item $K^n$において，標準基底ベクトルの組$(\bm{e}_1, \bm{e}_2, \dots, \bm{e}_n)$は基底である．
    これを$K^n$の\textbf{標準基底}\index{ひょうじゅんきてい@標準基底}という．

    \item $K^2$において，
    \[
      \bm{v}_1 = \begin{pmatrix} 1 \\ 1 \end{pmatrix}, \quad
      \bm{v}_2 = \begin{pmatrix} 1 \\ -1 \end{pmatrix}
    \]
    の組も基底である．このように，基底の取り方は一通りではない．

    \item $K$の元を成分とする$m \times n$行列全体のなす線型空間$M_{m,n}(K)$において，$(i,j)$成分のみが$1$で他が$0$である行列$E_{ij}$の組は基底である\footnote{この$E_{ij}$は\textbf{行列単位}とよばれる．単位行列$E_n$（\deref{零行列と単位行列}）と記号が似ているが別物である．なお，正方行列（$m = n$）の場合には$E_n = E_{11} + E_{22} + \cdots + E_{nn}$という関係がある．}．

    \item 次数が$n$以下の実数係数多項式全体のなす線型空間において，$(1, x, x^2, \dots, x^n)$は基底である．
  \end{enumerate}
\end{example}

\begin{definition}{有限生成}{有限生成}
  線型空間$V$に対して，有限個のベクトル$\bm{v}_1, \bm{v}_2, \dots, \bm{v}_n \in V$が存在して
  \[
    V = \Span(\bm{v}_1, \bm{v}_2, \dots, \bm{v}_n)
  \]
  となるとき，$V$は\textbf{有限生成}\index{ゆうげんせいせい@有限生成}であるという．
\end{definition}

\begin{mycolumn}
  有限生成でない線型空間も存在する．
  たとえば，実数係数多項式全体のなす線型空間$\mathbb{R}[x]$は有限生成ではない．
  実際，有限個の多項式$f_1, \dots, f_n$に対して，十分大きい整数$d$をとれば，$f_1, \dots, f_n$はいずれも次数$d$以下の多項式である．次数$d$以下の多項式（零多項式を含む）の線型結合はふたたび次数$d$以下の多項式であるから，たとえば$x^{d+1}$は$f_1, \dots, f_n$の線型結合として表せず，より高次の多項式は生成できない．

  このような空間にも「基底」を考えることはできるが，そこでは無限個のベクトルからなる集合を扱う必要があり，線型独立性も「その有限部分集合がつねに線型独立である」と定義し直さなければならない．
  さらに，一般の線型空間に基底が存在することの証明には選択公理（ツォルンの補題）が必要となる．

  本書では，無限次元の線型空間を本格的に扱うことはせず，有限生成性（のちの言葉でいえば有限次元性）を仮定する命題では，その都度これを明記する．
  また，本書で単に「基底」というときは，つねに\deref{基底}の意味での，有限個のベクトルからなる順序付きの組を指すものとし，無限集合としての基底は正式な定義の対象としない．
\end{mycolumn}

基底のもっとも重要な性質は，「表し方がただ一通りに定まる」ことである．

\begin{theorem}{基底による表示の一意性}{基底による表示の一意性}
  $V$を体$K$上の線型空間とし，$\bm{v}_1, \bm{v}_2, \dots, \bm{v}_n \in V$とする．
  このとき，$(\bm{v}_1, \bm{v}_2, \dots, \bm{v}_n)$が$V$の基底であるための必要十分条件は，
  任意の$\bm{v} \in V$が
  \[
    \bm{v} = x_1 \bm{v}_1 + x_2 \bm{v}_2 + \cdots + x_n \bm{v}_n
    \quad (x_i \in K)
  \]
  の形にただ一通りに表されることである．
\end{theorem}

\begin{leftbar}
\begin{proof}
  まず，$(\bm{v}_1, \dots, \bm{v}_n)$が基底であるとする．
  生成性より，任意の$\bm{v} \in V$は上の形に表せる．
  いま，
  \[
    \bm{v} = \sum_{i=1}^{n} x_i \bm{v}_i = \sum_{i=1}^{n} y_i \bm{v}_i
  \]
  と二通りに表せたとすると，辺々引いて
  \[
    \sum_{i=1}^{n} (x_i - y_i) \bm{v}_i = \bm{0}_V
  \]
  を得る．線型独立性より$x_i - y_i = 0$，すなわち$x_i = y_i$である．
  よって表し方は一通りである．

  逆に，任意の$\bm{v} \in V$がただ一通りに表されるとする．
  表せること自体から$V = \Span(\bm{v}_1, \dots, \bm{v}_n)$であり，生成性が成り立つ．
  また，
  \[
    \sum_{i=1}^{n} \lambda_i \bm{v}_i = \bm{0}_V
  \]
  とすると，$\bm{0}_V$の表し方としては$\lambda_1 = \cdots = \lambda_n = 0$とするものがすでにある．
  一意性より$\lambda_i = 0$であり，線型独立性が成り立つ．
  以上の考察により，定理が示された．
\end{proof}
\end{leftbar}

\thref{基底による表示の一意性}により，基底を一つ固定すると，$V$のベクトルに$n$個の数の組が一意に対応する．

\begin{definition}{座標ベクトル}{座標ベクトル}
  $\mathcal{B} = (\bm{v}_1, \bm{v}_2, \dots, \bm{v}_n)$を線型空間$V$の基底とする．
  $\bm{v} \in V$に対して，\thref{基底による表示の一意性}により一意に定まる$x_1, x_2, \dots, x_n \in K$を用いて
  \[
    [\bm{v}]_{\mathcal{B}}
    = \begin{pmatrix} x_1 \\ x_2 \\ \vdots \\ x_n \end{pmatrix} \in K^n
  \]
  と定め，これを基底$\mathcal{B}$に関する$\bm{v}$の\textbf{座標ベクトル}\index{ざひょうべくとる@座標ベクトル}という．
\end{definition}

\begin{prop}{座標ベクトルの性質}{座標ベクトルの性質}
  $\mathcal{B}$を線型空間$V$の基底とする．
  このとき，写像
  \[
    V \to K^n, \quad \bm{v} \mapsto [\bm{v}]_{\mathcal{B}}
  \]
  は全単射であり，任意の$\bm{v}, \bm{w} \in V$および任意の$\alpha \in K$に対して
  \[
    [\bm{v} + \bm{w}]_{\mathcal{B}} = [\bm{v}]_{\mathcal{B}} + [\bm{w}]_{\mathcal{B}},
    \quad
    [\alpha \bm{v}]_{\mathcal{B}} = \alpha [\bm{v}]_{\mathcal{B}}
  \]
  が成り立つ．
\end{prop}

\begin{leftbar}
\begin{proof}
  $\mathcal{B} = (\bm{v}_1, \dots, \bm{v}_n)$とする．
  全射性は生成性から，単射性は表示の一意性から従う．

  $[\bm{v}]_{\mathcal{B}} = (x_i)$，$[\bm{w}]_{\mathcal{B}} = (y_i)$とすると，
  \[
    \bm{v} + \bm{w}
    = \sum_{i=1}^{n} x_i \bm{v}_i + \sum_{i=1}^{n} y_i \bm{v}_i
    = \sum_{i=1}^{n} (x_i + y_i) \bm{v}_i
  \]
  であり，表示の一意性より$[\bm{v} + \bm{w}]_{\mathcal{B}} = (x_i + y_i)$である．
  スカラー倍についても同様である．
  これが証明すべきことであった．
\end{proof}
\end{leftbar}

\begin{remark}[基底による$K^n$との同一視]
  \prref{座標ベクトルの性質}は，「基底を一つ選べば，抽象的な線型空間$V$は数ベクトル空間$K^n$と同じものとみなせる」ことを述べている．
  たとえば，次数$2$以下の実数係数多項式全体は，基底$(1, x, x^2)$を選ぶことで$\mathbb{R}^3$と同一視され，
  \[
    a + bx + cx^2 \longleftrightarrow \begin{pmatrix} a \\ b \\ c \end{pmatrix}
  \]
  という対応のもとで，多項式の計算が数ベクトルの計算に翻訳される．

  ただし，この同一視は基底の選び方に依存する．
  「どの基底を選んでも変わらない性質」と「基底に依存する表示」を区別することが，以下の議論の主題である．
\end{remark}

\phantomsection
\subsection{次元}

基底の取り方は一通りではないが，「基底に含まれるベクトルの個数」は基底の取り方によらない．
これを示すのが本項の目標である．

\partref{線型独立性と階数}では，数ベクトル空間$K^m$に限って「線型独立な組の拡大」と「取替補題」を示した（\lemref{線型独立な組の拡大}，\lemref{取替補題}）．
本項では，同じ議論が一般の線型空間においてもそのまま通用することを確認する．
これは単なる繰り返しではなく，\partref{線型独立性と階数}の具体的な議論を抽象的な枠組みへ再構成する作業である．

まず，準備として次の補題を用意する．

\begin{lemma}{線型独立な組の拡大}{一般の線型独立な組の拡大}
  $\bm{v}_1, \bm{v}_2, \dots, \bm{v}_k$を線型空間$V$の線型独立なベクトルとし，
  \[
    \bm{w} \notin \Span(\bm{v}_1, \bm{v}_2, \dots, \bm{v}_k)
  \]
  とする．このとき，$\bm{v}_1, \bm{v}_2, \dots, \bm{v}_k, \bm{w}$は線型独立である．
\end{lemma}

\begin{leftbar}
\begin{proof}
  \[
    \lambda_1 \bm{v}_1 + \cdots + \lambda_k \bm{v}_k + \mu \bm{w} = \bm{0}_V
  \]
  とする．もし$\mu \ne 0$ならば，
  \[
    \bm{w} = \sum_{i=1}^{k} \left(-\frac{\lambda_i}{\mu}\right) \bm{v}_i
    \in \Span(\bm{v}_1, \dots, \bm{v}_k)
  \]
  となり仮定に反する．よって$\mu = 0$である．
  このとき$\sum_{i=1}^{k} \lambda_i \bm{v}_i = \bm{0}_V$となるので，$\bm{v}_1, \dots, \bm{v}_k$の線型独立性より$\lambda_1 = \cdots = \lambda_k = 0$である．
  これが証明すべきことであった．
\end{proof}
\end{leftbar}

次が，次元の概念を支える中心的な補題である．

\begin{lemma}{取替補題}{一般の取替補題}
  $V$を線型空間とし，$\bm{v}_1, \bm{v}_2, \dots, \bm{v}_n \in V$とする．
  $\bm{w}_1, \bm{w}_2, \dots, \bm{w}_m \in \Span(\bm{v}_1, \bm{v}_2, \dots, \bm{v}_n)$が線型独立ならば，
  \[
    m \le n
  \]
  が成り立つ．
\end{lemma}

\begin{leftbar}
\begin{proof}
  $W = \Span(\bm{v}_1, \dots, \bm{v}_n)$とおく．
  $k = 0, 1, \dots, m$に対して，次の主張$(\ast_k)$を$k$についての帰納法で示す．

  \begin{quote}
    $(\ast_k)$：$k \le n$であり，かつ$\bm{v}_1, \dots, \bm{v}_n$の番号を適当に付け替えることで
    \[
      W = \Span(\bm{w}_1, \dots, \bm{w}_k, \bm{v}_{k+1}, \dots, \bm{v}_n)
    \]
    とできる．
  \end{quote}

  $k = 0$のときは$W$の定義そのものであり，成り立つ．

  $k < m$として$(\ast_k)$を仮定する．
  $\bm{w}_{k+1} \in W$であるから，帰納法の仮定より
  \begin{equation}
    \bm{w}_{k+1}
    = \sum_{i=1}^{k} a_i \bm{w}_i + \sum_{j=k+1}^{n} b_j \bm{v}_j
    \label{eq:一般の取替補題の展開}
  \end{equation}
  と表せる．

  ここで，$b_{k+1} = \cdots = b_n = 0$と仮定してみる（$k = n$のときは第2の和が空なので自動的にこの場合である）．
  このとき\eqref{eq:一般の取替補題の展開}は
  \[
    a_1 \bm{w}_1 + \cdots + a_k \bm{w}_k + (-1)\bm{w}_{k+1} = \bm{0}_V
  \]
  と書き直せる．$\bm{w}_{k+1}$の係数は$-1 \ne 0$であるから，これは$\bm{w}_1, \dots, \bm{w}_{k+1}$の非自明な線型関係であり，線型独立性に矛盾する．

  よって$k < n$，すなわち$k + 1 \le n$であり，さらに$b_j \ne 0$となる$j$が存在する．
  残っている$\bm{v}_{k+1}, \dots, \bm{v}_n$の番号を付け替えて$b_{k+1} \ne 0$としてよい．
  このとき\eqref{eq:一般の取替補題の展開}を$\bm{v}_{k+1}$について解くと，
  \[
    \bm{v}_{k+1}
    = \frac{1}{b_{k+1}}\left( \bm{w}_{k+1} - \sum_{i=1}^{k} a_i \bm{w}_i - \sum_{j=k+2}^{n} b_j \bm{v}_j \right)
  \]
  となるので，
  \[
    \bm{v}_{k+1} \in \Span(\bm{w}_1, \dots, \bm{w}_{k+1}, \bm{v}_{k+2}, \dots, \bm{v}_n)
  \]
  である．
  したがって，右辺の線型包は$\bm{w}_1, \dots, \bm{w}_k, \bm{v}_{k+1}, \dots, \bm{v}_n$をすべて含む線型部分空間であり，\prref{線型包の最小性}より
  \[
    W = \Span(\bm{w}_1, \dots, \bm{w}_k, \bm{v}_{k+1}, \dots, \bm{v}_n)
    \subset \Span(\bm{w}_1, \dots, \bm{w}_{k+1}, \bm{v}_{k+2}, \dots, \bm{v}_n)
  \]
  を得る．逆の包含は明らかであるから，$(\ast_{k+1})$が成り立つ．

  以上より$(\ast_m)$が成り立ち，とくに$m \le n$である．
  これが証明すべきことであった．
\end{proof}
\end{leftbar}

\begin{theorem}{基底の個数の一意性}{基底の個数の一意性}
  $V$を線型空間とし，$(\bm{v}_1, \dots, \bm{v}_n)$と$(\bm{w}_1, \dots, \bm{w}_m)$がともに$V$の基底であるとする．
  このとき，$m = n$である．
\end{theorem}

\begin{leftbar}
\begin{proof}
  $(\bm{v}_1, \dots, \bm{v}_n)$の生成性より$\bm{w}_1, \dots, \bm{w}_m \in \Span(\bm{v}_1, \dots, \bm{v}_n)$であり，$\bm{w}_1, \dots, \bm{w}_m$は線型独立であるから，\lemref{一般の取替補題}より$m \le n$である．

  $(\bm{v}_1, \dots, \bm{v}_n)$と$(\bm{w}_1, \dots, \bm{w}_m)$の役割を入れ替えて同じ議論を行えば$n \le m$を得る．
  よって$m = n$である．
  これが証明すべきことであった．
\end{proof}
\end{leftbar}

\thref{基底の個数の一意性}により，次の定義が意味をもつ．

\begin{definition}{次元}{次元}
  線型空間$V$が基底$(\bm{v}_1, \bm{v}_2, \dots, \bm{v}_n)$をもつとき，
  $n$を$V$の\textbf{次元}\index{じげん@次元}といい，
  \[
    \dim_K V = n
  \]
  と表す．体$K$が文脈から明らかなときは，単に$\dim V$とかく．

  とくに，$\dim \{\bm{0}_V\} = 0$である．
\end{definition}

\begin{remark}[「次元」が意味をもつために]
  「次元が$n$である」という主張は，
  \begin{itemize}
    \item 基底が存在すること（存在性）
    \item その個数が基底の取り方によらないこと（一意性）
  \end{itemize}
  の二つがそろって初めて意味をもつ．
  \thref{基底の個数の一意性}は後者を保証するものであり，前者は，有限生成な線型空間については次の\thref{基底の存在}が保証する（一般の線型空間の場合については，\deref{有限生成}のあとのコラムで触れた）．

  定義の順序としては，まず基底を定め，つぎにその個数が基底の取り方によらないことを示し，そのうえで次元を定義する，という流れになっている．次元を先に定義することはできない点に注意してほしい．
\end{remark}

次元に関連して，用語を一つ補っておく．

\begin{definition}{有限次元と無限次元}{有限次元と無限次元}
  有限個のベクトルからなる基底をもつ線型空間を\textbf{有限次元}\index{ゆうげんじげん@有限次元}線型空間という．
  有限次元でない線型空間を\textbf{無限次元}\index{むげんじげん@無限次元}線型空間という．
\end{definition}

のちの\thref{基底の存在}で示すように，有限生成な線型空間はつねに（有限個のベクトルからなる）基底をもつ．
逆に，有限個のベクトルからなる基底をもつ線型空間は，明らかに有限生成である．
したがって「有限生成」と「有限次元」は同じクラスの線型空間を指す言葉である．本書では，有限次元性（有限生成性）を仮定する命題では，その都度これを明記する（\deref{有限生成}のあとのコラムも参照）．

\begin{example}{次元の例}{次元の例}
  \exref{基底の例}より，
  \[
    \dim_K K^n = n,
    \quad
    \dim_K M_{m,n}(K) = mn,
    \quad
    \dim_{\mathbb{R}} \mathbb{R}[x]_{\le n} = n + 1
  \]
  である．ここで$M_{m,n}(K)$は$K$の元を成分とする$m \times n$行列全体の集合（\deref{行列全体の集合}），$\mathbb{R}[x]_{\le n}$は次数が$n$以下の実数係数多項式全体を表す．

  なお，$\mathbb{C}$は$\mathbb{C}$上の線型空間としては$\dim_{\mathbb{C}} \mathbb{C} = 1$であるが，$\mathbb{R}$上の線型空間とみなすと$(1, i)$が基底となり$\dim_{\mathbb{R}} \mathbb{C} = 2$である．
  次元は，どの体上の線型空間とみるかに依存する．
\end{example}

\phantomsection
\subsection{基底の存在と延長}

\begin{theorem}{基底の存在}{基底の存在}
  $V$を有限生成な線型空間とすると，$V$は基底をもつ．
  より詳しく，$V = \Span(\bm{v}_1, \dots, \bm{v}_n)$ならば，$\bm{v}_1, \dots, \bm{v}_n$のうちの適当な部分の組が$V$の基底となる．
\end{theorem}

\begin{leftbar}
\begin{proof}
  $\bm{v}_1, \dots, \bm{v}_n$が線型独立ならば，これがそのまま基底である．

  線型独立でないとすると，\prref{線型従属の言い換え}より，ある$i$について$\bm{v}_i$は残りのベクトルの線型結合として表せる．
  このとき，
  \[
    \Span(\bm{v}_1, \dots, \bm{v}_n)
    = \Span(\bm{v}_1, \dots, \widehat{\bm{v}_i}, \dots, \bm{v}_n)
  \]
  が成り立つ（$\widehat{\phantom{x}}$はその項を除くことを表す）．
  実際，右辺は左辺に含まれ，逆に$\bm{v}_i$が右辺に属するので，\prref{線型包の最小性}より左辺は右辺に含まれる．

  この操作でベクトルの個数は$1$減るので，有限回で線型独立な組に到達する．
  すべてのベクトルが除かれた場合は$V = \Span(\varnothing) = \{\bm{0}_V\}$であり，空の組が基底である．
  これが証明すべきことであった．
\end{proof}
\end{leftbar}

\begin{theorem}{基底の延長}{基底の延長}
  $V$を有限生成な線型空間とし，$\bm{w}_1, \bm{w}_2, \dots, \bm{w}_k \in V$を線型独立なベクトルとする．
  このとき，適当な$\bm{u}_1, \dots, \bm{u}_l \in V$を付け加えて
  \[
    (\bm{w}_1, \dots, \bm{w}_k, \bm{u}_1, \dots, \bm{u}_l)
  \]
  を$V$の基底とすることができる．
\end{theorem}

\begin{leftbar}
\begin{proof}
  $V = \Span(\bm{v}_1, \dots, \bm{v}_n)$とする．

  $\bm{v}_1, \bm{v}_2, \dots, \bm{v}_n$を順に見ていき，
  「その時点までにできている組の線型包に属さないもの」だけを組に付け加える，という操作を行う．
  \lemref{一般の線型独立な組の拡大}より，付け加えるたびに組は線型独立なままである．
  最終的に得られる組を
  \[
    (\bm{w}_1, \dots, \bm{w}_k, \bm{u}_1, \dots, \bm{u}_l)
  \]
  とおく．

  構成法より，各$\bm{v}_j$はこの組の線型包に属する（付け加えられた場合は自明であり，付け加えられなかった場合はその時点の線型包に属していたからである）．
  よって\prref{線型包の最小性}より
  \[
    V = \Span(\bm{v}_1, \dots, \bm{v}_n)
    \subset \Span(\bm{w}_1, \dots, \bm{w}_k, \bm{u}_1, \dots, \bm{u}_l)
    \subset V
  \]
  となり，生成性が成り立つ．線型独立性はすでに示したので，この組は$V$の基底である．
  これが証明すべきことであった．
\end{proof}
\end{leftbar}

\begin{cor}{次元による基底の判定}{次元による基底の判定}
  $\dim V = n$とし，$\bm{v}_1, \bm{v}_2, \dots, \bm{v}_n \in V$とする（個数が$n$であることに注意）．
  このとき，次の三つは同値である．
  \begin{enumerate}[label=(\roman*),leftmargin=3.2em]
    \item $\bm{v}_1, \dots, \bm{v}_n$は線型独立である．
    \item $\Span(\bm{v}_1, \dots, \bm{v}_n) = V$である．
    \item $(\bm{v}_1, \dots, \bm{v}_n)$は$V$の基底である．
  \end{enumerate}
\end{cor}

\begin{leftbar}
\begin{proof}
  (iii)から(i)，(ii)が従うのは基底の定義そのものである．

  (i)$\limp$(iii)：\thref{基底の延長}より，この組はベクトルを付け加えて基底にできる．
  しかし\thref{基底の個数の一意性}より基底の個数は$n$でなければならないから，付け加えるベクトルは$0$個である．

  (ii)$\limp$(iii)：\thref{基底の存在}より，この組から適当な部分を取り出して基底にできる．
  同様に，基底の個数は$n$であるから，取り除くベクトルは$0$個である．
  これが証明すべきことであった．
\end{proof}
\end{leftbar}

\begin{remark}[次元による基底の判定の注意]
  \coref{次元による基底の判定}は実用上たいへん便利である．
  たとえば$\mathbb{R}^3$のベクトルが$3$個与えられたとき，基底であることを確かめるには線型独立性だけを調べればよく，生成性を別に確認する必要はない．

  ただし，これは「個数が次元と一致している」場合に限られる．
  個数が$n$と異なる組については，線型独立性と生成性は独立な条件であり，一方から他方は従わない．
\end{remark}

\begin{prop}{部分空間の次元}{部分空間の次元}
  $V$を有限生成な線型空間とし，$W$を$V$の線型部分空間とする．
  このとき$W$も有限生成であり，
  \[
    \dim W \le \dim V
  \]
  が成り立つ．さらに，$\dim W = \dim V$であることと$W = V$であることは同値である．
\end{prop}

\begin{leftbar}
\begin{proof}
  $\dim V = n$とする．
  \lemref{一般の取替補題}より，$W$の線型独立な組に含まれるベクトルの個数は高々$n$である．
  空の組も線型独立とみなすと，$W$の線型独立な組の長さは$0, 1, \dots, n$のいずれかであるから，その最大値$r$が存在する．
  そこで，長さが最大である$W$の線型独立な組を$(\bm{w}_1, \dots, \bm{w}_r)$とする（$r \le n$．$W = \{\bm{0}_V\}$のときは$r = 0$であり，空の組をとることになる）．

  もし$\Span(\bm{w}_1, \dots, \bm{w}_r) \ne W$ならば，$\bm{w} \in W$であって$\bm{w} \notin \Span(\bm{w}_1, \dots, \bm{w}_r)$となるものがとれる．
  \lemref{一般の線型独立な組の拡大}より$\bm{w}_1, \dots, \bm{w}_r, \bm{w}$は線型独立となり，$r$の最大性に矛盾する．
  よって$(\bm{w}_1, \dots, \bm{w}_r)$は$W$の基底であり，$\dim W = r \le n$である．

  最後の主張について，$W = V$ならば$\dim W = \dim V$は明らかである．
  逆に$\dim W = n$とすると，$W$の基底は$V$の$n$個の線型独立なベクトルであるから，\coref{次元による基底の判定}より$V$の基底でもある．
  よって
  \[
    V = \Span(\bm{w}_1, \dots, \bm{w}_n) \subset W \subset V
  \]
  となり，$W = V$である．
  これが証明すべきことであった．
\end{proof}
\end{leftbar}

\begin{theorem}{次元公式}{次元公式}
  $V$を有限生成な線型空間とし，$W_1, W_2$を$V$の線型部分空間とする．
  このとき，
  \[
    \dim(W_1 + W_2) + \dim(W_1 \cap W_2) = \dim W_1 + \dim W_2
  \]
  が成り立つ．
\end{theorem}

\begin{leftbar}
\begin{proof}
  $\dim(W_1 \cap W_2) = r$とし，$W_1 \cap W_2$の基底を$(\bm{u}_1, \dots, \bm{u}_r)$とする．
  \thref{基底の延長}より，これを延長して
  \[
    W_1 \text{の基底} \quad (\bm{u}_1, \dots, \bm{u}_r, \bm{x}_1, \dots, \bm{x}_s),
  \]
  \[
    W_2 \text{の基底} \quad (\bm{u}_1, \dots, \bm{u}_r, \bm{y}_1, \dots, \bm{y}_t)
  \]
  がとれる．すなわち$\dim W_1 = r + s$，$\dim W_2 = r + t$である．

  ここで，
  \[
    (\bm{u}_1, \dots, \bm{u}_r, \bm{x}_1, \dots, \bm{x}_s, \bm{y}_1, \dots, \bm{y}_t)
  \]
  が$W_1 + W_2$の基底であることを示す．

  生成性について，まず，上の組の各ベクトルは$W_1$または$W_2$に属するから$W_1 + W_2$に属し，$W_1 + W_2$は線型部分空間であるから，\prref{線型包の最小性}より上の組の線型包は$W_1 + W_2$に含まれる．
  逆に，$W_1$も$W_2$も上の組の線型包に含まれるから，\prref{和と線型包}の$W_1 + W_2 = \Span(W_1 \cup W_2)$と\prref{線型包の最小性}より，$W_1 + W_2$は上の組の線型包に含まれる．
  したがって，上の組の線型包は$W_1 + W_2$に等しい．

  線型独立性について，
  \begin{equation}
    \sum_{i=1}^{r} a_i \bm{u}_i + \sum_{j=1}^{s} b_j \bm{x}_j + \sum_{k=1}^{t} c_k \bm{y}_k = \bm{0}_V
    \label{eq:次元公式の線型関係}
  \end{equation}
  とする．このとき
  \[
    \bm{z} := \sum_{k=1}^{t} c_k \bm{y}_k
    = -\left( \sum_{i=1}^{r} a_i \bm{u}_i + \sum_{j=1}^{s} b_j \bm{x}_j \right)
  \]
  とおくと，左辺の表示より$\bm{z} \in W_2$，右辺の表示より$\bm{z} \in W_1$であるから，$\bm{z} \in W_1 \cap W_2$である．
  よって
  \[
    \bm{z} = \sum_{i=1}^{r} d_i \bm{u}_i
  \]
  と表せるので，
  \[
    \sum_{i=1}^{r} d_i \bm{u}_i - \sum_{k=1}^{t} c_k \bm{y}_k = \bm{0}_V
  \]
  となる．$(\bm{u}_1, \dots, \bm{u}_r, \bm{y}_1, \dots, \bm{y}_t)$は$W_2$の基底であるから線型独立であり，とくに$c_1 = \cdots = c_t = 0$である．

  これを\eqref{eq:次元公式の線型関係}に代入すると
  \[
    \sum_{i=1}^{r} a_i \bm{u}_i + \sum_{j=1}^{s} b_j \bm{x}_j = \bm{0}_V
  \]
  となり，$W_1$の基底の線型独立性より$a_i = 0$，$b_j = 0$である．

  以上より$\dim(W_1 + W_2) = r + s + t$であり，
  \[
    (r + s + t) + r = (r + s) + (r + t)
  \]
  から主張が従う．
  これが証明すべきことであった．
\end{proof}
\end{leftbar}

\begin{remark}[次元公式の読み方]
  次元公式は，有限集合の元の個数についての公式
  \[
    |A \cup B| + |A \cap B| = |A| + |B|
  \]
  と同じ形をしている．
  ただし，線型空間では合併$W_1 \cup W_2$ではなく和$W_1 + W_2$を考える点が異なる（\prref{和と線型包}のあとのRemarkを参照）．

  とくに$W_1 \cap W_2 = \{\bm{0}_V\}$のときは
  \[
    \dim(W_1 + W_2) = \dim W_1 + \dim W_2
  \]
  となる．この場合，和$W_1 + W_2$は（それ自身を一つの線型空間とみて）$W_1$と$W_2$の内部直和$W_1 \oplus W_2$である（\deref{内部直和}）．
  このとき各元が$\bm{w}_1 + \bm{w}_2$（$\bm{w}_i \in W_i$）の形にただ一通りに表されることは，\prref{内部直和と分解の一意性}で見た通りである．
  すなわち上の等式は，「内部直和の次元は各成分の次元の和である」と読むことができる．
\end{remark}

本項の締めくくりとして，もう一つの次元公式を証明しよう．
\partref{線型写像}で，線型写像$f \colon K^n \to K^m$に対して像$\Im f$と核$\Ker f$を定義した（\deref{像と核}）．
本章では，これらがそれぞれ$K^m$，$K^n$の線型部分空間であることを確かめた（\prref{像と核の線型部分空間性}）ので，その次元を考えることができる．

\begin{theorem}{線型写像の次元公式}{線型写像の次元公式}
  $f \colon K^n \to K^m$を線型写像とする．このとき，
  \[
    \dim \Ker f + \dim \Im f = n
  \]
  が成り立つ．
\end{theorem}

\begin{leftbar}
\begin{proof}
  $K^n$は有限生成であり（標準基底，\exref{基底の例}），$\Ker f$はその線型部分空間であるから，\prref{部分空間の次元}より$\Ker f$も有限生成である．
  $\dim \Ker f = k$とおき，$\Ker f$の基底を$(\bm{u}_1, \dots, \bm{u}_k)$とする（$k = 0$のときは空の組をとる）．
  \thref{基底の延長}より，これを延長して$K^n$の基底
  \[
    (\bm{u}_1, \dots, \bm{u}_k, \bm{w}_1, \dots, \bm{w}_l)
  \]
  がとれる．$\dim K^n = n$（\exref{次元の例}）であるから，\thref{基底の個数の一意性}より$k + l = n$である．

  ここで，$(f(\bm{w}_1), \dots, f(\bm{w}_l))$が$\Im f$の基底であることを示す．

  生成性について，任意の$\bm{v} \in \Im f$は，ある$\bm{x} \in K^n$を用いて$\bm{v} = f(\bm{x})$と表せる．
  $\bm{x}$を上の$K^n$の基底で
  \[
    \bm{x} = \sum_{i=1}^{k} a_i \bm{u}_i + \sum_{j=1}^{l} b_j \bm{w}_j
  \]
  と表すと，\prref{線型写像と線型結合}と$f(\bm{u}_i) = \bm{0}$（$\bm{u}_i \in \Ker f$）より，
  \[
    \bm{v} = f(\bm{x})
    = \sum_{i=1}^{k} a_i f(\bm{u}_i) + \sum_{j=1}^{l} b_j f(\bm{w}_j)
    = \sum_{j=1}^{l} b_j f(\bm{w}_j)
  \]
  となる．よって$\Im f = \Span\bigl(f(\bm{w}_1), \dots, f(\bm{w}_l)\bigr)$である．

  線型独立性について，
  \[
    \sum_{j=1}^{l} c_j f(\bm{w}_j) = \bm{0}
  \]
  とする．\prref{線型写像と線型結合}より$f\left(\sum_{j=1}^{l} c_j \bm{w}_j\right) = \bm{0}$，すなわち
  \[
    \sum_{j=1}^{l} c_j \bm{w}_j \in \Ker f
  \]
  である．$(\bm{u}_1, \dots, \bm{u}_k)$は$\Ker f$の基底であるから，
  \[
    \sum_{j=1}^{l} c_j \bm{w}_j = \sum_{i=1}^{k} d_i \bm{u}_i
  \]
  と表せる．移項すると
  \[
    \sum_{i=1}^{k} d_i \bm{u}_i - \sum_{j=1}^{l} c_j \bm{w}_j = \bm{0}
  \]
  となり，$(\bm{u}_1, \dots, \bm{u}_k, \bm{w}_1, \dots, \bm{w}_l)$の線型独立性より，とくに$c_1 = \cdots = c_l = 0$である．

  以上より$(f(\bm{w}_1), \dots, f(\bm{w}_l))$は$\Im f$の基底であり，$\dim \Im f = l$である．
  したがって，
  \[
    \dim \Ker f + \dim \Im f = k + l = n
  \]
  が成り立つ．これが証明すべきことであった．
\end{proof}
\end{leftbar}

\partref{線型独立性と階数}で「線型独立な列の最大個数」として導入した階数（\deref{階数}）は，この公式を通じて「像の次元」として捉え直される．

\begin{cor}{階数と像・核の次元}{階数と像の次元}
  $A$を$m \times n$行列とし，$T_A \colon K^n \to K^m$を$A$の定める線型写像とする．このとき，
  \[
    \dim \Im T_A = \rank A,
    \qquad
    \dim \Ker T_A = n - \rank A
  \]
  が成り立つ．
\end{cor}

\begin{leftbar}
\begin{proof}
  $A$の列ベクトルを$\bm{a}_1, \dots, \bm{a}_n \in K^m$とする．
  任意の$\bm{x} \in K^n$に対して，$\bm{x} = \sum_{j=1}^{n} x_j \bm{e}_j$，$A\bm{e}_j = \bm{a}_j$と\prref{線型写像と線型結合}より
  \[
    A\bm{x} = \sum_{j=1}^{n} x_j \bm{a}_j
  \]
  であるから，$\Im T_A = \Span(\bm{a}_1, \dots, \bm{a}_n)$である．

  $r = \rank A$とおく．
  階数の定義（\deref{階数}）より，$\bm{a}_1, \dots, \bm{a}_n$のうちから線型独立な$r$本$\bm{a}_{j_1}, \dots, \bm{a}_{j_r}$が選べる．
  任意の列$\bm{a}_j$に対して，$r + 1$本の組$\bm{a}_{j_1}, \dots, \bm{a}_{j_r}, \bm{a}_j$は$r$の最大性より線型従属であるから，\lemref{一般の線型独立な組の拡大}の対偶より
  \[
    \bm{a}_j \in \Span(\bm{a}_{j_1}, \dots, \bm{a}_{j_r})
  \]
  である．よって\prref{線型包の最小性}より
  \[
    \Im T_A = \Span(\bm{a}_1, \dots, \bm{a}_n)
    \subset \Span(\bm{a}_{j_1}, \dots, \bm{a}_{j_r})
    \subset \Im T_A
  \]
  となり，$\Im T_A = \Span(\bm{a}_{j_1}, \dots, \bm{a}_{j_r})$である．
  $(\bm{a}_{j_1}, \dots, \bm{a}_{j_r})$は線型独立かつ$\Im T_A$を生成するから，$\Im T_A$の基底であり，$\dim \Im T_A = r = \rank A$である．

  後半の等式は，\thref{線型写像の次元公式}より
  \[
    \dim \Ker T_A = n - \dim \Im T_A = n - \rank A
  \]
  として従う．これが証明すべきことであった．
\end{proof}
\end{leftbar}

\begin{remark}[次元公式の解釈と解の自由度]
  \thref{線型写像の次元公式}は，定義域のもつ$n$次元分の情報が，「核としてつぶれる分」と「像として生き残る分」とに過不足なく分かれることを述べている．
  \exref{像と核の例}の正射影$f_3$でいえば，$\dim \Ker f_3 = 1$（$y$軸方向がつぶれる）と$\dim \Im f_3 = 1$（$x$軸が残る）の和が，定義域$\mathbb{R}^2$の次元$2$に一致している．
  文献によっては，\thref{次元公式}ではなくこちらを指して「次元公式」あるいは「次元定理」とよぶことも多い．

  また，\coref{階数と像の次元}の$\dim \Ker T_A = n - \rank A$は，連立一次方程式の解の様子を定量的に述べている．
  斉次連立一次方程式$A\bm{x} = \bm{0}$の解空間は$\Ker T_A$であり（\prref{連立一次方程式と像と核}），その次元$n - \rank A$は，解を表すのに必要なパラメータの個数（解の自由度）にあたる．
  とくに$\rank A = n$のときは$\Ker T_A = \{\bm{0}\}$となり，\prref{単射性と核}より$T_A$は単射，すなわち$A\bm{x} = \bm{0}$は自明な解$\bm{x} = \bm{0}$しかもたない．

  なお，のちに一般の線型空間のあいだの線型写像を導入すれば，有限次元の線型空間$V$から線型空間$W$への線型写像$f$に対しても，まったく同じ証明によって
  \[
    \dim \Ker f + \dim \Im f = \dim V
  \]
  が成り立つ．
\end{remark}

\phantomsection
\section{基底の取替え}

ベクトルそのものは，どの基底を選ぶかとは無関係に定まっている．
一方，その座標ベクトルは，基底の選び方に応じて変化する「表示」である（\prref{座標ベクトルの性質}のあとのRemarkを参照）．
本節では，この「対象」と「表示」の区別を踏まえ，基底を取り替えたときに座標ベクトルがどのような規則で変換されるかを調べる．

\phantomsection
\subsection{基底の行列表記}

基底を取り替える操作を見通しよく扱うために，記法を一つ用意する．

\begin{definition}{ベクトルの組と行列の積}{ベクトルの組と行列の積}
  $V$を体$K$上の線型空間とし，$\bm{v}_1, \bm{v}_2, \dots, \bm{v}_n \in V$とする．
  $n \times m$行列$P = (p_{ij})$に対して，
  \[
    (\bm{v}_1, \bm{v}_2, \dots, \bm{v}_n) P
    :=
    \left( \sum_{i=1}^{n} p_{i1} \bm{v}_i,~ \sum_{i=1}^{n} p_{i2} \bm{v}_i,~ \dots,~ \sum_{i=1}^{n} p_{im} \bm{v}_i \right)
  \]
  と定める．すなわち，右辺の第$j$成分は$P$の第$j$列を係数とする$\bm{v}_1, \dots, \bm{v}_n$の線型結合である．

  同様に，$\bm{x} = (x_i) \in K^n$に対して
  \[
    (\bm{v}_1, \bm{v}_2, \dots, \bm{v}_n) \bm{x}
    := \sum_{i=1}^{n} x_i \bm{v}_i \in V
  \]
  と定める．
\end{definition}

\begin{remark}[形式的な行列記法について]
  これは，ベクトルを成分とする「$1 \times n$行列」と通常の行列との積を，形式的に行列の積の定義（\deref{行列の積}）と同じ規則で定めたものである．
  $V$が一般の線型空間の場合には成分どうしの積は定義されないので，あくまで記法の約束として理解してほしい．

  この記法のもとで，基底$\mathcal{B} = (\bm{v}_1, \dots, \bm{v}_n)$と座標ベクトルの関係は
  \[
    \bm{v} = \mathcal{B} \, [\bm{v}]_{\mathcal{B}}
  \]
  と簡潔にかける．
\end{remark}

\begin{prop}{ベクトルの組と行列の積の結合律}{ベクトルの組と行列の積の結合律}
  $\bm{v}_1, \dots, \bm{v}_n \in V$とし，$P$を$n \times m$行列，$Q$を$m \times l$行列とする．
  $\mathcal{V} = (\bm{v}_1, \dots, \bm{v}_n)$とかくと，
  \[
    (\mathcal{V} P) Q = \mathcal{V} (PQ)
  \]
  が成り立つ．
\end{prop}

\begin{leftbar}
\begin{proof}
  $P = (p_{ij})$，$Q = (q_{jk})$とし，$\mathcal{V}P$の第$j$成分を$\bm{w}_j = \sum_{i=1}^{n} p_{ij} \bm{v}_i$とおく．
  このとき，$(\mathcal{V}P)Q$の第$k$成分は
  \begin{align*}
    \sum_{j=1}^{m} q_{jk} \bm{w}_j
    &= \sum_{j=1}^{m} q_{jk} \left( \sum_{i=1}^{n} p_{ij} \bm{v}_i \right) \\
    &= \sum_{i=1}^{n} \left( \sum_{j=1}^{m} p_{ij} q_{jk} \right) \bm{v}_i
  \end{align*}
  である．最後の括弧の中は$PQ$の$(i,k)$成分にほかならないから，これは$\mathcal{V}(PQ)$の第$k$成分に等しい．
  これが証明すべきことであった．
\end{proof}
\end{leftbar}

\phantomsection
\subsection{基底の取替行列}

\begin{definition}{基底の取替行列}{取替行列}
  $\mathcal{B} = (\bm{v}_1, \dots, \bm{v}_n)$，$\mathcal{B}' = (\bm{v}'_1, \dots, \bm{v}'_n)$をともに線型空間$V$の基底とする．
  各$\bm{v}'_j$を$\mathcal{B}$で表したときの係数を$p_{ij}$とする．すなわち，
  \[
    \bm{v}'_j = \sum_{i=1}^{n} p_{ij} \bm{v}_i \quad (1 \le j \le n)
  \]
  とし，$P = (p_{ij})$とおく．
  このとき，
  \[
    \mathcal{B}' = \mathcal{B} P
  \]
  が成り立つ．この$P$を，$\mathcal{B}$から$\mathcal{B}'$への\textbf{基底の取替行列}\index{きていのとりかええぎょうれつ@基底の取替行列}という．
\end{definition}

\begin{remark}[取替行列の列の意味]
  取替行列$P$の第$j$列は，新しい基底ベクトル$\bm{v}'_j$の，古い基底$\mathcal{B}$に関する座標ベクトルである：
  \[
    P = \bigl( [\bm{v}'_1]_{\mathcal{B}},~ [\bm{v}'_2]_{\mathcal{B}},~ \dots,~ [\bm{v}'_n]_{\mathcal{B}} \bigr)
  \]
  「新しい基底を，古い基底で表したものを列に並べる」と覚えるとよい．
  添字の付き方が$\bm{v}'_j = \sum_i p_{ij} \bm{v}_i$と，和をとるのが第$1$添字である点に注意してほしい．
  これは$\mathcal{B}' = \mathcal{B}P$という表記が成り立つように定めた結果である．
\end{remark}

\begin{prop}{取替行列の正則性}{取替行列の正則性}
  $\mathcal{B}$から$\mathcal{B}'$への取替行列を$P$とする．
  このとき$P$は正則であり，$\mathcal{B}'$から$\mathcal{B}$への取替行列は$P^{-1}$である．
\end{prop}

\begin{leftbar}
\begin{proof}
  $\mathcal{B}'$から$\mathcal{B}$への取替行列を$Q$とすると，
  \[
    \mathcal{B} = \mathcal{B}' Q = (\mathcal{B}P)Q = \mathcal{B}(PQ)
  \]
  である（最後の等号は\prref{ベクトルの組と行列の積の結合律}による）．
  一方，明らかに$\mathcal{B} = \mathcal{B}E_n$である．
  $\mathcal{B}$は基底であるから表示は一意であり（\thref{基底による表示の一意性}），両辺の各成分を比較して
  \[
    PQ = E_n
  \]
  を得る．$\mathcal{B}$と$\mathcal{B}'$の役割を入れ替えれば$QP = E_n$も同様に従う．
  よって$P$は正則で$Q = P^{-1}$である．
  これが証明すべきことであった．
\end{proof}
\end{leftbar}

\phantomsection
\subsection{座標の変換}

基底を取り替えると，同じベクトルの座標ベクトルはどう変わるだろうか．

\begin{theorem}{座標の変換公式}{座標の変換公式}
  $\mathcal{B}$，$\mathcal{B}'$を線型空間$V$の基底とし，$P$を$\mathcal{B}$から$\mathcal{B}'$への取替行列とする．
  このとき，任意の$\bm{v} \in V$に対して
  \[
    [\bm{v}]_{\mathcal{B}} = P \, [\bm{v}]_{\mathcal{B}'},
    \quad \text{すなわち} \quad
    [\bm{v}]_{\mathcal{B}'} = P^{-1} [\bm{v}]_{\mathcal{B}}
  \]
  が成り立つ．
\end{theorem}

\begin{leftbar}
\begin{proof}
  $\bm{x} = [\bm{v}]_{\mathcal{B}}$，$\bm{x}' = [\bm{v}]_{\mathcal{B}'}$とおく．
  座標ベクトルの定義より
  \[
    \bm{v} = \mathcal{B}\bm{x} = \mathcal{B}'\bm{x}'
  \]
  である．ここで$\mathcal{B}' = \mathcal{B}P$であるから，\prref{ベクトルの組と行列の積の結合律}より
  \[
    \mathcal{B}'\bm{x}' = (\mathcal{B}P)\bm{x}' = \mathcal{B}(P\bm{x}')
  \]
  となる．よって
  \[
    \mathcal{B}\bm{x} = \mathcal{B}(P\bm{x}')
  \]
  であり，$\mathcal{B}$による表示の一意性（\thref{基底による表示の一意性}）から$\bm{x} = P\bm{x}'$を得る．
  \prref{取替行列の正則性}より$P$は正則であるから，$\bm{x}' = P^{-1}\bm{x}$である．
  これが証明すべきことであった．
\end{proof}
\end{leftbar}

\begin{remark}[基底と座標の変換の向き]
  ここで，向きが「逆になる」ことに注意してほしい．
  \[
    \text{基底：} \quad \mathcal{B}' = \mathcal{B}P,
    \qquad
    \text{座標：} \quad [\bm{v}]_{\mathcal{B}'} = P^{-1} [\bm{v}]_{\mathcal{B}}
  \]
  基底は$P$によって新しいものに移るのに対し，座標は$P^{-1}$によって移るのである．
  この点は混乱しやすいので，具体例で確かめておこう．

  $\mathbb{R}^1$において$\bm{v}_1 = 1$を基底とし，$\bm{v}'_1 = 2$を新しい基底とする．
  このとき$\bm{v}'_1 = 2\bm{v}_1$であるから$P = (2)$である．
  ベクトル$\bm{v} = 6$の座標は，$\mathcal{B}$では$6$，$\mathcal{B}'$では$3$である．
  基底ベクトルを$2$倍に「伸ばす」と，同じベクトルを測る目盛りの数は$1/2$倍に「縮む」わけである．

  ものさしの目盛りを細かくすれば読み取る数値は大きくなり，粗くすれば小さくなる，という日常の感覚と同じである．
  このように，基底の変換と座標の変換とに互いに逆の行列（$P$と$P^{-1}$）が現れることを指して，座標は基底の変換に対して\textbf{反変的}\index{はんぺんてき@反変的}に変換する，と表現することがある．
\end{remark}

\begin{example}{座標の変換}{座標の変換}
  $V = \mathbb{R}^2$とし，$\mathcal{B} = (\bm{e}_1, \bm{e}_2)$を標準基底，
  \[
    \mathcal{B}' = (\bm{v}'_1, \bm{v}'_2), \quad
    \bm{v}'_1 = \begin{pmatrix} 1 \\ 1 \end{pmatrix}, \quad
    \bm{v}'_2 = \begin{pmatrix} 1 \\ -1 \end{pmatrix}
  \]
  とする．
  $\bm{v}'_1 = \bm{e}_1 + \bm{e}_2$，$\bm{v}'_2 = \bm{e}_1 - \bm{e}_2$であるから，$\mathcal{B}$から$\mathcal{B}'$への取替行列は
  \[
    P = \begin{pmatrix} 1 & 1 \\ 1 & -1 \end{pmatrix}
  \]
  である．$\det P = -2 \ne 0$より$P$は正則で，
  \[
    P^{-1} = \frac{1}{2}\begin{pmatrix} 1 & 1 \\ 1 & -1 \end{pmatrix}
  \]
  である．

  いま，$\bm{v} = \begin{pmatrix} 5 \\ 1 \end{pmatrix}$を考える．
  $\mathcal{B}$は標準基底であるから$[\bm{v}]_{\mathcal{B}} = \begin{pmatrix} 5 \\ 1 \end{pmatrix}$であり，
  \[
    [\bm{v}]_{\mathcal{B}'}
    = P^{-1}\begin{pmatrix} 5 \\ 1 \end{pmatrix}
    = \frac{1}{2}\begin{pmatrix} 6 \\ 4 \end{pmatrix}
    = \begin{pmatrix} 3 \\ 2 \end{pmatrix}
  \]
  となる．実際，
  \[
    3\begin{pmatrix} 1 \\ 1 \end{pmatrix} + 2\begin{pmatrix} 1 \\ -1 \end{pmatrix}
    = \begin{pmatrix} 5 \\ 1 \end{pmatrix}
  \]
  であり，確かに成り立っている．
\end{example}

\begin{remark}[標準基底の場合と行列式による判定]
  $V = K^n$で$\mathcal{B}$が標準基底の場合，取替行列$P$は「新しい基底ベクトルを列として並べた行列」そのものである．
  この場合，$\bm{v} = P[\bm{v}]_{\mathcal{B}'}$という関係は，\partref{線型独立性と階数}のRemarkで見た「$A\bm{x}$は$A$の列ベクトルの線型結合であり，その係数は$\bm{x}$の成分である」という事実にほかならない．

  また，$K = \mathbb{R}$または$\mathbb{C}$とする．\partref{行列式}では行列式に関する諸性質を複素数行列について証明したが，実数行列についても同じ定義式と証明がそのまま通用する（\partref{行列式}冒頭の注意を参照）．よって，\coref{次元による基底の判定}と\prref{行列式と階数}を合わせると，
  $K^n$の$n$個のベクトル$\bm{v}'_1, \dots, \bm{v}'_n$が基底をなすことは，それらを列に並べた行列$P$について$\det P \ne 0$が成り立つことと同値である．
  基底かどうかの判定が，行列式の計算に帰着されるのである．
\end{remark}

\newpage
\phantomsection

\part*{参考文献}

\noindent
本書は
\cite{saitomasahiko1966}，
\cite{arai}，
\cite{yukiegunron}，
\cite{fujioka2021}
を参考にした．

\medskip

\printbibliography[heading=none]
% 索引
% 参考文献の次のページ（偶数ページ）から索引を開始する．
% 通常の \part は右ページ開始だが，索引だけは右ページ送りを無効にする．
\clearpage
\thispagestyle{empty}
\setcounter{section}{6}
\renewcommand{\indexname}{索引}

\begingroup
  \makeatletter
  \def\part@newpage{\clearpage\thispagestyle{empty}}
  \makeatother
  \let\section\part
  \pagestyle{empty}
  \printindex
\endgroup

% 裏表紙の直前に白紙を1ページ挿入
\clearpage
\thispagestyle{empty}
\null
\clearpage

\MathNoteBackCover{Linear Algebra}{2026 Edition}

\end{document}